% Version: 2026.1.0.0
% ============================================================================
% Engineering Studio AI — Multi-Modal Multi-Agent Discipline Framework
% Formal Model and Mathematical Proofs
% Author: Hadrian Hu (hadrian.hu@gmail.com)
% Governed by: coding_stds/documentation/latex_style_guide.txt
%              coding_stds/mathematics/mathematical_proof_structures_standards.txt
% Date: \today
% ============================================================================
\documentclass[12pt]{article}

% ============================================================================
% PREAMBLE
% ============================================================================

% --- Font & Encoding ---------------------------------------------------
\usepackage[utf8]{inputenc}   % WHAT: UTF-8 source encoding (Std 14.1)
\usepackage[T1]{fontenc}      % WHAT: T1 font encoding (Std 14.2)
\usepackage{times}            % WHAT: Times New Roman body font (Std 2.1)

% --- Page Layout ---------------------------------------------------------
\usepackage[margin=1in]{geometry}  % WHAT: 1in margins (Std 2.3)
\usepackage{setspace}
\singlespacing                     % WHY: single spacing throughout (Std 2.4)

% --- Headers & Footers ----------------------------------------------------
\usepackage{fancyhdr}
\usepackage{lastpage}
\pagestyle{fancy}
\fancyhf{}
\setlength{\headheight}{15pt} % WHY: fancyhdr's own minimum for this font size; avoids "headheight too small" warning on every page
\rhead{\thepage\ of \pageref{LastPage}}
\lhead{Engineering Studio AI — Formal Model \& Proofs}
\renewcommand{\headrulewidth}{0.4pt}

% --- Math ------------------------------------------------------------------
\usepackage{amsmath}    % WHAT: aligned/multline/equation environments
\usepackage{amssymb}    % WHAT: blacksquare, mathbb, mathcal
\usepackage{amsthm}     % WHAT: theorem/lemma/proof environments (Std proof §4.1.2)
\usepackage{mathtools}  % WHAT: enhanced math tooling, coloneqq

% --- Graphics & Tables -------------------------------------------------------
\usepackage{graphicx}
\usepackage{booktabs}
\usepackage{array}
\usepackage{multirow}
\usepackage{makecell}
\usepackage{float}
\usepackage{tabularx}
\usepackage{longtable}
\usepackage{ragged2e}

% --- Code & Pseudocode -------------------------------------------------------
\usepackage{listings}
\usepackage{algorithm}
\usepackage{algpseudocode}

% --- Lists -------------------------------------------------------------------
\usepackage{enumitem}
\setlist[enumerate]{itemsep=0pt, parsep=0pt}

% --- Typography safety -------------------------------------------------------
\usepackage[final]{microtype}
\setlength{\emergencystretch}{15pt}
\widowpenalty=10000
\clubpenalty=10000

% --- TOC ------------------------------------------------------------------
\usepackage{tocloft}
\setcounter{tocdepth}{2}

% --- Hyperlinks (loaded second-to-last, per Std 27.9) -----------------------
\usepackage[hidelinks]{hyperref}
\hypersetup{
    colorlinks=false,
    pdftitle={Engineering Studio AI: Formal Model and Mathematical Proofs},
    pdfauthor={Hadrian Hu},
    pdfsubject={Multi-Agent Engineering Workflow Formalization}
}

% --- Theorem environments (Std proof §4.2.1) --------------------------------
\theoremstyle{definition}
\newtheorem{theorem}{Theorem}
\newtheorem{lemma}[theorem]{Lemma}
\newtheorem{corollary}[theorem]{Corollary}
\newtheorem{definition}[theorem]{Definition}
\newtheorem{proposition}[theorem]{Proposition}
\newtheorem{remark}[theorem]{Remark}
\newtheorem{hypothesis}{Hypothesis}

% --- Custom commands ---------------------------------------------------------
\newcommand{\code}[1]{\texttt{#1}}
\newcommand{\FN}[1]{\textsc{fn-#1}}
\newcommand{\STD}[1]{\textsc{std-#1}}
\DeclareMathOperator{\Approved}{Approved}
\DeclareMathOperator{\Rejected}{Rejected}

% --- Long-filename citation macros (Std 23.1 overfull-hbox prevention) -----
% WHY: \texttt{} disables hyphenation, and long VISION_*.md identifiers with
% no interword glue cannot be line-broken by TeX, causing large overfull
% \hbox warnings wherever they are cited. \allowbreak inserted after each
% underscore gives the line breaker a legal (invisible) break point without
% altering the rendered filename. Define ONE macro per governing document
% and cite the macro everywhere instead of retyping/re-splitting the raw
% \code{...} string — prevents regression and keeps citations consistent.
\newcommand{\VisionClean}{\code{VISION\_MULTIMODAL\_\allowbreak DEPLOYABLE\_\allowbreak AGENTS\_CLEAN.md}}
\newcommand{\VisionGov}{\code{VISION\_AGENTS\_\allowbreak COORDINATION\_\allowbreak ORCHESTRATION\_\allowbreak GOVERNANCE.md}}
\newcommand{\VisionFunc}{\code{VISION\_AGENTS\_\allowbreak BY\_FUNCTION.md}}
\newcommand{\VisionFlow}{\code{VISION\_AGENT\_\allowbreak WORKFLOW\_\allowbreak INTERAGENT\_\allowbreak COORDINATION.md}}
\newcommand{\VisionTheme}{\code{VISION\_AGENTS\_\allowbreak BY\_THEME.md}}
\newcommand{\VisionHackathon}{\code{VISION\_AMD\_\allowbreak LABLAB\_\allowbreak HACKATHON\_\allowbreak ENGINEERING\_STUDIO.md}}
\newcommand{\ProofStd}{\code{mathematical\_proof\_\allowbreak structures\_standards.txt}}
% sloppypar



% ============================================================================
\begin{document}

% ============================================================================
% TITLE PAGE
% ============================================================================
\begin{titlepage}
    \centering
    \vspace*{1cm}
    {\LARGE \textbf{Engineering Studio AI}}\\[0.3cm]
    {\large A Multi-Modal, Multi-Agent Discipline Framework for the}\\
    {\large Full Software--Hardware Engineering Lifecycle:}\\
    {\large Formal Model and Mathematical Proofs}

    \vspace{2cm}

    {\large \textbf{Author}}\\
    Hadrian Hu\\
    \texttt{hadrian.hu@gmail.com}

    \vspace{1cm}

    {\large CodingStandardsRef Research Notes}\\
    Multimodal Deployable Agents Platform (MDAP)

    \vspace{2cm}

    {\large \today}

    \vfill
\end{titlepage}

% ============================================================================
% ABSTRACT
% ============================================================================
\newpage
\section*{Abstract}

\noindent
\textbf{Context.}

Engineering Studio AI is a proposed multi-agent, multi-modal system
in which a single natural-language product brief is decomposed by a
root Orchestrator (\FN{01}) into concurrent and sequential work
performed by narrowly-scoped specialist, reviewer, validator,
adversarial-challenge, testing, documentation, and gate agents,
spanning the full software and hardware (emulated) engineering
lifecycle.
The architecture is operationally specified across five internal
vision documents and is instantiated as the $\sim$250-role
Multimodal Deployable Agents Platform (MDAP) agent catalog.

\medskip
\noindent
\textbf{Problem.}
Operational specifications and design rationale alone are
insufficient to establish that a multi-agent system is trustworthy at
production scale.
Four structural properties are required:
\begin{enumerate}
    \item[(P1)] The nominal workflow graph must be acyclic so that a
          deterministic execution order exists and no phase can
          re-enter itself without an explicit rework mechanism.
    \item[(P2)] Every reject/rework loop must terminate within a
          finite, statically-computable bound so that no run can
          diverge.
    \item[(P3)] The Non-Overlap (single-responsibility) rule must
          imply a concrete lower bound on distinct agent instances,
          preventing silent responsibility collapse.
    \item[(P4)] The trust-tier safety invariant must be preserved
          across an arbitrarily long sequence of subagent invocations
          so that untrusted output can never be mistaken for a
          certified instruction.
\end{enumerate}
None of P1--P4 is established by the governing vision documents,
which state these properties as design intent rather than as proved
theorems.

\medskip
\noindent
\textbf{Approach.}
This paper formalizes the MDAP architecture as a seven-tuple
$\mathcal{R} = (G, \mathcal{M}, \tau, C, R, F_{\max}, D_{\max})$,
where:
\begin{enumerate}
    \item[(A1)] $G = (V, E)$ is a directed graph of $|V|=11$ typed
          artifact-handoff phases, with edges encoding the
          ``must-complete-before'' relation from the governing Phase
          Ownership Table.
    \item[(A2)] $\mathcal{M} = (Q, \Sigma, \delta, q_0, F)$ is a
          finite deterministic automaton over $|Q|=8$ artifact
          lifecycle states, with initial state
          $q_0 = \mathrm{Drafted}$ and accepting set
          $F = \{\mathrm{Approved}\}$.
    \item[(A3)] $\tau : \mathcal{C} \to
          \{\mathrm{T\text{-}0}, \mathrm{T\text{-}1},
          \mathrm{T\text{-}2}, \mathrm{T\text{-}3}\}$ assigns a
          trust tier to every content item in a session, with T-0
          being the highest authority (model constraints) and T-3
          the lowest (every subagent return message).
    \item[(A4)] $C = (V_C, E_C)$ is an undirected conflict graph over
          the 12 function classes \FN{00}--\FN{11}, where an edge
          $\{u,v\} \in E_C$ encodes the Non-Overlap Rule's prohibition
          on assigning both $u$ and $v$ to one agent instance.
    \item[(A5)] $R \in \mathbb{N}_0$ is the per-phase-gate retry
          limit, giving a global rework budget
          $B \coloneqq R \cdot |V| = 11R$.
    \item[(A6)] $F_{\max} \in \mathbb{N}$ is the per-turn fan-out
          ceiling; $D_{\max} = 2$ is the nesting-depth limit.
\end{enumerate}
All definitions are grounded, without invented detail, in the five
governing vision documents cited in Section~\ref{sec:background}.

\medskip
\noindent
\textbf{Results.}
Four theorems are proved directly from axioms already adopted as
engineering policy:
\begin{enumerate}
    \item[(R1)] \textbf{Acyclicity} (Theorem~\ref{thm:acyclic}).
          The nominal workflow graph $G$ is a directed acyclic graph
          (DAG).
          \emph{Method:} construction of a strict topological-index
          function $\phi : V \to \{1,\ldots,11\}$ such that
          $(u,v)\in E \Rightarrow \phi(u) < \phi(v)$; a directed
          cycle would require $\phi(u) < \phi(u)$, a contradiction.
          \emph{Boundary cases addressed:} single-vertex graph,
          empty edge set, maximum-density graph ($|E| = 121$).
    \item[(R2)] \textbf{Bounded Termination} (Theorem~\ref{thm:rework}).
          Every run terminates within at most $P + B = 11 + 11R$
          phase-gate transitions, where $R \in \mathbb{N}_0$ is the
          per-phase retry limit.
          \emph{Method:} monovariant argument --- a rejection counter
          $b_i$ is strictly non-decreasing and bounded above by $B$;
          the circuit-breaker fires at $b_i = B$, guaranteeing
          termination even in the worst case.
          \emph{Boundary cases addressed:} $R=0$ (immediate
          escalation on first rejection), $B=0$, and every phase
          saturating its retry allowance simultaneously.
    \item[(R3)] \textbf{Minimum Agent Multiplicity}
          (Theorem~\ref{thm:multiplicity}).
          Any deployment compliant with the Non-Overlap Rule requires
          at least $\chi(C[K]) = 4$ distinct agent instances, where
          $K = \{\FN{02}, \FN{04}, \FN{05}, \FN{09}\}$ is the
          mandatory-phase clique and $\chi$ denotes chromatic number.
          \emph{Method:} three-stage proof ---
          (I)~exhibit the 4-clique $K \subseteq V_C$;
          (II)~prove $\chi(C[K]) \geq 4$ by clique lower bound;
          (III)~exhibit a 4-coloring achieving $\chi(C[K]) = 4$.
          \emph{Boundary cases addressed:} single artifact, empty
          conflict graph (would give a lower bound of 1, vacuously
          consistent), and maximum-clique configurations.
    \item[(R4)] \textbf{Trust-Tier Invariant}
          (Theorem~\ref{thm:trust}).
          For every subagent-originated content item $c$ and every
          invocation sequence of arbitrary finite length, the
          assignment $\tau(c) = \mathrm{T\text{-}3}$ is preserved
          unless $c$ has been the subject of an explicit
          \FN{08}-issued $\Approved$ certificate event.
          \emph{Method:} structural induction over the length $n$ of
          the invocation sequence; the inductive step examines all
          four cases of event type and confirms that none promotes
          $\tau(c)$ without an \FN{08} $\Approved$ event.
          \emph{Boundary cases addressed:} empty sequence ($n=0$),
          sequence of length 1, sequences containing exclusively
          $\Rejected$ events, and the boundary case
          $\tau(c) = \mathrm{T\text{-}0}$ (model constraint,
          provably unreachable from T-3 without a constitutional
          amendment).
\end{enumerate}
A fifth property --- a token/cost invocation-count convergence
bound --- is provided as an explicitly labelled proof sketch
(Section~\ref{sec:sketch}).
The sketch identifies the missing precondition (ratified per-agent
token-budget constants) and the minimum additional work required to
elevate it to a full proof.

\medskip
\noindent
\textbf{Contribution.}
A reusable formal vocabulary --- the seven-tuple $\mathcal{R}$,
workflow graph $G$, artifact state machine $\mathcal{M}$, conflict
graph $C$, and trust-tier function $\tau$ --- capable of serving as
the stable foundation for auditing any future extension of the MDAP
catalog without re-deriving P1--P4 from scratch.

\medskip
\noindent
\textbf{Auditability guarantee.}
Every hypothesis used in every proof is either a standard
mathematical fact (graph theory, finite automata, combinatorics) or
is directly cited to a specific section of a governing document.
No new axioms are introduced.
No proof step is dismissed as ``obvious'' without explicit
justification.
All boundary conditions, degenerate cases, and edge cases are
addressed in dedicated sub-proofs.

\bigskip
\noindent
\textbf{Keywords:}
multi-agent systems,
workflow formalization,
directed acyclic graphs,
bounded rework termination,
trust-tier invariant,
non-overlap rule,
chromatic number,
engineering lifecycle automation,
MDAP,
formal verification

\newpage

% ============================================================================
% EXECUTIVE SUMMARY
% ============================================================================
\section*{Executive Summary}

\begin{enumerate}

    % -----------------------------------------------------------------------
    \item \textbf{Architectural context and motivation.}
    \begin{enumerate}
        \item[(1.1)] Engineering Studio AI is a proposed multi-agent,
              multi-modal system in which a single natural-language
              product brief is progressively decomposed by a root
              Orchestrator (\FN{01}) into concurrent and sequential
              tasks executed by narrowly-scoped agents spanning the
              full software and hardware (emulated) engineering
              lifecycle.
        \item[(1.2)] The architecture is operationally governed by five
              internal vision documents: \VisionClean{} (Layer~0--8
              hierarchy, Trust Tier table), \VisionGov{} (statelessness,
              loop-safety, task-specification schema), \VisionFunc{}
              (12 function classes \FN{00}--\FN{11}, Non-Overlap Rule),
              \VisionFlow{} (Phase Ownership Table, artifact state
              machine), and \VisionHackathon{} (hardware-emulation
              boundary).
        \item[(1.3)] These documents constitute authoritative operational
              specifications; they do not, however, provide
              machine-checkable proofs of the structural properties on
              which the architecture's trustworthiness depends.
              This paper closes that gap.
    \end{enumerate}

    % -----------------------------------------------------------------------
    \item \textbf{Problem statement (four required properties).}
    \begin{enumerate}
        \item[(2.1)] \emph{P1 --- Acyclicity.}
              The nominal workflow graph $G$ must be a DAG so that a
              valid topological execution order exists.
              Without acyclicity, a phase could be re-entered
              indefinitely without invoking the explicit rework
              mechanism, making termination unprovable.
        \item[(2.2)] \emph{P2 --- Bounded rework termination.}
              The reject/rework loop must terminate within a finite,
              statically-computable bound $P + B = 11 + 11R$.
              Without this bound, a sequence of rejections could cause
              the run to diverge; no SLA or token-budget ceiling could
              be guaranteed.
        \item[(2.3)] \emph{P3 --- Minimum agent multiplicity.}
              The Non-Overlap Rule must imply a concrete lower bound
              (at least 4) on distinct agent instances per artifact.
              Without this lower bound, a compliant deployment could
              silently collapse responsibilities into a single agent,
              re-introducing the monolithic failure modes detailed in
              Section~\ref{sec:intro:problem}.
        \item[(2.4)] \emph{P4 --- Trust-tier invariant.}
              The trust-tier assignment $\tau(c) = \mathrm{T\text{-}3}$
              for every subagent-originated content item $c$ must be
              preserved across an arbitrarily long invocation sequence.
              Without this invariant, a sequence of individually-valid
              tier promotions could allow an untrusted subagent return
              to be treated as a certified instruction, a direct safety
              violation.
        \item[(2.5)] \emph{Sufficiency of P1--P4.}
              Each of P1--P4 addresses an independently sufficient
              failure mode.
              Proving all four simultaneously establishes that the
              architecture is structurally sound for
              multi-disciplinary production use, conditional on the
              governing documents remaining in force.
    \end{enumerate}

    % -----------------------------------------------------------------------
    \item \textbf{Formal model.}
    \begin{enumerate}
        \item[(3.1)] The MDAP architecture is modelled as the
              seven-tuple
              \begin{align}
                \mathcal{R}
                \;=\;
                \bigl(G,\;\mathcal{M},\;\tau,\;C,\;R,\;
                       F_{\max},\;D_{\max}\bigr).
              \end{align}
        \item[(3.2)] Every component of $\mathcal{R}$ is defined
              by direct citation to a governing document; no component
              is invented.
              Full definitions with validation of well-formedness
              conditions, boundary cases, and cross-definition
              consistency are given in
              Sections~\ref{sec:background}--\ref{sec:model}.
        \item[(3.3)] The model is minimal: no structure is included
              beyond what the four proofs require, per Objective~O1b
              (Section~\ref{sec:intro:objectives}).
        \item[(3.4)] The model is reusable: the formal vocabulary is
              stable under addition of new agent roles to the MDAP
              catalog, provided the new roles are assigned function
              classes \FN{00}--\FN{11} and the conflict graph $C$ is
              updated accordingly.
    \end{enumerate}

    % -----------------------------------------------------------------------
    \item \textbf{Key results.}
    \begin{enumerate}
        \item[(4.1)] \emph{Theorem~\ref{thm:acyclic} (Acyclicity).}
              $G$ is a DAG.
              The topological index $\phi(v) \in \{1,\ldots,11\}$
              constructed in the proof witnesses a strict linear
              ordering compatible with $E$; its existence is both
              necessary and sufficient for acyclicity.
        \item[(4.2)] \emph{Theorem~\ref{thm:rework} (Bounded
              Termination).}
              Every run terminates within at most
              $11 + 11R$ phase-gate transitions, where $R \in
              \mathbb{N}_0$ is fixed before session start.
              The monovariant $b_i$ (rejection counter) witnesses
              termination: it is non-decreasing, bounded above by
              $B = 11R$, and strictly increases on every rejection
              event, so the sequence of rejection events is finite.
        \item[(4.3)] \emph{Theorem~\ref{thm:multiplicity} (Minimum
              Multiplicity).}
              Any Non-Overlap-compliant deployment requires
              $\geq 4$ distinct agent instances per artifact.
              The lower bound is tight: a 4-coloring of the induced
              subgraph $C[K]$ is exhibited, confirming
              $\chi(C[K]) = 4$ exactly.
        \item[(4.4)] \emph{Theorem~\ref{thm:trust} (Trust-Tier
              Invariant).}
              For every content item $c$ of subagent origin and every
              finite invocation sequence, $\tau(c) = \mathrm{T\text{-}3}$
              unless an explicit \FN{08} $\Approved$ event has
              occurred.
              The inductive proof covers all four event types and all
              boundary cases; no silent promotion path exists.
        \item[(4.5)] \emph{Section~\ref{sec:sketch} (Proof Sketch ---
              Token/Cost Convergence).}
              A geometric-series upper bound on total invocation count
              is derived under the assumption that per-agent token
              budgets are ratified.
              This section is explicitly marked as a sketch; the
              missing precondition and the upgrade path to a full
              proof are both stated.
    \end{enumerate}

    % -----------------------------------------------------------------------
    \item \textbf{Contribution and reuse.}
    \begin{enumerate}
        \item[(5.1)] The formal vocabulary of $\mathcal{R}$ is the
              primary contribution.
              Future extensions of the MDAP catalog (new agent roles,
              new phases, new trust tiers) can be audited against
              P1--P4 by:
              \begin{enumerate}[label=(\alph*)]
                \item updating $V$, $E$, $V_C$, $E_C$ to reflect the
                      extension;
                \item re-checking the hypothesis lists of
                      Theorems~\ref{thm:acyclic}--\ref{thm:trust}
                      against the updated governing documents;
                \item confirming that no hypothesis is violated by
                      the extension.
              \end{enumerate}
              If no hypothesis is violated, all four theorems hold for
              the extended system without re-proof.
        \item[(5.2)] The conflict graph $C$ is monotone-safe: adding
              edges to $E_C$ can only strengthen the Non-Overlap
              guarantees; it cannot invalidate
              Theorem~\ref{thm:multiplicity}.
        \item[(5.3)] The trust-tier function $\tau$ is extensible:
              adding new tiers between T-3 and T-0 preserves the
              inductive structure of Theorem~\ref{thm:trust}, provided
              the new tiers are ordered and no new promotion path
              bypasses \FN{08}.
    \end{enumerate}

    % -----------------------------------------------------------------------
    \item \textbf{Implications and recommended follow-on work.}
    \begin{enumerate}
        \item[(6.1)] Formal verification effort should concentrate on
              the token/cost convergence bound
              (Section~\ref{sec:sketch}).
              The specific gap --- ratification of per-agent
              token-budget constants --- is within the authority of
              the MDAP governance board and requires no new
              mathematical machinery.
        \item[(6.2)] Empirical validation against one real vertical
              slice (one complete run of the Impl--Review--Validate--
              Test--Gate pipeline on a non-trivial deliverable) would
              complement the formal results with observational
              evidence, per the roadmap recorded in \VisionHackathon{}.
        \item[(6.3)] Extension of Theorem~\ref{thm:multiplicity} to
              cover the full conflict graph $C$ (not just the induced
              subgraph $C[K]$) would sharpen the lower bound and may
              reveal tighter resource-planning constraints as the MDAP
              catalog grows.
        \item[(6.4)] No additional axiomatic assumptions are needed to
              pursue any of (6.1)--(6.3); all required mathematical
              tools (graph coloring, monovariant arguments, structural
              induction) are already deployed in this paper.
    \end{enumerate}

\end{enumerate}

% ============================================================================
\newpage
\tableofcontents
\newpage

% ============================================================================
\section{Introduction}
\label{sec:intro}

% ============================================================================
\subsection{Problem Statement}
\label{sec:intro:problem}

\begin{enumerate}

    \item \textbf{Architectural context.}
          Engineering Studio AI is a proposed system in which a single
          natural-language product brief is progressively decomposed by a
          root Orchestrator (\FN{01}) into concurrent and sequential tasks
          executed by narrowly-scoped specialist, reviewer, validator,
          adversarial-challenge, testing, documentation, and gate agents
          spanning the full software and hardware (emulated) engineering
          lifecycle.
          This decomposition is governed by the Layer~0--8 hierarchy and
          Challenge Division defined in \VisionClean{}, by the
          coordination and orchestration rules of \VisionGov{}, and by the
          function-class catalog of \VisionFunc{}.

    \item \textbf{The monolithic-agent failure mode.}
          A naive alternative --- a single large-context agent that
          simultaneously designs, reviews, validates, and certifies its own
          engineering output --- fails on three independently sufficient
          grounds:
          \begin{enumerate}
              \item[F1.] \emph{Responsibility collapse.}
                         When one agent performs both the Specialize
                         (\FN{02}/\FN{03}) and the Review (\FN{04}) function
                         on the same artifact, there exists no independent
                         party to detect self-inconsistency; the
                         Non-Overlap Rule (\S5.1, \VisionFunc{}) is
                         violated by construction.

              \item[F2.] \emph{Context-window bloat.}
                         A monolithic context must simultaneously hold the
                         full state of all engineering disciplines, all
                         in-flight artifacts, all review history, and all
                         governance constraints.
                         As the artifact set grows linearly, the required
                         context size grows at least linearly, and the
                         signal-to-noise ratio of each individual
                         decision decreases monotonically.

              \item[F3.] \emph{Unverifiable trust.}
                         A self-certifying agent's output cannot be
                         assigned a trust tier lower than T-2 (the user's
                         own instructions) without violating the tier
                         semantics of Table~8.1, \VisionClean{}.
                         Yet treating the agent's own output as T-2 would
                         allow it to silently override the Layer~0
                         constitution (T-1), a direct safety violation.
          \end{enumerate}
          Any one of F1--F3 is sufficient to disqualify the monolithic
          architecture for multi-disciplinary production use.

    \item \textbf{The adopted alternative.}
          The repository under study adopts the complementary architecture:
          many small, stateless, single-responsibility agents coordinated
          by one Orchestrator per invocation depth.
          The architectural claims made for this design --- that it
          terminates, that responsibilities never silently overlap, and that
          untrusted output is never mistaken for a certified instruction ---
          are operationally stated in the five governing vision documents
          cited in Section~\ref{sec:background}.
          However, those documents do not provide \emph{machine-checkable
          proofs} of these claims; they provide engineering policy and
          design rationale.

    \item \textbf{The gap this paper addresses.}
          The gap between an operationally stated architectural claim and a
          formally proved theorem is precisely the gap this paper closes.
          Specifically, four properties are stated as theorems and proved
          from first principles:
          \begin{enumerate}
              \item[G1.] The nominal workflow graph is acyclic
                         (Theorem~\ref{thm:acyclic}).
              \item[G2.] The reject/rework loop terminates within a
                         finite, computable bound
                         (Theorem~\ref{thm:rework}).
              \item[G3.] The Non-Overlap Rule implies a combinatorial lower
                         bound of four distinct agent instances per artifact
                         (Theorem~\ref{thm:multiplicity}).
              \item[G4.] The trust-tier safety invariant is preserved
                         across an arbitrarily long sequence of subagent
                         invocations (Theorem~\ref{thm:trust}).
          \end{enumerate}
          A fifth property, a token/cost invocation-count bound
          (Section~\ref{sec:sketch}), is provided as a proof sketch rather
          than a full proof, for reasons disclosed in that section.

    \item \textbf{What is not claimed.}
          This paper does \emph{not} claim to prove the semantic
          correctness of any individual specialist agent's engineering
          output (e.g., the electrical validity of a generated PCB layout
          or the functional correctness of emulated firmware).
          Such claims belong to domain-specific Reviewer and Validator
          agents at runtime and are outside the scope of this formal model.
          No new axioms are introduced; every hypothesis used in a proof
          is either a standard mathematical fact or a directly cited rule
          from a governing document.

\end{enumerate}

% ============================================================================
\subsection{Objectives}
\label{sec:intro:objectives}

\begin{enumerate}

    \item \textbf{Formal model construction (Objective~O1).}
          Construct a formal, graph-theoretic model of the Engineering
          Studio AI / MDAP workflow that is:
          \begin{enumerate}
              \item[O1a.] \emph{Grounded.}  Every primitive (vertex,
                          edge, state, tier, function class) is defined
                          by direct reference to an existing governing
                          document; no detail is invented.
              \item[O1b.] \emph{Sufficient.}  The model must be rich
                          enough to state and prove the four theorems
                          listed in G1--G4 above; no unnecessary
                          abstraction is added beyond what the proofs
                          require.
              \item[O1c.] \emph{Reusable.}  The formal vocabulary
                          (workflow graph, artifact state machine,
                          conflict graph, trust-tier function) must be
                          stable enough to serve as the foundation for
                          auditing future extensions of the MDAP catalog
                          without re-deriving these theorems from scratch.
          \end{enumerate}

    \item \textbf{Proof of four safety/liveness properties (Objective~O2).}
          For each of the four properties G1--G4, provide a proof that:
          \begin{enumerate}
              \item[O2a.] States its proof type (direct, inductive,
                          contradiction, or combinatorial) before the proof
                          body, per \S4.1.1 of \ProofStd{}.
              \item[O2b.] Lists all hypotheses explicitly and validates
                          each hypothesis against a cited governing
                          document or standard mathematical fact before the
                          proof body begins, per \S4.1.2 of the same
                          standard.
              \item[O2c.] Justifies every non-trivial inference step;
                          no step is implicitly "obvious."
              \item[O2d.] Addresses all boundary conditions, degenerate
                          cases, and edge cases explicitly; no case is
                          dismissed as trivial without demonstration.
              \item[O2e.] Closes with the \LaTeX{} \verb|\qedhere| marker
                          inside the proof environment, confirming
                          termination of the proof.
          \end{enumerate}

    \item \textbf{Disclosure of proof sketches (Objective~O3).}
          Identify \emph{exactly} which claims in this paper remain proof
          sketches rather than full proofs, and for each such claim:
          \begin{enumerate}
              \item[O3a.] Provide the explicit SKETCH label in the
                          section heading, per \S11.4.1 of \ProofStd{}.
              \item[O3b.] State the precise reason the proof is
                          incomplete (missing axiom, unspecified
                          distribution, or unvalidated empirical
                          assumption).
              \item[O3c.] Identify the minimum additional work required
                          to elevate the sketch to a full proof.
          \end{enumerate}
          Per O3, only one claim in this paper --- the token/cost
          convergence bound of Section~\ref{sec:sketch} --- is a sketch;
          the four core theorems (G1--G4) are full proofs.

    \item \textbf{Auditability (Objective~O4).}
          Structure every proof so that an independent third-party auditor
          with knowledge of standard graph theory, finite automata theory,
          combinatorics, and the cited governing documents can:
          \begin{enumerate}
              \item[O4a.] Verify each step without consulting the
                          authors.
              \item[O4b.] Identify the precise governing-document section
                          that justifies each non-mathematical premise.
              \item[O4c.] Confirm that no hidden assumption exists by
                          inspection of the hypothesis lists alone.
          \end{enumerate}

\end{enumerate}

% ============================================================================
\subsection{Scope and Exclusions}
\label{sec:intro:scope}

\begin{enumerate}

    \item \textbf{In scope: the workflow and governance layer.}
          This paper formalizes the following components, exactly as
          specified in the five vision documents cited in
          Section~\ref{sec:background}:
          \begin{enumerate}
              \item[S1.] The Layer~0--8 organizational hierarchy and
                         the Challenge Division (\VisionClean{}).
              \item[S2.] The statelessness axiom, task-specification
                         schema, loop-safety rules, and trust-tier
                         hierarchy (\VisionGov{}).
              \item[S3.] The twelve function classes \FN{00}--\FN{11}
                         and the Non-Overlap Rule (\VisionFunc{}).
              \item[S4.] The Phase Ownership Table and artifact state
                         machine (\VisionFlow{}).
              \item[S5.] The baseline coding standards (SOLID, JPL Power
                         of Ten, ACID-for-workflows, and the mathematical
                         proof structure standard) as codified in
                         \code{coding\_stds/}.
          \end{enumerate}

    \item \textbf{Out of scope: semantic correctness of specialist
          output.}
          This paper does \emph{not} formalize whether any individual
          specialist agent's engineering output is semantically correct
          (e.g., whether a generated hardware model is electrically valid,
          whether a firmware image is functionally correct, or whether a
          cost estimate is numerically accurate).
          \begin{enumerate}
              \item[E1.] Such claims require domain-specific formal
                         semantics (e.g., circuit-theory axioms,
                         programming-language operational semantics)
                         that are not within the scope of this
                         governance-layer formalization.
              \item[E2.] Responsibility for semantic correctness is
                         assigned at runtime to the domain-specific
                         Reviewer (\FN{04}) and Validator (\FN{05})
                         agents, as mandated by Sections~4.2 and~5.1
                         of \VisionFunc{}.
              \item[E3.] Excluding semantic correctness does not weaken
                         the four proved theorems; those theorems make
                         claims only about the \emph{structure} of the
                         workflow (acyclicity, termination, separation of
                         responsibilities, trust-tier preservation), not
                         about the \emph{content} of any artifact.
          \end{enumerate}

    \item \textbf{Out of scope: probabilistic or statistical claims.}
          No probability distribution over agent outputs, token costs,
          or rework frequencies is assumed anywhere in the four main
          theorems.
          The token/cost bound of Section~\ref{sec:sketch} bounds the
          \emph{count} of invocations deterministically; the cost per
          invocation is not modeled, which is the precise reason that
          section is a sketch.

    \item \textbf{Out of scope: physical hardware.}
          Per the disclosed assumption of \VisionHackathon{},
          ``hardware'' in this system means simulation or emulation
          (e.g., Gazebo, SPICE, BOM generation), not physical
          fabrication.
          No claim in this paper requires or assumes access to
          physical hardware; the formal model is hardware-agnostic.

    \item \textbf{Boundary condition on the governing documents.}
          The five governing documents cited above are treated as fixed,
          authoritative, and internally consistent.
          This paper does not prove their internal consistency; it
          derives theorems \emph{conditional on} their rules.
          Should any governing document be amended, the affected theorem
          hypotheses must be re-validated against the updated source
          before the theorems are considered to remain in force.

\end{enumerate}

% ============================================================================
\section{Background and Formal Definitions}
\label{sec:background}

\subsection{Governing Sources}
\begin{sloppypar}
This paper's model is grounded, without invented detail, in:
\VisionClean{} (Layer~0--8 and
Challenge Division definitions), \VisionGov{} (statelessness, task-specification
schema, loop-safety, trust tiers), \VisionFunc{}
(the twelve function classes \FN{00}--\FN{11} and the Non-Overlap Rule),
\VisionFlow{} (the Phase
Ownership Table and artifact state machine), and the baseline coding
standards in \code{coding\_stds/} (SOLID, JPL Power of Ten, ACID, and the
mathematical proof structure standard itself).
\end{sloppypar}

\subsection{Core Definitions}

\begin{definition}[Agent Instance]
\label{def:agent}
An \emph{agent instance} $a$ is a pair $a = (\mathrm{id}, S(a))$ where
$\mathrm{id}$ is a unique identifier and $S(a) \subseteq
\{\FN{00},\ldots,\FN{11}\}$ is the (non-empty) set of function classes
(Section~3, \VisionFunc{}) that instance is
authorized to perform.
\end{definition}

\begin{definition}[Workflow Graph]
\label{def:workflow-graph}
The \emph{nominal workflow graph} is a directed graph
$G = (V, E)$ where $V$ is the finite set of artifact-lifecycle phases
(Table~3.1--3.2, \VisionFlow{}):
\begin{equation*}
  \begin{aligned}
    V = \{\ & \mathrm{Plan}, \mathrm{EnvVerify}, \mathrm{StdMap}, \mathrm{Impl}, \mathrm{Review}, \\
         & \mathrm{Challenge}, \mathrm{Validate}, \mathrm{Test}, \mathrm{Doc}, \mathrm{CommitSync}, \mathrm{Gate}\ \}
  \end{aligned}
\end{equation*}
and $(u,v) \in E$ iff phase $v$ may legally begin only after phase $u$
completes (Sections~4.2, 7.1--7.2 of the same document).
\end{definition}

\begin{definition}[Artifact State Machine]
\label{def:state-machine}
The \emph{artifact state machine} is the finite automaton
$\mathcal{M} = (Q, \Sigma, \delta, q_0, F)$, exactly as diagrammed in
Figure~5.1 of \VisionFlow{}, with:
\begin{equation*}
  \begin{aligned}
    Q &= \{\mathrm{Drafted}, \mathrm{UnderReview}, \mathrm{UnderValidation}, \\
      &\phantom{{}={}} \mathrm{UnderTest}, \mathrm{Documented}, \mathrm{Gated}, \mathrm{Approved}, \mathrm{Rejected}\}, \\
    q_0 &= \mathrm{Drafted}, \qquad F = \{\mathrm{Approved}\}
  \end{aligned}
\end{equation*}
\end{definition}

\begin{definition}[Trust Tier Function]
\label{def:trust}
Let $\tau : \mathcal{C} \to \{\mathrm{T\text{-}0}, \mathrm{T\text{-}1},
\mathrm{T\text{-}2}, \mathrm{T\text{-}3}\}$ assign a trust tier to every
piece of content $c \in \mathcal{C}$ that exists in a session, per the
Trust Tier Hierarchy (Table~8.1, \VisionClean{}): T-0 model constraints, T-1 the Layer~0
constitution, T-2 the verified user's own instructions, T-3 everything
else, including every subagent return message (\S8.1.1,
\VisionGov{}).
\end{definition}

\begin{definition}[Conflict Graph]
\label{def:conflict-graph}
The \emph{function conflict graph} is the undirected graph
$C = (\{\FN{00},\ldots,\FN{11}\}, E_C)$ where $\{u,v\} \in E_C$ iff the
Non-Overlap Rule (\S5.1, \VisionFunc{}) forbids
one agent instance from performing both $u$ and $v$ on the same artifact.
By that rule, at minimum:
\begin{equation*}
  \{\FN{02},\FN{04}\},\ \{\FN{02},\FN{05}\},\ \{\FN{03},\FN{04}\},\ \{\FN{03},\FN{05}\},\ \{\FN{04},\FN{09}\} \ \in\ E_C.
\end{equation*}
\end{definition}

% ============================================================================
\subsection{Definition Validation and Structural Well-Formedness}
\label{sec:def-validation}
% WHAT: Rigorously validates every assumption, edge case, boundary condition,
%       and structural invariant embedded in
%       Definitions~\ref{def:agent}--\ref{def:conflict-graph}.
% WHY:  Per mathematical_proof_structures_standards.txt §P.1.2, all
%       hypotheses must be explicitly discharged before any theorem may
%       appeal to a definition; unvalidated assumptions constitute
%       proof gaps that would fail a third-party audit.
% HOW:  Each sub-section addresses one definition in the order it was
%       introduced, enumerating every condition to be verified.
% ============================================================================

The five definitions of Section~\ref{sec:background} carry implicit
well-formedness conditions that must be discharged before the formal
proofs of Section~\ref{sec:proofs} may appeal to them.
This section enumerates those conditions explicitly,
validates each one against the governing sources, and records the
precise boundary or edge case that each condition guards against.
No assumption is left implicit.

% ----------------------------------------------------------------------------
\subsubsection{Validation of Definition~\ref{def:agent} (Agent Instance)}
\label{sec:val-agent}
% ----------------------------------------------------------------------------

\begin{enumerate}[label=\textbf{A\arabic*.}, leftmargin=*, widest=A9]

  \item \textbf{Identifier uniqueness (boundary: single-agent system).}
        The identifier $\mathrm{id}$ must be globally unique across
        \emph{all} agent instances that are simultaneously active within
        one Engineering Studio AI session.
        \emph{Edge case:} a session containing exactly one agent instance
        trivially satisfies uniqueness; the requirement is non-trivial
        only when $\lvert\text{active instances}\rvert \geq 2$, which is
        the normal operating mode.
        \emph{Validation:} \VisionGov{} \S3.1 mandates that the
        orchestrator assigns a unique task-specification identifier to
        every subagent invocation; this identifier serves as $\mathrm{id}$.

  \item \textbf{Non-emptiness of the function-class set (boundary: null
        agent).}
        The requirement $S(a) \neq \emptyset$ excludes degenerate agents
        with no authorized function.
        \emph{Edge case:} an agent instantiated with $S(a) = \emptyset$
        could receive a task assignment but could never legally execute
        any step, making it a deadlock source.
        \emph{Validation:} \VisionFunc{} \S2 states that every agent role
        in the catalog is assigned at least one of \FN{00}--\FN{11};
        the catalog contains no ``null-role'' entries.

  \item \textbf{Function-class subset bound (boundary: all-function
        agent).}
        $S(a) \subseteq \{\FN{00},\ldots,\FN{11}\}$ is a finite
        superset of size 12; hence $S(a)$ has at most $2^{12} = 4096$
        distinct realizations.
        \emph{Edge case:} $S(a) = \{\FN{00},\ldots,\FN{11}\}$ would
        assign every function to one instance, violating the Non-Overlap
        Rule (see Definition~\ref{def:conflict-graph}).
        \emph{Validation:} The Non-Overlap Rule (\S5.1, \VisionFunc{})
        explicitly forbids multi-function assignments for all pairs in
        $E_C$; that constraint is formalized in
        Definition~\ref{def:conflict-graph} and enforced in
        Theorem~\ref{thm:multiplicity} (Section~\ref{sec:proofs}).

  \item \textbf{Statelessness of agent instances (boundary: session
        boundary).}
        An agent instance $a$ carries no persistent internal state
        across session boundaries; $a = (\mathrm{id}, S(a))$ contains
        no memory component by construction.
        \emph{Edge case:} if a resumed session reuses an old $\mathrm{id}$
        for a new agent, the uniqueness invariant (A1) is violated.
        \emph{Validation:} \VisionGov{} \S4.2 mandates statelessness and
        requires fresh identifiers on every invocation.

  \item \textbf{Layer~0 agent boundary (boundary: constitutional
        constraint).}
        The function class \FN{00} corresponds to the Layer~0 global
        constitution.  An agent with $\FN{00} \in S(a)$ is subject to
        the highest-priority constraints and \emph{must not} be
        overridden by any subagent instruction.
        \emph{Edge case:} $S(a) = \{\FN{00}\}$ produces a pure
        constitution-enforcement agent with no deliverable capability.
        This is a valid degenerate case (the guardrails agent).
        \emph{Validation:} \VisionClean{} \S2 explicitly permits a
        Layer~0-only agent role.

\end{enumerate}

% ----------------------------------------------------------------------------
\subsubsection{Validation of Definition~\ref{def:workflow-graph}
               (Workflow Graph)}
\label{sec:val-workflow}
% ----------------------------------------------------------------------------

\begin{enumerate}[label=\textbf{W\arabic*.}, leftmargin=*, widest=W9]

  \item \textbf{Vertex set completeness (boundary: missing phase).}
        $V$ must contain every phase named in Tables~3.1--3.2 of
        \VisionFlow{}.
        \emph{Edge case:} omitting even one phase (e.g., \textit{Gate})
        would make the acceptance criterion unreachable, violating the
        liveness argument in Theorem~\ref{thm:rework}.
        \emph{Validation:} The 11 elements of $V$ are listed verbatim
        from Tables~3.1--3.2; the count matches the governing source.

  \item \textbf{Well-definedness of the edge predicate (boundary: no
        edges).}
        $(u,v)\in E$ iff phase $v$ may \emph{legally} begin only after
        phase $u$ completes.  The empty-edge case ($E = \emptyset$) is
        degenerate: it yields a graph of 11 isolated vertices with no
        reachable path from any start node to \textit{Gate}.
        \emph{Validation:} Sections~4.2 and 7.1--7.2 of \VisionFlow{}
        give an explicit ordered phase sequence; at minimum the chain
        $\mathrm{Plan} \to \cdots \to \mathrm{Gate}$ is
        mandated, so $E \neq \emptyset$.

  \item \textbf{Finiteness (boundary: infinite graph).}
        $|V| = 11$ and $|E| \leq |V|^2 = 121$; both are finite.
        \emph{Edge case:} dynamic phase insertion during a run is
        \emph{not} permitted because the vertex set is fixed at session
        start (\VisionFlow{} \S4.1).  Theorem~\ref{thm:acyclic}
        relies on finiteness for its topological-sort argument.

  \item \textbf{Determinism of phase ordering (boundary: concurrent
        start).}
        Two phases $u, v \in V$ with neither $(u,v) \in E$ nor
        $(v,u) \in E$ may execute concurrently.
        \emph{Edge case:} if the governing source specifies a
        \emph{strict} total order on all phases, then no concurrency
        exists and $G$ is a Hamiltonian path.
        \emph{Validation:} \VisionFlow{} \S7.2 explicitly permits
        parallel specialist work within the \textit{Impl} phase, so
        $G$ is a DAG but not necessarily a path.

  \item \textbf{Single entry point (boundary: multiple roots).}
        $G$ must have exactly one source vertex (in-degree zero),
        namely \textit{Plan}.
        \emph{Edge case:} a second source vertex would allow a workflow
        run that bypasses planning, violating the decomposition
        invariant.
        \emph{Validation:} \VisionFlow{} \S4.2 names \textit{Plan} as
        the mandatory first phase with no prerequisites.

  \item \textbf{Single terminal phase (boundary: multiple sinks).}
        $G$ must have exactly one sink vertex (out-degree zero), namely
        \textit{Gate}.
        \emph{Edge case:} multiple sinks would allow a run to terminate
        without passing the quality gate, violating the approval
        invariant.
        \emph{Validation:} \VisionFlow{} \S7.2 names \textit{Gate} as
        the sole phase whose completion constitutes workflow acceptance.

\end{enumerate}

% ----------------------------------------------------------------------------
\subsubsection{Validation of Definition~\ref{def:state-machine}
               (Artifact State Machine)}
\label{sec:val-state-machine}
% ----------------------------------------------------------------------------

\begin{enumerate}[label=\textbf{S\arabic*.}, leftmargin=*, widest=S9]

  \item \textbf{State set completeness (boundary: missing state).}
        $Q$ must contain all states named in Figure~5.1 of
        \VisionFlow{}.
        \emph{Validation:} The eight elements of $Q$ are transcribed
        verbatim from that figure:
        \begin{align*}
          Q = \{\ &\mathrm{Drafted},\ \mathrm{UnderReview},\
                   \mathrm{UnderValidation},\ \mathrm{UnderTest}, \\
                  &\mathrm{Documented},\ \mathrm{Gated},\
                   \mathrm{Approved},\ \mathrm{Rejected}\ \}.
        \end{align*}
        Count: $|Q| = 8$.  No state from the governing figure is
        omitted.

  \item \textbf{Initial state membership (boundary: $q_0 \notin Q$).}
        $q_0 = \mathrm{Drafted} \in Q$; confirmed by inspection of
        the set above.
        \emph{Edge case:} if $q_0 \notin Q$, the automaton would be
        ill-formed and no trace could begin.

  \item \textbf{Accepting-state membership (boundary: $F \not\subseteq
        Q$).}
        $F = \{\mathrm{Approved}\} \subseteq Q$; confirmed by
        inspection.
        \emph{Edge case:} $F = \emptyset$ would make every trace
        non-accepting, meaning no artifact could ever be formally
        approved.

  \item \textbf{Reachability of the accepting state (boundary:
        unreachable $F$).}
        There must exist at least one sequence of transitions
        $\sigma \in \Sigma^{*}$ such that
        $\delta^{*}(q_0, \sigma) \in F$.
        \emph{Validation:} Figure~5.1 of \VisionFlow{} shows the
        path
        \begin{align*}
          \mathrm{Drafted}
            &\xrightarrow{\,\sigma_1\,} \mathrm{UnderReview}
             \xrightarrow{\,\sigma_2\,} \mathrm{UnderValidation}
             \xrightarrow{\,\sigma_3\,} \mathrm{UnderTest} \\
            &\xrightarrow{\,\sigma_4\,} \mathrm{Documented}
             \xrightarrow{\,\sigma_5\,} \mathrm{Gated}
             \xrightarrow{\,\sigma_6\,} \mathrm{Approved},
        \end{align*}
        confirming that $\mathrm{Approved}$ is reachable.

  \item \textbf{Role of the \textit{Rejected} state (boundary:
        absorbing sink).}
        $\mathrm{Rejected} \notin F$ and $\mathrm{Rejected}$ is a
        terminal state from which no transition leads back to
        $\mathrm{Approved}$ \emph{within the same artifact
        instance}.
        \emph{Edge case:} if a transition from $\mathrm{Rejected}$
        to $\mathrm{Drafted}$ existed, the automaton would support
        in-place revision; this is not modeled here (a new artifact
        instance is created instead, per \VisionFlow{} \S5.2).
        \emph{Validation:} Figure~5.1 shows \textit{Rejected} as a
        sink; the governing source does not depict any outgoing edge.

  \item \textbf{Determinism of $\delta$ (boundary: nondeterministic
        transition).}
        The automaton $\mathcal{M}$ is defined as a finite automaton;
        $\delta : Q \times \Sigma \to Q$ is a \emph{total} function.
        \emph{Edge case:} partial $\delta$ (undefined on some
        $(q, \sigma)$ pair) would require a ``trap state'' or
        explicit error handling.
        \emph{Validation:} \VisionFlow{} \S5.1 specifies a transition
        for every valid event from every reachable state; states not
        depicted as having an outgoing edge for a given symbol implicitly
        self-loop or are unreachable under that symbol, maintaining
        totality.

\end{enumerate}

% ----------------------------------------------------------------------------
\subsubsection{Validation of Definition~\ref{def:trust}
               (Trust Tier Function)}
\label{sec:val-trust}
% ----------------------------------------------------------------------------

\begin{enumerate}[label=\textbf{T\arabic*.}, leftmargin=*, widest=T9]

  \item \textbf{Domain completeness (boundary: content outside
        $\mathcal{C}$).}
        $\tau$ is defined on \emph{every} piece of content
        $c \in \mathcal{C}$ present in a session.
        \emph{Edge case:} content injected after session initialization
        (e.g., a tool-return payload) must be assigned a tier before
        it is processed.
        \emph{Validation:} \VisionGov{} \S8.1.1 states that all
        subagent return messages are classified T-3 by default,
        providing a catch-all assignment that ensures $\tau$ is total.

  \item \textbf{Tier ordering and strict precedence (boundary:
        tier inversion).}
        The tiers are totally ordered:
        $\mathrm{T\text{-}0} \succ \mathrm{T\text{-}1}
         \succ \mathrm{T\text{-}2} \succ \mathrm{T\text{-}3}$.
        A conflict between content at tiers $t_i$ and $t_j$ with
        $t_i \succ t_j$ is always resolved in favor of $t_i$.
        \emph{Edge case:} two pieces of content at the \emph{same}
        tier that conflict (intra-tier conflict) are not automatically
        resolved by $\tau$ alone; such conflicts must be escalated to
        the Layer~0 constitution (\VisionClean{} \S8.2).
        \emph{Validation:} Table~8.1 of \VisionClean{} defines the
        four-tier hierarchy with explicit precedence.

  \item \textbf{Immutability of T-0 constraints (boundary: T-0
        override attempt).}
        No content at any tier $\mathrm{T\text{-}1}$ through
        $\mathrm{T\text{-}3}$ may override a T-0 model constraint.
        \emph{Edge case:} a T-2 user instruction that attempts to
        disable a T-0 safety constraint must be silently dropped and
        logged, not executed.
        \emph{Validation:} \VisionClean{} \S2 (Layer~0 mandate) and
        \VisionGov{} \S8.1 both confirm this invariant explicitly.

  \item \textbf{Session-scoped validity (boundary: cross-session
        reuse).}
        $\tau$ is defined per session; a content item from a prior
        session carries no tier assignment into a new session.
        \emph{Validation:} statelessness of agent instances
        (A4 above) implies no cross-session content persistence;
        every new session initializes $\mathcal{C}$ fresh.

  \item \textbf{Codomain exhaustiveness (boundary: untiered
        content).}
        The codomain
        $\{\mathrm{T\text{-}0}, \mathrm{T\text{-}1},
         \mathrm{T\text{-}2}, \mathrm{T\text{-}3}\}$
        has exactly four elements; every content item receives
        exactly one tier (function, not relation).
        \emph{Edge case:} a content item that matches multiple tier
        criteria simultaneously (e.g., a user instruction that also
        embeds a model constraint) is assigned the \emph{highest}
        applicable tier to preserve safety.
        \emph{Validation:} \VisionClean{} \S8.1 specifies
        highest-tier-wins as the tie-breaking rule.

\end{enumerate}

% ----------------------------------------------------------------------------
\subsubsection{Validation of Definition~\ref{def:conflict-graph}
               (Conflict Graph)}
\label{sec:val-conflict}
% ----------------------------------------------------------------------------

\begin{enumerate}[label=\textbf{C\arabic*.}, leftmargin=*, widest=C9]

  \item \textbf{Vertex set completeness (boundary: missing function
        class).}
        The vertex set of $C$ is
        $\{\FN{00}, \FN{01}, \ldots, \FN{11}\}$, exactly 12 elements,
        one per function class defined in \VisionFunc{} \S2.
        \emph{Edge case:} if a new function class \FN{12} were added
        to the governing catalog, $C$ must be updated and all conflict
        pairs re-evaluated.
        \emph{Validation:} the catalog in \VisionFunc{} \S2 enumerates
        exactly 12 function classes; the vertex count matches.

  \item \textbf{Undirectedness and irreflexivity (boundary: self-loops
        and directed edges).}
        $C$ is an \emph{undirected} graph:
        $\{u,v\} \in E_C \Rightarrow \{v,u\} \in E_C$, and
        $\{u,u\} \notin E_C$ for all $u$ (a single agent cannot
        ``conflict with itself'' in the non-overlap sense).
        \emph{Validation:} the Non-Overlap Rule (\S5.1, \VisionFunc{})
        is stated symmetrically (``functions $u$ and $v$ may not
        coexist on the same artifact''); self-conflict is
        semantically vacuous.

  \item \textbf{Minimality vs.\ completeness of the stated edge set
        (boundary: under-specification).}
        Definition~\ref{def:conflict-graph} states that
        \emph{at minimum} the five listed pairs belong to $E_C$.
        This is a lower bound, not an exhaustive enumeration.
        The full $E_C$ may contain additional pairs mandated by
        \VisionFunc{} \S5.1.
        \emph{Edge case:} if proofs require only the five listed
        pairs, they remain valid even if $E_C$ is strictly larger,
        because adding edges can only \emph{strengthen} the
        non-overlap guarantees.
        \emph{Validation:} Theorem~\ref{thm:multiplicity} in
        Section~\ref{sec:proofs} is stated with respect to
        $E_C$ in full generality, not merely its lower bound.

  \item \textbf{Explicit enumeration of the minimum edge set.}
        For complete auditability the five minimum pairs are restated
        individually:
        \begin{enumerate}[label=(\alph*), leftmargin=*]
          \item $\{\FN{02}, \FN{04}\}$: Domain Specialist and
                Reviewer may not share an artifact instance.
          \item $\{\FN{02}, \FN{05}\}$: Domain Specialist and
                Validator may not share an artifact instance.
          \item $\{\FN{03}, \FN{04}\}$: Micro Specialist and
                Reviewer may not share an artifact instance.
          \item $\{\FN{03}, \FN{05}\}$: Micro Specialist and
                Validator may not share an artifact instance.
          \item $\{\FN{04}, \FN{09}\}$: Reviewer and Adversarial
                Challenger may not share an artifact instance.
        \end{enumerate}
        \emph{Rationale:} items (a)--(d) implement separation of
        construction from critique; item (e) implements separation
        of structured review from adversarial stress-testing.
        Both rationales are grounded in \VisionFunc{} \S5.1.

  \item \textbf{Layer~0 vertex isolation (boundary: conflicting
        the constitution).}
        $\FN{00}$ represents the global constitution; it has no
        deliverable-level conflict with any other function class
        because it operates at a meta-level above artifact
        authorship.
        \emph{Edge case:} one might ask whether
        $\{\FN{00}, \FN{04}\}$ should be in $E_C$.
        The governing source does not place \FN{00} in a
        deliverable-authorship role, so no conflict arises on
        a per-artifact basis.
        \emph{Validation:} \VisionFunc{} \S2 describes \FN{00} as
        a constitutional constraint layer, not an artifact producer;
        therefore the Non-Overlap Rule's ``same artifact'' predicate
        never fires for \FN{00}.

  \item \textbf{Independence of non-adjacent function classes
        (boundary: implicit conflict).}
        Function-class pairs $\{u,v\} \notin E_C$ are explicitly
        \emph{permitted} to coexist on the same artifact.
        For example, $\{\FN{06}, \FN{07}\}$ (Documentation,
        Testing) is not in the minimum $E_C$ and is not prohibited
        by \VisionFunc{} \S5.1.
        \emph{Validation:} absence of an edge in $C$ is an
        affirmative assertion of compatibility, not merely silence;
        this interpretation is consistent with the closed-world
        assumption used in \VisionFunc{} \S5.1.

\end{enumerate}

% ----------------------------------------------------------------------------
\subsubsection{Cross-Definition Consistency Checks}
\label{sec:val-cross}
% ----------------------------------------------------------------------------

The five definitions must be mutually consistent.
The following checks confirm there are no contradictions at
their interfaces.

\begin{enumerate}[label=\textbf{X\arabic*.}, leftmargin=*, widest=X9]

  \item \textbf{Agent functions vs.\ conflict graph (interface:
        Definitions~\ref{def:agent} and~\ref{def:conflict-graph}).}
        For every agent instance $a$ and every pair
        $\{f_i, f_j\} \in E_C$, the assignment
        $\{f_i, f_j\} \subseteq S(a)$ is forbidden.
        This is a direct restatement of the Non-Overlap Rule in
        terms of the two definitions.
        \emph{Consistency:} both definitions use the same
        function-class labels \FN{00}--\FN{11}; no label clash
        or mismatch exists.

  \item \textbf{Workflow phases vs.\ artifact states (interface:
        Definitions~\ref{def:workflow-graph}
        and~\ref{def:state-machine}).}
        Each phase $v \in V$ must be associated with at least one
        artifact state transition in $\mathcal{M}$.
        \begin{align*}
          \mathrm{Review}\ &\leftrightarrow\
            \mathrm{Drafted} \to \mathrm{UnderReview} \\
          \mathrm{Validate}\ &\leftrightarrow\
            \mathrm{UnderReview} \to \mathrm{UnderValidation} \\
          \mathrm{Test}\ &\leftrightarrow\
            \mathrm{UnderValidation} \to \mathrm{UnderTest} \\
          \mathrm{Doc}\ &\leftrightarrow\
            \mathrm{UnderTest} \to \mathrm{Documented} \\
          \mathrm{Gate}\ &\leftrightarrow\
            \mathrm{Gated} \to \mathrm{Approved}\ \text{or}\
            \mathrm{Gated} \to \mathrm{Rejected}
        \end{align*}
        \emph{Consistency:} every accepting or rejecting event in
        $\mathcal{M}$ corresponds to a phase transition in $G$;
        neither definition introduces states or phases that are
        unreferenced in the other.

  \item \textbf{Trust tiers vs.\ workflow phases (interface:
        Definitions~\ref{def:trust}
        and~\ref{def:workflow-graph}).}
        $\tau$ applies to content produced or consumed at
        \emph{every} phase $v \in V$.
        \emph{Edge case:} at the \textit{Plan} phase, the primary
        input is the user's product brief, which is T-2 content.
        At \textit{Gate}, the quality-gate agent's decision message
        is T-3 content (subagent return), but the Layer~0
        constitutional check that precedes acceptance is T-0.
        \emph{Consistency:} $\tau$ is defined on all of
        $\mathcal{C}$; since content exists at every phase,
        $\tau$ is always applicable.

  \item \textbf{Agent functions vs.\ workflow phases (interface:
        Definitions~\ref{def:agent}
        and~\ref{def:workflow-graph}).}
        Every phase $v \in V$ must be executable by at least one
        function class $f \in \{\FN{00},\ldots,\FN{11}\}$, and
        every function class must be invokable during at least one
        phase.
        \emph{Validation:}
        \begin{align*}
          \mathrm{Plan}\ &\mapsto\ \FN{01}\ \text{(Orchestrator)} \\
          \mathrm{EnvVerify}\ &\mapsto\ \FN{10}\ \text{(Guardrails)} \\
          \mathrm{StdMap}\ &\mapsto\ \FN{01}, \FN{11}
            \ \text{(Orchestrator, Utility)} \\
          \mathrm{Impl}\ &\mapsto\ \FN{02}, \FN{03}
            \ \text{(Domain, Micro Specialist)} \\
          \mathrm{Review}\ &\mapsto\ \FN{04}\ \text{(Reviewer)} \\
          \mathrm{Challenge}\ &\mapsto\ \FN{09}
            \ \text{(Adversarial)} \\
          \mathrm{Validate}\ &\mapsto\ \FN{05}
            \ \text{(Validator)} \\
          \mathrm{Test}\ &\mapsto\ \FN{07}\ \text{(Testing)} \\
          \mathrm{Doc}\ &\mapsto\ \FN{06}
            \ \text{(Documentation)} \\
          \mathrm{CommitSync}\ &\mapsto\ \FN{01}, \FN{11} \\
          \mathrm{Gate}\ &\mapsto\ \FN{08}
            \ \text{(Quality Gate)}
        \end{align*}
        \FN{00} (Layer~0 constitution) is active across all phases
        as a meta-level constraint, not a phase-specific actor.
        \emph{Consistency:} all 12 function classes participate in
        at least one phase; all 11 phases have at least one
        responsible function class.

  \item \textbf{Finiteness across all definitions (boundary:
        infinite structures).}
        All five definitions produce finite mathematical objects:
        \begin{enumerate}[label=(\alph*), leftmargin=*]
          \item $|S(a)| \leq 12$ (agent function set);
          \item $|V| = 11$, $|E| \leq 121$ (workflow graph);
          \item $|Q| = 8$, $|\Sigma|$ finite by \VisionFlow{}
                \S5.1 (state machine);
          \item $|\text{codomain}(\tau)| = 4$ (trust tiers);
          \item $|\text{vertices of }C| = 12$,
                $|E_C| \leq 66$ (conflict graph, $\binom{12}{2}$).
        \end{enumerate}
        \emph{Consistency:} finiteness is preserved across all
        definitions; no definition introduces an unbounded structure
        that could invalidate the termination argument.

\end{enumerate}

% ============================================================================
\section{System Architecture (Summary)}
\label{sec:architecture}

Table~\ref{tab:layers} summarizes the Layer~0--8 organizational model and
the Challenge Division that this paper's proofs are stated over; full
mission statements are governed by \VisionClean{} and are not re-derived here, per the token-efficiency
standard already applied throughout this vision family.

\begin{table}[H]
\centering
\caption{Layer/Division to Function-Class Summary.}
\label{tab:layers}
\begin{tabular}{@{}llp{6.5cm}@{}}
\toprule
\textbf{Layer/Division} & \textbf{Function} & \textbf{Role} \\
\midrule
Layer 0 & \FN{00} & Global constitution; no deliverables \\
Layer 1 & \FN{01} & Orchestrator: decompose, assign, merge \\
Layer 2/3 & \FN{02}/\FN{03} & Domain / Micro Specialists: build \\
Layer 4 & \FN{04} & Reviewer: critique against one standard \\
Layer 5 & \FN{05} & Validator: detect contradiction, never fix \\
Layer 6 & \FN{06} & Documentation \\
Layer 7 & \FN{07} & Testing \\
Layer 8 & \FN{08} & Quality Gate: sole approval authority \\
Challenge Div. & \FN{09} & Adversarial: "how does this fail?" \\
Guardrails & \FN{10} & Cross-cutting mandate enforcement \\
Utility & \FN{11} & Shared, read-only support capability \\
\bottomrule
\end{tabular}
\end{table}

\section{Formal Model of the Workflow}
\label{sec:model}

Combining Definitions~\ref{def:workflow-graph}--\ref{def:conflict-graph}, an
Engineering Studio AI run over one deliverable is the tuple
$(G, \mathcal{M}, \tau, C, R, F_{\max}, D_{\max})$ where $R$ is the
per-phase-gate retry limit, $F_{\max}$ is the per-turn fan-out ceiling, and
$D_{\max} = 2$ is the nesting-depth limit (\S5.1.1,
\VisionGov{}). Section
\ref{sec:proofs} proves properties of this tuple.

% ============================================================================
% FORMAL MODEL EXPANSION
% ============================================================================

\medskip
\noindent\textbf{Notation.}
Throughout this section the following symbols are fixed and refer
exclusively to the objects introduced in
Definitions~\ref{def:workflow-graph}--\ref{def:conflict-graph}.
No symbol is reused with a different meaning.

\begin{enumerate}
  \item $G = (V, E)$: the workflow graph
        (Definition~\ref{def:workflow-graph}).
  \item $\mathcal{M} = (Q, \Sigma, \delta, q_0, F)$: the artifact
        state machine (Definition~\ref{def:state-machine}).
  \item $\tau : \mathrm{Content} \to
        \{\mathrm{T\text{-}1},\mathrm{T\text{-}2},\mathrm{T\text{-}3}\}$:
        the trust-tier assignment function
        (Definition~\ref{def:trust}).
  \item $C = (V_C, E_C)$: the conflict graph
        (Definition~\ref{def:conflict-graph}).
  \item $R \in \mathbb{N}_0$: per-phase-gate retry limit
        (\S8.2.1, \VisionGov{}).
  \item $F_{\max} \in \mathbb{N}$: per-turn fan-out ceiling
        (\S5.1.1, \S5.2.3, \VisionGov{}).
  \item $D_{\max} = 2$: nesting-depth limit
        (\S5.1.1, \VisionGov{}).
  \item $P \coloneqq |V| = 11$: number of workflow phases
        (Table~3.2, \VisionFlow{}).
  \item $B \coloneqq R \cdot P$: global pooled rework budget.
\end{enumerate}

\medskip
\noindent\textbf{Definition of the Run Tuple.}
An Engineering Studio AI run over a single deliverable is the
seven-tuple
\begin{align}
  \label{eq:run-tuple}
  \mathcal{R}
  \;\coloneqq\;
  \bigl(
    G,\;
    \mathcal{M},\;
    \tau,\;
    C,\;
    R,\;
    F_{\max},\;
    D_{\max}
  \bigr).
\end{align}

\noindent The individual components and their domains are defined
and validated in the following enumeration.  Every assumption is
stated explicitly; no component is left undefined.

\begin{enumerate}
  % ------------------------------------------------------------------
  \item \textbf{Component 1: Workflow Graph $G = (V,E)$.}
  \begin{enumerate}
    \item[1.1] \emph{Definition.}
               $G$ is the directed graph of
               Definition~\ref{def:workflow-graph}.  Its vertex set
               is the finite, ordered enumeration
               \begin{align}
                 V \;=\;
                 \{&\mathrm{Plan},\;
                   \mathrm{EnvVerify},\;
                   \mathrm{StdMap},\;
                   \mathrm{Impl},\notag\\
                 &\mathrm{Review},\;
                   \mathrm{Challenge},\;
                   \mathrm{Validate},\;
                   \mathrm{Test},\notag\\
                 &\mathrm{Doc},\;
                   \mathrm{CommitSync},\;
                   \mathrm{Gate}\}. \label{eq:V-def}
               \end{align}
    \item[1.2] \emph{Cardinality.}
               Counting the elements of \eqref{eq:V-def}:
               $|V| = 11 = P$.
    \item[1.3] \emph{Edge set.}
               $E$ is exactly the set of ``must-complete-before''
               constraints from Table~3.2 and the join-point
               constraints from \S7.2 of \VisionFlow{}.
               No edge is introduced by this paper beyond those listed
               in those sources; no edge is omitted.
    \item[1.4] \emph{Acyclicity.}
               $G$ is a directed acyclic graph (DAG).
               Proof: Theorem~\ref{thm:acyclic} (Section
               \ref{sec:proofs}).
    \item[1.5] \emph{Finiteness.}
               $|V| = 11 < \infty$ and
               $|E| \leq |V|^2 = 121 < \infty$.
               Both bounds are finite; no infinite structure is
               introduced.
    \item[1.6] \emph{Boundary conditions.}
               An empty graph ($|V|=0$ or $|E|=\emptyset$) would
               still be a valid DAG.  The actual $G$ has $|V|=11$
               and at least the ten edges corresponding to the main
               sequential chain
               (Plan $\to$ \ldots $\to$ Gate), so $|E| \geq 10$.
  \end{enumerate}

  % ------------------------------------------------------------------
  \item \textbf{Component 2: Artifact State Machine
        $\mathcal{M} = (Q,\Sigma,\delta,q_0,F)$.}
  \begin{enumerate}
    \item[2.1] \emph{Definition.}
               $\mathcal{M}$ is the finite deterministic automaton of
               Definition~\ref{def:state-machine}.
    \item[2.2] \emph{State set.}
               $Q$ is the finite set of artifact processing states;
               its elements are enumerated in
               Definition~\ref{def:state-machine}.  Concretely,
               \begin{align}
                 Q \;=\;
                 \{&\mathrm{Drafted},\;
                   \mathrm{UnderReview},\;
                   \mathrm{UnderChallenge},\notag\\
                 &\mathrm{UnderValidation},\;
                   \mathrm{UnderTest},\notag\\
                 &\mathrm{Approved},\;
                   \mathrm{Rejected},\;
                   \mathrm{Escalated}\},
                   \label{eq:Q-def}
               \end{align}
               giving $|Q| = 8$.
    \item[2.3] \emph{Input alphabet.}
               $\Sigma$ is the finite set of gate-decision symbols.
               At minimum $\Sigma$ contains the symbols
               $\{\mathtt{approve},\,\mathtt{reject}\}$;
               additional symbols may be introduced by the governance
               document without affecting the bound proofs, provided
               $|\Sigma| < \infty$.
    \item[2.4] \emph{Transition function.}
               $\delta : Q \times \Sigma \to Q$ is the deterministic
               transition function.  Determinism ensures that at every
               state $q \in Q$ and every input $\sigma \in \Sigma$,
               exactly one successor state $\delta(q,\sigma)$ is
               defined.  No non-determinism is introduced.
    \item[2.5] \emph{Initial state.}
               $q_0 = \mathrm{Drafted} \in Q$.
    \item[2.6] \emph{Accepting states.}
               $F = \{\mathrm{Approved},\mathrm{Escalated}\} \subseteq Q$.
               Both terminal states represent run completion:
               $\mathrm{Approved}$ denotes successful certification by
               the Quality Gate; $\mathrm{Escalated}$ denotes
               circuit-breaker activation (\S5.3.2,
               \VisionGov{}).
    \item[2.7] \emph{Finiteness.}
               $|Q| = 8 < \infty$; $|\Sigma| < \infty$ (by
               Definition~\ref{def:state-machine} and the governance
               document); $|\delta| = |Q \times \Sigma| < \infty$.
    \item[2.8] \emph{Termination.}
               Every run terminates within at most $P + B$
               state-machine transitions.
               Proof: Theorem~\ref{thm:rework}
               (Section~\ref{sec:proofs}).
  \end{enumerate}

  % ------------------------------------------------------------------
  \item \textbf{Component 3: Trust-Tier Function $\tau$.}
  \begin{enumerate}
    \item[3.1] \emph{Definition.}
               $\tau$ is defined in Definition~\ref{def:trust}.  Its
               domain is the set of all content items produced or
               consumed during a session; its codomain is the
               three-element trust lattice:
               \begin{align}
                 \mathrm{T\text{-}1}
                 \;<\; \mathrm{T\text{-}2}
                 \;<\; \mathrm{T\text{-}3},
                 \label{eq:trust-order}
               \end{align}
               where ``$<$'' denotes decreasing trust (T-1 is most
               trusted; T-3 is least trusted).
    \item[3.2] \emph{Assignment rules.}
               The three tiers are assigned as follows:
               \begin{enumerate}
                 \item[(T-1)] System-prompt and hard-coded constraints
                              (\S8.1.1, \VisionGov{}).
                 \item[(T-2)] User-supplied requests (same section).
                 \item[(T-3)] Every subagent return message at the
                              moment of production (same section).
               \end{enumerate}
    \item[3.3] \emph{Re-tagging restriction.}
               Only \FN{08} (Quality Gate) may re-tag content to a
               lower (more trusted) tier.  No other agent may alter
               $\tau(c)$ downward.  Authority: Table~8.1,
               \VisionClean{}.
    \item[3.4] \emph{Invariant.}
               Every subagent-originated content item $c$ satisfies
               $\tau(c) = \mathrm{T\text{-}3}$ throughout a session
               unless it has been the subject of an explicit
               \FN{08} $\Approved$ certificate event.
               Proof: Theorem~\ref{thm:trust}
               (Section~\ref{sec:proofs}).
    \item[3.5] \emph{Well-definedness.}
               $\tau$ is a total function: every content item receives
               a tier assignment at the moment of production (H2 of
               Theorem~\ref{thm:trust}).  No content item is left
               unclassified.
  \end{enumerate}

  % ------------------------------------------------------------------
  \item \textbf{Component 4: Conflict Graph $C = (V_C, E_C)$.}
  \begin{enumerate}
    \item[4.1] \emph{Definition.}
               $C$ is the undirected conflict graph of
               Definition~\ref{def:conflict-graph}.  Its vertex set
               is the full function-class catalog:
               \begin{align}
                 V_C
                 \;=\;
                 \{\FN{00},\,\FN{01},\,\FN{02},\,\ldots,\,\FN{11}\},
                 \quad
                 |V_C| = 12.
                 \label{eq:VC-def}
               \end{align}
    \item[4.2] \emph{Edge semantics.}
               $\{u,v\} \in E_C$ iff assigning both function classes
               $u$ and $v$ to the same agent instance violates the
               Non-Overlap Rule (\S5.1, \VisionFunc{}) or the
               no-self-certification clause (\S1.2.3, \VisionClean{}).
    \item[4.3] \emph{Mandatory clique.}
               The four mandatory-phase function classes
               \begin{align}
                 K
                 \;=\;
                 \{\FN{02},\;\FN{04},\;\FN{05},\;\FN{09}\}
                 \label{eq:K-def}
               \end{align}
               form a $4$-clique in $C$: every pair of distinct
               elements of $K$ is an edge of $C$
               (Stage~II of Theorem~\ref{thm:multiplicity}).
    \item[4.4] \emph{Minimum multiplicity.}
               Any compliant deployment covering $K$ requires at least
               $|K| = 4$ distinct agent instances.
               Proof: Theorem~\ref{thm:multiplicity}
               (Section~\ref{sec:proofs}).
    \item[4.5] \emph{Finiteness.}
               $|V_C| = 12 < \infty$;
               $|E_C| \leq \binom{12}{2} = 66 < \infty$.
  \end{enumerate}

  % ------------------------------------------------------------------
  \item \textbf{Component 5: Per-Phase Retry Limit $R$.}
  \begin{enumerate}
    \item[5.1] \emph{Domain.}
               $R \in \mathbb{N}_0 = \{0,1,2,\ldots\}$.
               The value $R = 0$ is admissible: it means the first
               rejection immediately triggers the circuit-breaker.
    \item[5.2] \emph{Authority.}
               Declared at \S8.2.1, \VisionGov{}.
               This paper treats $R$ as an arbitrary but fixed element
               of $\mathbb{N}_0$; no numeric value is assumed beyond
               finiteness.
    \item[5.3] \emph{Derived budget.}
               The global pooled rework budget is
               \begin{align}
                 B \;\coloneqq\; R \cdot P
                 \;=\; R \cdot 11
                 \;\in\; \mathbb{N}_0.
                 \label{eq:B-def}
               \end{align}
    \item[5.4] \emph{Budget exhaustion.}
               Once $b_i = B$ rejection events have occurred in a run,
               no further rejection is processed; the escalation
               mechanism fires (\S5.3.2, \VisionGov{}).
               This is the sole mechanism by which an artifact exits
               the run in the $\mathrm{Escalated}$ state.
    \item[5.5] \emph{Boundary cases.}
               \begin{enumerate}
                 \item[(a)] $R = 0$: $B = 0$; the first rejection is
                            immediately escalated.
                 \item[(b)] $R \to \infty$: excluded by H1 of
                            Theorem~\ref{thm:rework}, which requires
                            $R \in \mathbb{N}_0$ (finite).  The model
                            is not defined for $R = \infty$.
               \end{enumerate}
  \end{enumerate}

  % ------------------------------------------------------------------
  \item \textbf{Component 6: Fan-Out Ceiling $F_{\max}$.}
  \begin{enumerate}
    \item[6.1] \emph{Domain.}
               $F_{\max} \in \mathbb{N} = \{1,2,3,\ldots\}$.
               In particular $F_{\max} \geq 1$; zero fan-out is
               excluded because an Orchestrator that never invokes
               subagents cannot drive the workflow.
    \item[6.2] \emph{Authority.}
               Declared at \S5.1.1 and \S5.2.3,
               \VisionGov{}.
               No specific numeric value is assigned by the governance
               document; the proofs treat $F_{\max}$ as an arbitrary
               but fixed element of $\mathbb{N}$.
    \item[6.3] \emph{Per-turn bound.}
               At any single Orchestrator turn $t$, the number of
               distinct subagent invocations satisfies
               \begin{align}
                 \mathrm{invocations}(t) \;\leq\; F_{\max}.
                 \label{eq:fmax-bound}
               \end{align}
    \item[6.4] \emph{Interaction with $D_{\max}$.}
               The per-turn bound \eqref{eq:fmax-bound} applies to
               the \emph{direct} children of the Orchestrator in the
               invocation tree.  Because $D_{\max} = 2$, no subagent
               invokes further subagents; the total invocations per
               turn equals exactly the direct-children count, which is
               bounded by $F_{\max}$.
    \item[6.5] \emph{Boundary cases.}
               \begin{enumerate}
                 \item[(a)] $F_{\max} = 1$: serial execution; each
                            turn invokes at most one subagent.
                 \item[(b)] Saturation at $F_{\max}$: admissible;
                            represents the worst-case parallel
                            fan-out.
               \end{enumerate}
  \end{enumerate}

  % ------------------------------------------------------------------
  \item \textbf{Component 7: Nesting-Depth Limit $D_{\max}$.}
  \begin{enumerate}
    \item[7.1] \emph{Value.}
               $D_{\max} = 2$ (a specific, numerically fixed constant,
               not a free parameter).
    \item[7.2] \emph{Authority.}
               Declared at \S5.1.1, \S5.1.2,
               \VisionGov{}.
    \item[7.3] \emph{Semantics.}
               The invocation hierarchy has at most two levels:
               \begin{enumerate}
                 \item[(a)] Depth~0: the Orchestrator (\FN{01}).
                 \item[(b)] Depth~1: subagents directly invoked by the
                            Orchestrator (any \FN{02}--\FN{11}).
               \end{enumerate}
               A depth-1 subagent may \emph{not} itself invoke further
               subagents.  Formally, for any invocation event $e$ at
               depth $d$:
               \begin{align}
                 d \;\leq\; D_{\max} - 1 \;=\; 1.
                 \label{eq:depth-bound}
               \end{align}
    \item[7.4] \emph{Consequence.}
               \eqref{eq:depth-bound} ensures that the invocation tree
               rooted at any Orchestrator turn has height exactly~1
               (root plus one level of leaves).  This is the key
               structural assumption exploited in Stage~I of the proof
               of Theorem~5 (Section~\ref{sec:sketch}).
    \item[7.5] \emph{Boundary cases.}
               \begin{enumerate}
                 \item[(a)] $D_{\max} = 1$: the Orchestrator itself
                            is the only active agent; no subagent
                            invocations are permitted.  The bound
                            $N \leq F_{\max} \cdot P \cdot (1+R)$
                            still holds (vacuously, with
                            $\mathrm{invocations}(t) = 0$).
                 \item[(b)] $D_{\max} > 2$: outside the scope of this
                            model; the governance document does not
                            currently permit deeper nesting.
               \end{enumerate}
  \end{enumerate}
\end{enumerate}

\medskip
\noindent\textbf{Parameter Summary and Domain Table.}

\begin{table}[H]
\centering
\caption{Run-Tuple Parameters, Domains, and Governing References.}
\label{tab:params}
\begin{tabular}{@{}l l l p{4.0cm}@{}}
\toprule
\textbf{Symbol} & \textbf{Domain} & \textbf{Value / Constraint}
  & \textbf{Authority} \\
\midrule
$P$
  & $\mathbb{N}$
  & $11$ (fixed)
  & Table~3.2, \VisionFlow{} \\
$R$
  & $\mathbb{N}_0$
  & $\geq 0$ (free parameter)
  & \S8.2.1, \VisionGov{} \\
$B$
  & $\mathbb{N}_0$
  & $= R \cdot P$
  & Derived from $R$, $P$ \\
$F_{\max}$
  & $\mathbb{N}$
  & $\geq 1$ (free parameter)
  & \S5.1.1, \VisionGov{} \\
$D_{\max}$
  & $\mathbb{N}$
  & $2$ (fixed)
  & \S5.1.1, \VisionGov{} \\
$|V|$
  & $\mathbb{N}$
  & $11$ (fixed)
  & Definition~\ref{def:workflow-graph} \\
$|V_C|$
  & $\mathbb{N}$
  & $12$ (fixed)
  & Definition~\ref{def:conflict-graph} \\
\bottomrule
\end{tabular}
\end{table}

\medskip
\noindent\textbf{Derived Bounds (closed-form, from the tuple).}

The following bounds are derived directly from the tuple components
without any additional assumptions.  Each derivation is fully
explicit.

\begin{enumerate}
  \item \emph{Maximum state-machine transitions per run:}
        \begin{align}
          \text{transitions} \;\leq\; P + B
          \;=\; P + R \cdot P
          \;=\; P(1 + R)
          \;=\; 11(1 + R).
          \label{eq:derived-transitions}
        \end{align}
        Justification: Theorem~\ref{thm:rework}, proved in
        Section~\ref{sec:proofs}.

  \item \emph{Maximum total subagent invocations per run:}
        \begin{align}
          N
          \;\leq\;
          F_{\max} \cdot (P + B)
          \;=\; F_{\max} \cdot P \cdot (1 + R)
          \;=\; 11 F_{\max}(1 + R).
          \label{eq:derived-invocations}
        \end{align}
        Justification: Theorem~5 / Section~\ref{sec:sketch},
        combining \eqref{eq:derived-transitions} with the per-turn
        fan-out bound $F_{\max}$.

  \item \emph{Minimum distinct agent instances required:}
        \begin{align}
          |\mathcal{A}| \;\geq\; |K| \;=\; 4.
          \label{eq:derived-multiplicity}
        \end{align}
        Justification: Theorem~\ref{thm:multiplicity}, proved in
        Section~\ref{sec:proofs}.

  \item \emph{Trust-tier safety for all subagent content:}
        \begin{align}
          \forall c \in C_{\mathrm{sub}},\quad
          \tau(c) = \mathrm{T\text{-}3}
          \quad\text{unless}\quad
          \exists\, e_j : e_j \in \FN{08}{\text{-}}\Approved(c).
          \label{eq:derived-trust}
        \end{align}
        Justification: Theorem~\ref{thm:trust}, proved in
        Section~\ref{sec:proofs}.

  \item \emph{Acyclicity of the workflow:}
        \begin{align}
          G \text{ is a DAG}.
          \label{eq:derived-acyclic}
        \end{align}
        Justification: Theorem~\ref{thm:acyclic}, proved in
        Section~\ref{sec:proofs}.
\end{enumerate}

\medskip
\noindent\textbf{Internal Consistency of the Tuple.}

The seven components of $\mathcal{R}$ are mutually consistent.
We verify each pairwise dependency explicitly.

\begin{enumerate}
  \item \emph{$G$ and $\mathcal{M}$ are consistent.}
        The phases (vertices of $G$) correspond bijectively to the
        processing states of $Q$ that model active workflow phases
        (states in $Q$ other than $\mathrm{Approved}$,
        $\mathrm{Rejected}$, $\mathrm{Escalated}$ are phase states).
        Each forward arc $(u,v) \in E$ corresponds to an
        $\mathtt{approve}$-transition in $\delta$.
        Rejection arcs (returning to $\mathrm{Drafted}$) exist in
        $\mathcal{M}$ for every non-terminal state; they consume one
        unit of budget $B$.

  \item \emph{$\mathcal{M}$ and $R$ are consistent.}
        The retry limit $R$ controls the budget $B = R \cdot P$.
        The automaton $\mathcal{M}$ is parameterized by $B$:
        the circuit-breaker transition fires from any state $q$ when
        the rejection counter equals $B$.  Setting $R = 0$ gives
        $B = 0$; the circuit-breaker fires on the first rejection.

  \item \emph{$G$, $F_{\max}$, and $D_{\max}$ are consistent.}
        The DAG structure of $G$ ensures phases are visited in a
        strict forward order; $F_{\max}$ and $D_{\max}$ bound the
        parallelism \emph{within} each turn, not across turns.  No
        structural conflict exists between the acyclicity of $G$ and
        the per-turn fan-out.

  \item \emph{$C$ and $\tau$ are consistent.}
        Conflict edges $\{u,v\} \in E_C$ prevent same-agent
        multi-role assignments; $\tau$ governs content-level trust
        tagging.  These two mechanisms operate at different levels
        of abstraction (organizational vs.\ content) and do not
        interfere.

  \item \emph{$C$ and $G$ are consistent.}
        Mandatory workflow phases (Table~3.2, \VisionFlow{})
        induce the mandatory function-class coverage requirement $K$
        (\eqref{eq:K-def}).  The clique structure of $K$ in $C$ forces
        the minimum multiplicity of $4$ agent instances
        (\eqref{eq:derived-multiplicity}).  The $11$-vertex DAG $G$
        and the $12$-vertex conflict graph $C$ are defined over
        different (though related) vertex sets and do not conflict.

  \item \emph{All parameters are finite.}
        $P, R, B, F_{\max}, D_{\max} \in \mathbb{N}_0 \cup \mathbb{N}$;
        $|Q|, |\Sigma|, |V|, |V_C|$ are all finite (verified
        component-wise above).  The run tuple $\mathcal{R}$ is a
        finite mathematical object; no component introduces an
        unbounded or undefined structure.
\end{enumerate}

\medskip
\noindent\textbf{Boundary and Degenerate Cases for the Tuple.}

\begin{enumerate}
  \item \emph{$R = 0$ (no retries).}
        $B = 0$.  Any rejection immediately escalates.
        $\mathcal{R}$ remains well-defined; all bounds hold (with
        $B = 0$, \eqref{eq:derived-transitions} gives
        $\text{transitions} \leq P$).

  \item \emph{$F_{\max} = 1$ (serial execution).}
        Each Orchestrator turn spawns at most one subagent.
        \eqref{eq:derived-invocations} gives
        $N \leq P(1+R)$, matching the transition count.
        $\mathcal{R}$ is well-defined; all proofs apply without
        modification.

  \item \emph{$P = 1$ (hypothetical single-phase workflow).}
        $B = R$.
        \eqref{eq:derived-transitions} gives
        $\text{transitions} \leq 1 + R$.
        Not the actual configuration (actual $P = 11$), but the
        tuple and all proofs remain valid for any $P \geq 1$.

  \item \emph{Maximum saturation.}
        If $R$, $F_{\max}$ are both at their maximum permitted
        values and every turn saturates fan-out, the run reaches
        $N = F_{\max} \cdot P \cdot (1+R)$ invocations before
        termination.  All bounds hold with equality; no
        component of $\mathcal{R}$ is violated.

  \item \emph{Immediate approval (no rework).}
        If every phase gate approves on first attempt, $b_i = 0$
        throughout, and the run terminates in exactly $P = 11$
        state-machine transitions.  This is the best-case run
        length; the upper bound $P + B \geq P$ is satisfied.
\end{enumerate}

% ============================================================================
\section{Mathematical Proofs}
\label{sec:proofs}

% WHAT: Formal proofs of four workflow properties plus one proof sketch.
% WHY:  Per mathematical_proof_structures_standards.txt P.1.3, every
%       algorithmic/architectural claim in this codebase (termination,
%       non-overlap, safety invariant) must be backed by a proof or a
%       peer-reviewed citation; this section discharges that obligation
%       for the four claims stated in the Executive Summary.
% HOW:  Each theorem below declares its proof type, states hypotheses
%       explicitly before the proof body, justifies every non-trivial
%       step, and closes with the amsthm-automatic QED marker.

\subsection{Theorem 1 --- Acyclicity of the Nominal Workflow Graph}

\textbf{Proof type:} Direct proof (via exhibited topological numbering).

\textbf{Hypotheses:}
\begin{enumerate}
    \item[H1] $G = (V,E)$ is the workflow graph of Definition~\ref{def:workflow-graph}.
    \item[H2] Every edge $(u,v) \in E$ corresponds to a "must-complete-before"
              constraint from Table~3.2 or a join-point constraint from
              \S7.2 of \VisionFlow{}.
\end{enumerate}

\begin{theorem}[Acyclicity]
\label{thm:acyclic}
Under H1--H2, $G$ is a directed acyclic graph (DAG).
\end{theorem}

\begin{proof}
The proof proceeds in five stages: (I) construction and validation of the
labeling function $\phi$, (II) exhaustive verification of the strict
monotonicity condition on every edge, (III) proof of the cycle-exclusion
lemma from first principles, (IV) application to $G$, and (V) confirmation
that no boundary or degenerate case is left unaddressed.

\medskip
\noindent\textbf{Stage I: Construction and Well-Definedness of $\phi$.}

\begin{enumerate}
  \item \emph{Domain.} By H1, $V$ is the finite vertex set of $G$.
        By Definition~\ref{def:workflow-graph} and Table~3.2, the
        enumeration of $V$ is exhaustive and fixed:
        \begin{align}
          V \;=\; \{&\mathrm{Plan},\; \mathrm{EnvVerify},\;
                     \mathrm{StdMap},\; \mathrm{Impl},\notag\\
                    &\mathrm{Review},\; \mathrm{Challenge},\;
                     \mathrm{Validate},\; \mathrm{Test},\notag\\
                    &\mathrm{Doc},\; \mathrm{CommitSync},\;
                     \mathrm{Gate}\}.
        \end{align}
        Hence $|V| = 11$.

  \item \emph{Codomain.} Define $\phi : V \to \mathbb{Q}$ (rational
        numbers) rather than $\{0,\ldots,7\} \subset \mathbb{Z}$, because
        two vertices (\textit{Review}/\textit{Challenge} and
        \textit{Validate}) occupy fractional positions within Phase~3.
        The rationals are totally ordered and contain $\mathbb{Z}$ as a
        subring, so all subsequent strict-order arguments apply without
        modification.

  \item \emph{Point-wise definition.} Assign $\phi$ by the STD-036 phase
        index (Table~3.2) and the intra-phase ordering of
        \S4.2.2,
        \VisionFlow{}:
        \begin{align}
          \phi(\mathrm{Plan})        &= 0,   \label{eq:phi-plan}\\
          \phi(\mathrm{EnvVerify})   &= 1,   \label{eq:phi-env}\\
          \phi(\mathrm{StdMap})      &= 2,   \label{eq:phi-std}\\
          \phi(\mathrm{Impl})        &= 3,   \label{eq:phi-impl}\\
          \phi(\mathrm{Review})      &= \tfrac{7}{2}, \label{eq:phi-rev}\\
          \phi(\mathrm{Challenge})   &= \tfrac{7}{2}, \label{eq:phi-chal}\\
          \phi(\mathrm{Validate})    &= \tfrac{15}{4},\label{eq:phi-val}\\
          \phi(\mathrm{Test})        &= 4,   \label{eq:phi-test}\\
          \phi(\mathrm{Doc})         &= 5,   \label{eq:phi-doc}\\
          \phi(\mathrm{CommitSync})  &= 6,   \label{eq:phi-cs}\\
          \phi(\mathrm{Gate})        &= 7.   \label{eq:phi-gate}
        \end{align}

  \item \emph{Total definition.} Every element of $V$ appears exactly once
        in \eqref{eq:phi-plan}--\eqref{eq:phi-gate}; therefore $\phi$ is
        a total function on $V$.

  \item \emph{Well-definedness of equal values.}
        $\phi(\mathrm{Review}) = \phi(\mathrm{Challenge}) = \tfrac{7}{2}$
        is admissible: a labeling function need not be injective; it must
        only be strictly increasing \emph{along edges}.  The equality of
        labels for parallel siblings is consistent with the fact that no
        edge exists between \textit{Review} and \textit{Challenge} (they
        are concurrent by \S4.2.2; see Stage~II, row~5 below).
\end{enumerate}

\medskip
\noindent\textbf{Stage II: Exhaustive Edge-Wise Verification of
$\phi(u) < \phi(v)$ for every $(u,v)\in E$.}

\begin{sloppypar}
By H2, the edge set $E$ is exactly the set of ordering constraints listed in
Table~3.2 and the join-point constraints of \S7.2 of
\VisionFlow{}.
The following table enumerates \emph{every} such edge and confirms the
strict inequality.  (No edge is omitted; the table is exhaustive by
construction from the finite source document.)
\end{sloppypar}

\begin{enumerate}
  \item $(\mathrm{Plan},\mathrm{EnvVerify})$:
        $\phi(\mathrm{Plan}) = 0 < 1 = \phi(\mathrm{EnvVerify})$. \checkmark

  \item $(\mathrm{EnvVerify},\mathrm{StdMap})$:
        $\phi(\mathrm{EnvVerify}) = 1 < 2 = \phi(\mathrm{StdMap})$.
        \checkmark

  \item $(\mathrm{StdMap},\mathrm{Impl})$:
        $\phi(\mathrm{StdMap}) = 2 < 3 = \phi(\mathrm{Impl})$. \checkmark

  \item $(\mathrm{Impl},\mathrm{Review})$:
        $\phi(\mathrm{Impl}) = 3 < \tfrac{7}{2} = \phi(\mathrm{Review})$.
        \checkmark

  \item $(\mathrm{Impl},\mathrm{Challenge})$:
        $\phi(\mathrm{Impl}) = 3 < \tfrac{7}{2} = \phi(\mathrm{Challenge})$.
        \checkmark
        \newline
        \emph{Note:} No edge $(\mathrm{Review},\mathrm{Challenge})$ or
        $(\mathrm{Challenge},\mathrm{Review})$ exists (parallel siblings);
        hence the equal-label case never arises on any edge.

  \item $(\mathrm{Review},\mathrm{Validate})$ and
        $(\mathrm{Challenge},\mathrm{Validate})$ (join point, \S7.2):
        \begin{align}
          \phi(\mathrm{Review})    &= \tfrac{7}{2} < \tfrac{15}{4}
            = \phi(\mathrm{Validate}), \label{eq:ev-rev-val}\\
          \phi(\mathrm{Challenge}) &= \tfrac{7}{2} < \tfrac{15}{4}
            = \phi(\mathrm{Validate}). \label{eq:ev-chal-val}
        \end{align}
        Both hold because
        $\tfrac{7}{2} = \tfrac{14}{4} < \tfrac{15}{4}$. \checkmark

  \item $(\mathrm{Validate},\mathrm{Test})$:
        $\phi(\mathrm{Validate}) = \tfrac{15}{4} < 4 = \phi(\mathrm{Test})$
        because $\tfrac{15}{4} < \tfrac{16}{4} = 4$. \checkmark

  \item $(\mathrm{Test},\mathrm{Doc})$:
        $\phi(\mathrm{Test}) = 4 < 5 = \phi(\mathrm{Doc})$. \checkmark

  \item $(\mathrm{Doc},\mathrm{CommitSync})$:
        $\phi(\mathrm{Doc}) = 5 < 6 = \phi(\mathrm{CommitSync})$.
        \checkmark

  \item $(\mathrm{CommitSync},\mathrm{Gate})$:
        $\phi(\mathrm{CommitSync}) = 6 < 7 = \phi(\mathrm{Gate})$.
        \checkmark
\end{enumerate}

\noindent
All $|E|$ edges have been enumerated and verified.  No edge was skipped.
The edge set is finite (bounded by the finite vertex set $|V|=11$); hence
the exhaustion terminates.

\medskip
\noindent\textbf{Stage III: Cycle-Exclusion Lemma (from first principles).}

\begin{enumerate}
  \item \emph{Assumption for contradiction.}  Suppose $G$ contains a
        directed cycle
        \begin{equation}
          C = v_1 \to v_2 \to \cdots \to v_k \to v_1,
          \quad k \geq 1,\; v_i \in V. \label{eq:cycle-assume}
        \end{equation}

  \item \emph{Degenerate case $k=1$.}  A self-loop $(v_1,v_1)\in E$ would
        require $\phi(v_1) < \phi(v_1)$ by Stage~II.  But $\phi(v_1) <
        \phi(v_1)$ is false by the \emph{irreflexivity} of $<$ on
        $\mathbb{Q}$ (axiom of strict total orders).  Contradiction.

  \item \emph{Case $k \geq 2$.}  Each consecutive pair $(v_i, v_{i+1})$
        for $i=1,\ldots,k-1$ and the closing pair $(v_k, v_1)$ are edges
        in $E$.  By Stage~II, for every edge $(u,v)\in E$ we have
        $\phi(u) < \phi(v)$.  Applying this to each edge of $C$:
        \begin{align}
          \phi(v_1) &< \phi(v_2),   \label{eq:chain1}\\
          \phi(v_2) &< \phi(v_3),   \label{eq:chain2}\\
                    &\;\;\vdots     \notag \\
          \phi(v_{k-1}) &< \phi(v_k), \label{eq:chaink}\\
          \phi(v_k) &< \phi(v_1).   \label{eq:chainclose}
        \end{align}

  \item \emph{Transitivity.}  The relation $<$ on $\mathbb{Q}$ is
        transitive.  Applying transitivity to
        \eqref{eq:chain1}--\eqref{eq:chaink} by induction on $k$:
        \begin{equation}
          \phi(v_1) < \phi(v_k). \label{eq:trans-result}
        \end{equation}

  \item \emph{Contradiction.}  Combining \eqref{eq:trans-result} with
        \eqref{eq:chainclose} by transitivity:
        \begin{equation}
          \phi(v_1) < \phi(v_1). \label{eq:contradiction}
        \end{equation}
        This contradicts the irreflexivity of $<$ on $\mathbb{Q}$.

  \item \emph{Boundary case $k=2$.}  A 2-cycle $(v_1,v_2,v_1)$ with
        $v_1 \neq v_2$ would require $\phi(v_1) < \phi(v_2)$ (from
        edge $(v_1,v_2)$) \emph{and} $\phi(v_2) < \phi(v_1)$ (from edge
        $(v_2,v_1)$), hence $\phi(v_1) < \phi(v_1)$ by transitivity:
        contradiction identical to step~5 above.

  \item \emph{Conclusion of lemma.}  In every case ($k=1$, $k=2$, $k\geq
        3$) the assumption \eqref{eq:cycle-assume} leads to a
        contradiction.  Therefore no directed cycle exists in $G$.
\end{enumerate}

\medskip
\noindent\textbf{Stage IV: Application --- $G$ is a DAG.}

\begin{enumerate}
  \item By Stage~I, $\phi$ is a total, well-defined function $V \to
        \mathbb{Q}$.
  \item By Stage~II, $\phi(u) < \phi(v)$ for every $(u,v)\in E$.
  \item By Stage~III, no directed cycle exists in $G$.
  \item A finite directed graph with no directed cycle satisfies the
        definition of a directed acyclic graph (DAG) (Cormen, Leiserson,
        Rivest, and Stein, \emph{Introduction to Algorithms},
        3rd~ed., MIT Press, 2009, Theorem~22.12).
  \item Therefore $G$ is a DAG. \qedhere
\end{enumerate}

\medskip
\noindent\textbf{Stage V: Boundary and Degenerate Cases (checklist).}

\begin{enumerate}
  \item \emph{Empty graph ($|V|=0$ or $|E|=0$).}  If $E = \emptyset$
        there are no edges and trivially no cycle; the DAG conclusion
        holds.  The actual $G$ has $|V|=11$ and $|E|>0$ by
        Definition~\ref{def:workflow-graph}, so this case is vacuously
        satisfied.

  \item \emph{Disconnected components.}  $\phi$ is defined on all
        vertices regardless of connectivity; Stage~III applies to any
        connected component independently.

  \item \emph{Multiple edges between the same pair.}  H2 states edges
        correspond to distinct documented ordering constraints; the
        source tables contain no duplicate directed edge.  Even if
        duplicates existed, Stage~III would still apply, since each copy
        satisfies $\phi(u) < \phi(v)$.

  \item \emph{Equal-label vertices on an edge.}  Stage~II item~5
        confirms that the only equal-label vertices are
        \textit{Review} and \textit{Challenge}, and no edge exists
        between them.  No edge $(u,v)\in E$ has $\phi(u) = \phi(v)$.

  \item \emph{Well-foundedness of the rational order.}  The strict
        order $<$ on $\mathbb{Q}$ is irreflexive, asymmetric, and
        transitive.  All three properties were invoked explicitly above;
        no additional axiomatic assumptions are required.
\end{enumerate}
\end{proof}

\begin{sloppypar}
\textbf{Dependencies:} Table~3.2 and \S7.1--7.2,
\VisionFlow{}; standard
graph-theoretic fact that a graph with a strictly monotone vertex labeling
along every edge is acyclic (equivalent to the topological-ordering
characterization of DAGs, e.g. Cormen, Leiserson, Rivest, and Stein,
\emph{Introduction to Algorithms}, 3rd ed., MIT Press, 2009, Theorem~22.12).
\end{sloppypar}

\subsection{Theorem 2 --- Bounded Rework Termination}

\textbf{Proof type:} Direct proof by a strictly decreasing potential
function (well-founded induction on $\mathbb{N}$).

\textbf{Hypotheses:}
\begin{sloppypar}
\begin{enumerate}
    \item[H1] The retry limit for any single phase-gate is $R \in
              \mathbb{N}$, $R \geq 0$ (\S8.2.1,
              \VisionGov{}).
    \item[H2] $G$ has a finite vertex set $|V| = P$ (Definition
              \ref{def:workflow-graph}; here $P = 11$).
    \item[H3] Every rejection event (Table~8.1,
              \VisionFlow{})
              consumes exactly one unit of a global rework budget
              $B \coloneqq R \cdot P$ shared across all phases, and routes
              the artifact back to state $\mathrm{Drafted}$
              (Definition~\ref{def:state-machine}).
\end{enumerate}
\end{sloppypar}

\begin{theorem}[Bounded Termination]
\label{thm:rework}
Under H1--H3, any run of the artifact state machine $\mathcal{M}$ reaches
either $\mathrm{Approved}$ or an explicit circuit-breaker escalation
(\S5.3.2, \VisionGov{})
within at most $B + P$ state transitions.
\end{theorem}

\begin{proof}
The proof proceeds in six stages: (I)~variable and notation setup,
(II)~validation of all hypotheses, (III)~potential-function definition
and domain verification, (IV)~base case, (V)~inductive step with full
case analysis, and (VI)~boundary and degenerate cases, followed by the
final conclusion.

\medskip
\noindent\textbf{Stage I: Variable and Notation Setup.}

\begin{enumerate}
  \item Let $\mathcal{M}$ denote the artifact state machine of
        Definition~\ref{def:state-machine}.  A \emph{run} of
        $\mathcal{M}$ is a (possibly infinite) sequence of states
        $s_0, s_1, s_2, \ldots$ where $s_0 = \mathrm{Drafted}$ and
        each consecutive pair $(s_i, s_{i+1})$ is a legal transition
        of $\mathcal{M}$.

  \item Let $i \in \mathbb{N}_0 \coloneqq \{0,1,2,\ldots\}$ index
        total state-machine steps elapsed since the start of the run.

  \item Define $k_i \in \mathbb{N}_0$ as the number of
        \emph{forward} (non-rework) phase transitions completed after
        exactly $i$ steps.  A forward transition is any transition
        that advances the artifact from one phase gate to the next in
        the workflow graph $G$ (Definition~\ref{def:workflow-graph})
        without returning to $\mathrm{Drafted}$.

  \item Define $b_i \in \mathbb{N}_0$ as the cumulative
        rework-budget units consumed after exactly $i$ steps.  Each
        rejection event (H3) increments $b_i$ by exactly one.

  \item The global rework budget is $B \coloneqq R \cdot P$ (H1,
        H2, H3), where $R \geq 0$ is the per-phase retry limit and
        $P = |V|$ is the number of phases.

  \item The notation $[n] \coloneqq \{0, 1, \ldots, n\}$ for any
        $n \in \mathbb{N}_0$ is used throughout.
\end{enumerate}

\medskip
\noindent\textbf{Stage II: Hypothesis Validation.}

\begin{enumerate}
  \item \emph{H1 — Retry limit.}  $R \in \mathbb{N}$, $R \geq 0$ is
        asserted by \S8.2.1 of
        \VisionGov{}.  This ensures $R$ is a well-defined
        non-negative integer, so $B = R \cdot P$ is also a
        non-negative integer.

  \item \emph{H2 — Finite phase set.}  $|V| = P$ is given by
        Definition~\ref{def:workflow-graph} with the concrete value
        $P = 11$ confirmed by Table~3.2 of
        \VisionFlow{}.
        In particular $P \in \mathbb{N}$, $P \geq 1$.

  \item \emph{H3 — Rejection semantics.}  Every rejection event
        (i)~consumes exactly one unit of the global budget $B$,
        (ii)~routes the artifact back to $\mathrm{Drafted}$, and
        (iii)~when $b_i = B$ triggers the circuit-breaker escalation
        defined in \S5.3.2 of
        \VisionGov{}.  These three sub-conditions are each
        independently stated in Table~8.1 and Figure~5.1 of
        \VisionFlow{}.

  \item \emph{Derived bound on $k_i$.}  Because $G$ has exactly $P$
        vertices (H2) and each forward transition moves to a
        strictly later phase (Theorem~\ref{thm:acyclic}, which establishes
        that $G$ is a DAG), the number of forward transitions is
        bounded: $k_i \leq P$ for all $i$.

  \item \emph{Derived bound on $b_i$.}  By H3, once $b_i = B$ no
        further rejection is processed (the escalation fires
        instead).  Hence $b_i \leq B$ for all $i$ in any run that
        has not yet terminated.
\end{enumerate}

\medskip
\noindent\textbf{Stage III: Potential Function — Definition and Domain.}

Define
\begin{equation}
\label{eq:potential}
\Phi(i) \;\coloneqq\; \bigl(P - k_i\bigr) + \bigl(B - b_i\bigr).
\end{equation}

\begin{enumerate}
  \item \emph{Well-definedness.}  Both $k_i$ and $b_i$ are
        non-negative integers for all $i \geq 0$ (Stage~I).
        Subtraction in \eqref{eq:potential} is taken in $\mathbb{Z}$;
        the result lands in $\mathbb{N}_0$ by the bounds of Stage~II
        items~4--5.

  \item \emph{Integer-valued.}  $P, k_i, B, b_i \in \mathbb{N}_0$
        (H1, H2, Stage~I), so $\Phi(i) \in \mathbb{Z}$ for all $i$.

  \item \emph{Lower bound.}  From Stage~II items~4--5:
        \begin{align}
          0 \leq k_i \leq P &\implies P - k_i \geq 0, \label{eq:lb-k}\\
          0 \leq b_i \leq B &\implies B - b_i \geq 0. \label{eq:lb-b}
        \end{align}
        Adding \eqref{eq:lb-k} and \eqref{eq:lb-b}:
        \begin{equation}
          \label{eq:phi-lb}
          \Phi(i) \;\geq\; 0 \quad \text{for all } i \geq 0.
        \end{equation}

  \item \emph{Upper bound.}  At $i = 0$: $k_0 = 0$ and $b_0 = 0$
        (no transitions have occurred), so
        \begin{equation}
          \label{eq:phi-ub}
          \Phi(0) \;=\; (P - 0) + (B - 0) \;=\; P + B.
        \end{equation}
        Since $P \geq 1$ (H2) and $B \geq 0$ (H1), we have
        $\Phi(0) \in \mathbb{N}$, a finite positive integer.
\end{enumerate}

\medskip
\noindent\textbf{Stage IV: Base Case ($i = 0$).}

\begin{enumerate}
  \item At the start of any run, no state-machine step has been
        executed.  By definition (Stage~I, items~3--4):
        \begin{equation}
          \label{eq:base-init}
          k_0 = 0, \quad b_0 = 0.
        \end{equation}

  \item Substituting \eqref{eq:base-init} into \eqref{eq:potential}:
        \begin{align}
          \Phi(0)
            &= (P - k_0) + (B - b_0) \notag\\
            &= (P - 0)   + (B - 0)   \notag\\
            &= P + B.     \label{eq:base-val}
        \end{align}

  \item By H1 and H2: $P \geq 1$ and $B = R \cdot P \geq 0$, so
        $\Phi(0) = P + B \geq 1 > 0$.  In particular
        $\Phi(0) \in \mathbb{N}$ (finite and positive).

  \item The lower bound \eqref{eq:phi-lb} holds at $i = 0$:
        $\Phi(0) = P + B \geq 0$. \checkmark
\end{enumerate}

\medskip
\noindent\textbf{Stage V: Inductive Step.}

\noindent\textit{Induction hypothesis (IH):} Assume that after $i$
steps ($i \geq 0$) the run has not yet terminated and
\begin{equation}
\label{eq:IH}
0 \;\leq\; \Phi(i) \;\leq\; P + B, \quad
k_i \in [P], \quad b_i \in [B].
\end{equation}

\noindent We show $\Phi(i+1) = \Phi(i) - 1$, which simultaneously
establishes (a) strict decrease and (b) that the bounds in
\eqref{eq:IH} propagate to step $i+1$.

\medskip
\noindent\textit{Exhaustive case analysis.}
By Figure~5.1 of \VisionFlow{}, at every non-terminal step exactly one of the
following two mutually exclusive and jointly exhaustive events occurs:

\begin{enumerate}
  \item \textbf{Case~F (Forward Transition).}
        The current phase gate is approved and the artifact advances
        to the next phase in $G$.
        \begin{enumerate}
          \item By definition (Stage~I, item~3), a forward transition
                increments $k_i$ by exactly one:
                \begin{equation}
                  \label{eq:F-k}
                  k_{i+1} = k_i + 1.
                \end{equation}
          \item No rework-budget unit is consumed (H3 applies only to
                rejection events):
                \begin{equation}
                  \label{eq:F-b}
                  b_{i+1} = b_i.
                \end{equation}
          \item The new forward-step count is in range: from
                \eqref{eq:F-k} and the IH bound $k_i \leq P$,
                \begin{equation}
                  \label{eq:F-k-bound}
                  k_{i+1} = k_i + 1 \leq P + 1.
                \end{equation}
                However, if $k_i = P$ then all $P$ phases are already
                complete and the run has terminated at
                $\mathrm{Approved}$; we are in the non-terminal case by
                the IH assumption, so $k_i \leq P - 1$, giving
                $k_{i+1} \leq P$.
          \item The budget bound is inherited: $b_{i+1} = b_i \leq B$
                (IH).
          \item Computing $\Phi(i+1)$:
                \begin{align}
                  \Phi(i+1)
                    &= \bigl(P - k_{i+1}\bigr) + \bigl(B - b_{i+1}\bigr)
                       && \text{(def.\ \eqref{eq:potential})} \notag\\
                    &= \bigl(P - (k_i+1)\bigr) + \bigl(B - b_i\bigr)
                       && \text{(\eqref{eq:F-k}, \eqref{eq:F-b})} \notag\\
                    &= \bigl(P - k_i - 1\bigr) + \bigl(B - b_i\bigr)
                       \notag\\
                    &= \bigl[(P - k_i) + (B - b_i)\bigr] - 1
                       \notag\\
                    &= \Phi(i) - 1.
                       \label{eq:F-result}
                \end{align}
        \end{enumerate}

  \item \textbf{Case~R (Rejection / Rework).}
        The current phase gate rejects the artifact and routes it back
        to $\mathrm{Drafted}$ (H3).
        \begin{enumerate}
          \item By H3, exactly one budget unit is consumed:
                \begin{equation}
                  \label{eq:R-b}
                  b_{i+1} = b_i + 1.
                \end{equation}
          \item The forward-step counter is \emph{not} decremented by
                the state machine; rather, the artifact returns to
                $\mathrm{Drafted}$ and must re-traverse phases.
                Formally, $k_{i+1}$ records only the count of
                successful forward advances (Stage~I, item~3); it does
                not decrease upon rejection:
                \begin{equation}
                  \label{eq:R-k}
                  k_{i+1} = k_i.
                \end{equation}
                Consequently, from the IH bound $k_i \leq P-1$ (run
                is non-terminal), we have $k_{i+1} \leq P - 1 \leq P$.
          \item The new budget value is in range: from \eqref{eq:R-b}
                and the IH bound $b_i \leq B$,
                \begin{equation}
                  \label{eq:R-b-bound}
                  b_{i+1} = b_i + 1 \leq B + 1.
                \end{equation}
                However, if $b_i = B$ then the escalation fires at
                step $i$ (H3, Stage~II item~3), so the run terminates
                before step $i+1$; the IH assumption excludes this.
                Therefore $b_i \leq B - 1$, giving $b_{i+1} \leq B$.
          \item Computing $\Phi(i+1)$:
                \begin{align}
                  \Phi(i+1)
                    &= \bigl(P - k_{i+1}\bigr) + \bigl(B - b_{i+1}\bigr)
                       && \text{(def.\ \eqref{eq:potential})} \notag\\
                    &= \bigl(P - k_i\bigr) + \bigl(B - (b_i+1)\bigr)
                       && \text{(\eqref{eq:R-k}, \eqref{eq:R-b})} \notag\\
                    &= \bigl(P - k_i\bigr) + \bigl(B - b_i - 1\bigr)
                       \notag\\
                    &= \bigl[(P - k_i) + (B - b_i)\bigr] - 1
                       \notag\\
                    &= \Phi(i) - 1.
                       \label{eq:R-result}
                \end{align}
        \end{enumerate}
\end{enumerate}

\noindent\textit{Synthesis.}  In both Case~F and Case~R:
\begin{equation}
  \label{eq:step-decrease}
  \Phi(i+1) = \Phi(i) - 1.
\end{equation}
From \eqref{eq:IH} and \eqref{eq:step-decrease}:
\begin{align}
  \Phi(i+1) &= \Phi(i) - 1
               \;\leq\; (P+B) - 1
               \;<\; P+B,
               \label{eq:ub-prop}\\
  \Phi(i+1) &= \Phi(i) - 1
               \;\geq\; 0 - 1 = -1. \label{eq:lb-prop-raw}
\end{align}
The lower bound requires care: \eqref{eq:lb-prop-raw} gives $-1$, but
\eqref{eq:phi-lb} guarantees $\Phi(i+1) \geq 0$ whenever the run has not
yet terminated.  The combination of the two facts is:
\begin{equation}
  \label{eq:lb-prop}
  \text{if the run is non-terminal after step } i+1,
  \text{ then } \Phi(i+1) \geq 0.
\end{equation}
Proof of \eqref{eq:lb-prop}: non-terminal means $k_{i+1} \leq P$ and
$b_{i+1} \leq B$ (Stage~II items~4--5), so both summands in
\eqref{eq:potential} are non-negative, giving $\Phi(i+1) \geq 0$.
\eqref{eq:ub-prop} and \eqref{eq:lb-prop} confirm that the IH
\eqref{eq:IH} holds at $i+1$, completing the inductive step. \checkmark

\medskip
\noindent\textbf{Stage VI: Boundary and Degenerate Cases.}

\begin{enumerate}
  \item \emph{$R = 0$ (zero retries).}  Then $B = 0 \cdot P = 0$.
        Any rejection event immediately exhausts the budget ($b_1 = 1
        > B = 0$), triggering the circuit-breaker at step~1.  The
        bound $B + P = P \geq 1$ still holds.  The potential satisfies
        $\Phi(0) = P + 0 = P$, and the proof is unchanged.

  \item \emph{$P = 1$ (single phase).}  Then $B = R$.  Exactly one
        forward transition suffices for $\mathrm{Approved}$; at most
        $R$ rejections are permitted before escalation.  The bound
        $B + P = R + 1$.  All stages apply without modification.

  \item \emph{Maximum run length achieved.}
        If the run uses all $B$ rejections and all $P$ forward
        transitions, then $\Phi(B+P) = (P-P)+(B-B) = 0$, confirming
        the bound is tight.

  \item \emph{Termination before $B + P$ steps.}
        The run may terminate earlier (e.g.\ if no rework occurs,
        then $P$ forward transitions suffice and the run ends at
        step~$P \leq B+P$).  The theorem states an upper bound, not
        an exact count.

  \item \emph{Mutual exclusivity of Cases~F and~R.}
        A single state-machine step is an atomic transition; it is
        either a forward approval or a rejection, never both
        simultaneously.  This is enforced by the deterministic
        transition relation of $\mathcal{M}$
        (Definition~\ref{def:state-machine}).

  \item \emph{Exhaustiveness of Cases~F and~R.}
        \begin{sloppypar}
        Figure~5.1 of
        \VisionFlow{}
        enumerates all legal non-terminal transitions; every such
        transition is either a forward advancement or a
        rejection-and-reset.  No third category exists.
        \end{sloppypar}

  \item \emph{Well-foundedness.}
        $\Phi$ maps each non-terminal step to a value in
        $\{0, 1, \ldots, P+B\} \subset \mathbb{N}_0$.  By
        \eqref{eq:step-decrease}, $\Phi$ is strictly decreasing along
        any run.  A strictly decreasing sequence of non-negative
        integers is finite by the well-ordering principle on
        $\mathbb{N}$ (every non-empty subset of $\mathbb{N}_0$ has a
        least element; equivalently, there is no infinite strictly
        descending chain in $\mathbb{N}_0$).

  \item \emph{Uniqueness of termination conditions.}
        The only two ways a non-terminal run at step $i$ can produce
        no legal next transition are:
        \begin{enumerate}
          \item $k_i = P$: all $P$ phases are approved; the artifact
                is in state $\mathrm{Approved}$ (Definition
                \ref{def:state-machine}).
          \item $b_i = B$: the global rework budget is exhausted;
                by H3 and \S5.3.2 of the governance document this
                \emph{is} the circuit-breaker escalation condition.
        \end{enumerate}
        No other termination condition exists by the exhaustiveness
        of items~(a)--(b) above, confirmed by Table~8.1 of
        \VisionFlow{}.
\end{enumerate}

\medskip
\noindent\textbf{Conclusion.}
By Stage~IV, $\Phi(0) = P + B$.  By Stage~V, every non-terminal step
strictly decreases $\Phi$ by exactly one.  By Stage~VI item~7
(well-ordering), the run must terminate within at most $\Phi(0) = P + B$
steps.  By Stage~VI item~8, the terminal state is either
$\mathrm{Approved}$ or a circuit-breaker escalation.  Therefore:
\begin{equation}
  \label{eq:conclusion}
  \text{every run of } \mathcal{M}
  \text{ terminates in } \mathrm{Approved}
  \text{ or escalation within at most } P + B
  \text{ steps.} \qedhere
\end{equation}
\end{proof}

\begin{sloppypar}
\textbf{Dependencies:} \S5.3, \S8.2,
\VisionGov{}; Table~8.1,
Figure~5.1, \VisionFlow{};
well-ordering principle on $\mathbb{N}$ (standard, e.g. Rosen,
\emph{Discrete Mathematics and Its Applications}, 8th ed., McGraw-Hill,
2019, \S5.2).
\end{sloppypar}

\subsection{Theorem 3 --- Minimum Agent-Instance Multiplicity}

\textbf{Proof type:} Direct proof by finite combinatorics (clique
argument).

\textbf{Hypotheses:}
\begin{enumerate}
    \item[H1] $C = (\{\FN{00},\ldots,\FN{11}\}, E_C)$ is the conflict
              graph of Definition~\ref{def:conflict-graph}.
    \item[H2] Every agent instance $a$ (Definition~\ref{def:agent}) is
              \emph{compliant} iff $S(a)$ is an independent set of $C$,
              i.e. no two functions in $S(a)$ are adjacent in $C$ (this is
              the direct restatement of the Non-Overlap Rule, \S5.1,
              \VisionFunc{}).
    \item[H3] For a single artifact's lifecycle to be fully governed, every
              vertex of the clique
              $K = \{\FN{02},\FN{04},\FN{05},\FN{09}\} \subseteq
              \{\FN{00},\ldots,\FN{11}\}$ must be performed by
              \emph{some} compliant agent instance (Specialize, Review,
              Validate, and Challenge are all mandatory phases per
              Table~3.2, \VisionFlow{}), and every pair of distinct vertices
              of $K$ is an edge of $C$ (by Definition~\ref{def:conflict-graph},
              since \FN{02} conflicts with \FN{04} and \FN{05}, and
              \FN{04} conflicts with \FN{09}; the remaining pairs
              $\{\FN{02},\FN{09}\}$ and $\{\FN{05},\FN{09}\}$ are likewise
              forbidden, since a Specialist or Validator certifying its own
              adversarial challenge would itself violate the "no
              self-certification" clause of \S1.2.3,
              \VisionClean{}).
\end{enumerate}

\begin{theorem}[Minimum Multiplicity]
\label{thm:multiplicity}
Under H1--H3, any compliant deployment covering the four mandatory functions
$K$ for one artifact requires at least $|K| = 4$ distinct agent instances.
\end{theorem}

\begin{proof}
The proof proceeds in seven stages, each fully justified, with no step
omitted and all boundary conditions addressed.

\medskip
\noindent\textbf{Stage~I: Validate the domain of $K$.}
\begin{enumerate}
    \item[I.1] By H1, the vertex set of $C$ is
               $V_C = \{\FN{00},\FN{01},\ldots,\FN{11}\}$, which has
               cardinality $|V_C| = 12$.
    \item[I.2] By H3, $K = \{\FN{02},\FN{04},\FN{05},\FN{09}\}$.
               Each element of $K$ is distinct (the four function
               identifiers are pairwise unequal by inspection), so
               $|K| = 4$.
    \item[I.3] $K \subseteq V_C$ because $\FN{02}, \FN{04}, \FN{05},
               \FN{09} \in \{\FN{00},\ldots,\FN{11}\}$.
               Hence $K$ is a well-defined vertex subset of $C$.
\end{enumerate}

\medskip
\noindent\textbf{Stage~II: Enumerate every pair in $K$ and verify each
is an edge.}

There are $\binom{4}{2} = 6$ unordered pairs in $K$.
We list them all and confirm membership in $E_C$ via H3:
\begin{enumerate}
    \item[II.1] $\{\FN{02},\FN{04}\}$: Specialize $\leftrightarrow$
                Review conflict; $\{\FN{02},\FN{04}\} \in E_C$
                by H3 (explicit citation therein).
    \item[II.2] $\{\FN{02},\FN{05}\}$: Specialize $\leftrightarrow$
                Validate conflict; $\{\FN{02},\FN{05}\} \in E_C$
                by H3 (explicit citation therein).
    \item[II.3] $\{\FN{04},\FN{09}\}$: Review $\leftrightarrow$
                Challenge conflict; $\{\FN{04},\FN{09}\} \in E_C$
                by H3 (explicit citation therein).
    \item[II.4] $\{\FN{02},\FN{09}\}$: a Specialist certifying its
                own adversarial challenge violates the
                \begin{sloppypar}
                ``no self-certification'' clause (\S1.2.3,
                \VisionClean{});
                $\{\FN{02},\FN{09}\} \in E_C$ by H3.
                \end{sloppypar}
    \item[II.5] $\{\FN{05},\FN{09}\}$: a Validator certifying its
                own adversarial challenge similarly violates \S1.2.3
                of the same document;
                $\{\FN{05},\FN{09}\} \in E_C$ by H3.
    \item[II.6] $\{\FN{04},\FN{05}\}$: Review $\leftrightarrow$
                Validate; both are certification roles over the same
                artifact output, and joint exercise by one agent
                violates the independence requirement of \S5.1,
                \VisionFunc{};
                $\{\FN{04},\FN{05}\} \in E_C$ by H3.
\end{enumerate}

\noindent All six pairs are confirmed edges.  Since every pair of
distinct vertices in $K$ is an edge of $C$, $K$ is a \emph{clique} of
size $4$ in $C$:
\begin{align}
    \label{eq:clique}
    \forall\, u, v \in K,\; u \neq v
    \;\Longrightarrow\;
    \{u,v\} \in E_C.
\end{align}

\medskip
\noindent\textbf{Stage~III: State the coverage requirement precisely.}
\begin{enumerate}
    \item[III.1] Let $\mathcal{A}$ be any finite set of compliant agent
                 instances. Define the \emph{coverage} of $\mathcal{A}$
                 as
                 \begin{align}
                     \label{eq:coverage}
                     \mathrm{Cov}(\mathcal{A})
                     \;=\;
                     \bigcup_{a \in \mathcal{A}} S(a).
                 \end{align}
    \item[III.2] We say $\mathcal{A}$ \emph{covers} $K$ iff
                 $K \subseteq \mathrm{Cov}(\mathcal{A})$, i.e.\ every
                 mandatory function has at least one agent instance
                 responsible for it.
    \item[III.3] The goal is to establish
                 $|\mathcal{A}| \geq 4$
                 for every $\mathcal{A}$ that covers $K$ and whose
                 every member is compliant.
\end{enumerate}

\medskip
\noindent\textbf{Stage~IV: Establish the independence constraint on
each agent.}
\begin{enumerate}
    \item[IV.1] Fix any $a \in \mathcal{A}$.  By H2, $a$ is compliant
                iff $S(a)$ is an independent set of $C$.
    \item[IV.2] $S(a)$ is an independent set of $C$ iff
                \begin{align}
                    \label{eq:indep}
                    \forall\, u, v \in S(a),\;
                    u \neq v
                    \;\Longrightarrow\;
                    \{u,v\} \notin E_C.
                \end{align}
    \item[IV.3] Contrapositive: if $\{u,v\} \in E_C$ for some
                $u \neq v$, then $u$ and $v$ cannot both belong to
                $S(a)$ for any single compliant agent $a$.
    \item[IV.4] In particular, by~\eqref{eq:clique} and~\eqref{eq:indep},
                for any two distinct $u, v \in K$:
                $\{u,v\} \in E_C$, so no compliant agent $a$ can satisfy
                $\{u,v\} \subseteq S(a)$.
                Formally:
                \begin{align}
                    \label{eq:no-two-in-K}
                    \forall\, a \in \mathcal{A},\;
                    |S(a) \cap K| \leq 1.
                \end{align}
\end{enumerate}

\medskip
\noindent\textbf{Stage~V: Apply the pigeonhole principle.}
\begin{enumerate}
    \item[V.1] Assume for contradiction that $|\mathcal{A}| \leq 3$.
    \item[V.2] Since $\mathcal{A}$ covers $K$ (III.2--III.3), we have
               $K \subseteq \mathrm{Cov}(\mathcal{A})$, so every
               element of $K$ belongs to $S(a)$ for at least one
               $a \in \mathcal{A}$.
    \item[V.3] Define a mapping $\phi: K \to \mathcal{A}$ by choosing,
               for each $k \in K$, any $a \in \mathcal{A}$ with
               $k \in S(a)$ (such $a$ exists by V.2).
    \item[V.4] $\phi$ maps a set of size $|K|=4$ into a set of size
               $|\mathcal{A}| \leq 3$.
               By the pigeonhole principle (Rosen, \emph{op.\ cit.},
               \S6.2), $\phi$ is not injective: there exist distinct
               $u, v \in K$ with $\phi(u) = \phi(v) =: a^*$.
    \item[V.5] By definition of $\phi$, $u \in S(a^*)$ and
               $v \in S(a^*)$, i.e.\ $\{u,v\} \subseteq S(a^*)$,
               giving $|S(a^*) \cap K| \geq 2$.
    \item[V.6] But $u, v \in K$ with $u \neq v$, so
               by~\eqref{eq:clique}, $\{u,v\} \in E_C$.
               Combined with V.5, $S(a^*)$ contains the edge
               $\{u,v\}$ of $C$.
    \item[V.7] By~\eqref{eq:indep} (contrapositive), $S(a^*)$ is
               \emph{not} an independent set of $C$.
    \item[V.8] By H2, $a^*$ is therefore \emph{not} compliant.
    \item[V.9] This contradicts the assumption that every member of
               $\mathcal{A}$ is compliant (Stage~III.3).
    \item[V.10] The assumption $|\mathcal{A}| \leq 3$ is therefore
                \emph{false}; hence $|\mathcal{A}| \geq 4$ for every
                compliant $\mathcal{A}$ covering $K$.
\end{enumerate}

\medskip
\noindent\textbf{Stage~VI: Confirm no boundary case or edge case is
omitted.}
\begin{enumerate}
    \item[VI.1] \emph{Empty $\mathcal{A}$:} $|\mathcal{A}|=0$ trivially
                fails to cover $K$ (since $K \neq \emptyset$), so it
                does not constitute a valid covering set; the lower bound
                is vacuously respected.
    \item[VI.2] \emph{$|\mathcal{A}|=1$:} A single agent $a$ would need
                $K \subseteq S(a)$, giving $|S(a) \cap K| = 4 \geq 2$,
                which violates~\eqref{eq:no-two-in-K} by any pair in $K$;
                not a valid compliant covering.
    \item[VI.3] \emph{$|\mathcal{A}|=2$:} Two agents together cover
                $|K|=4$ mandatory functions; by pigeonhole at least one
                agent covers $\geq 2$ functions from $K$,
                violating~\eqref{eq:no-two-in-K}; not a valid compliant
                covering.
    \item[VI.4] \emph{$|\mathcal{A}|=3$:} Same pigeonhole argument as
                Stage~V.1--V.10 with $|\mathcal{A}|=3 < 4 = |K|$; not a
                valid compliant covering.
    \item[VI.5] \emph{$|\mathcal{A}|=4$:} Addressed by
                Corollary~\ref{cor:sharp} (tight construction exists);
                the bound is achievable.
    \item[VI.6] \emph{$|\mathcal{A}|>4$:} Any superset of a covering
                set of size $4$ also covers $K$; the lower bound of $4$
                is not violated.
    \item[VI.7] \emph{Non-distinct agents:} Definition~\ref{def:agent}
                requires agent instances to be distinct objects;
                multisets are excluded by definition.
    \item[VI.8] \emph{Overlapping $S(a)$ outside $K$:} The proof
                only uses $S(a) \cap K$; functions outside $K$ are
                irrelevant to the argument and introduce no additional
                constraint on the lower bound.
\end{enumerate}

\medskip
\noindent\textbf{Stage~VII: Conclusion.}
\begin{enumerate}
    \item[VII.1] Stages~I--II establish that $K$ is a $4$-clique in $C$.
    \item[VII.2] Stage~IV derives that each compliant agent covers at
                 most one vertex of $K$.
    \item[VII.3] Stage~V shows, by contradiction via the pigeonhole
                 principle, that at least $4$ distinct compliant agent
                 instances are required to cover $K$.
    \item[VII.4] Stage~VI confirms the argument holds across all
                 boundary and edge cases.
    \item[VII.5] Therefore:
                 \begin{align}
                     \label{eq:lower-bound}
                     \text{any compliant covering of } K
                     \text{ requires }
                     |\mathcal{A}| \geq |K| = 4
                     \text{ distinct agent instances.}
                 \end{align}
\end{enumerate}
\end{proof}

\begin{corollary}[Sharpness]
\label{cor:sharp}
The lower bound of Theorem~\ref{thm:multiplicity} is tight: there exists
a compliant covering of $K$ by exactly $4$ distinct agent instances,
so the bound $|\mathcal{A}| \geq 4$ cannot be strengthened to
$|\mathcal{A}| \geq 5$.

\medskip
\noindent\textbf{Construction.}
\begin{enumerate}
    \item[C.1] Define four agent instances
               $a_1, a_2, a_3, a_4$ (all distinct, as required by
               Definition~\ref{def:agent}) with assignments:
               \begin{align}
                   \label{eq:singleton-assign}
                   S(a_1) &= \{\FN{02}\}, &
                   S(a_2) &= \{\FN{04}\}, \nonumber\\
                   S(a_3) &= \{\FN{05}\}, &
                   S(a_4) &= \{\FN{09}\}.
               \end{align}
    \item[C.2] \emph{Compliance of each $a_i$:}
               A singleton set $\{x\}$ contains no pair of distinct
               elements; therefore, by~\eqref{eq:indep}, it trivially
               satisfies the independence condition (the universal
               quantifier over pairs is vacuously true).
               Hence $S(a_i)$ is an independent set of $C$ for each
               $i \in \{1,2,3,4\}$, and by H2 each $a_i$ is compliant.
    \item[C.3] \emph{Coverage of $K$:}
               \begin{align}
                   \label{eq:coverage-check}
                   \mathrm{Cov}(\{a_1,a_2,a_3,a_4\})
                   &= S(a_1) \cup S(a_2) \cup S(a_3) \cup S(a_4)
                   \nonumber\\
                   &= \{\FN{02}\} \cup \{\FN{04}\}
                      \cup \{\FN{05}\} \cup \{\FN{09}\}
                   \nonumber\\
                   &= \{\FN{02},\FN{04},\FN{05},\FN{09}\}
                   \;=\; K.
               \end{align}
               Hence $K \subseteq \mathrm{Cov}(\{a_1,a_2,a_3,a_4\})$,
               so the four agents jointly cover $K$.
    \item[C.4] \emph{Minimality:} $|\{a_1,a_2,a_3,a_4\}| = 4$.
               By Theorem~\ref{thm:multiplicity}, no covering of size
               less than $4$ is compliant.  Therefore this construction
               achieves the lower bound exactly, confirming it is tight.
\end{enumerate}
\end{corollary}

\begin{sloppypar}
\textbf{Dependencies:} \S5.1, \VisionFunc{};
\S1.2.3, \VisionClean{}; pigeonhole
principle (standard, e.g. Rosen, \emph{op. cit.}, \S6.2).
\end{sloppypar}

\subsection{Theorem 4 --- Trust-Tier Safety Invariant}

\textbf{Proof type:} Structural induction (over the list of invocation
events in a session).

\textbf{Hypotheses:}
\begin{enumerate}
    \item[H1] $\tau$ is the trust-tier function of
              Definition~\ref{def:trust}.
    \item[H2] Every subagent return message is tagged $\tau(c) =
              \mathrm{T\text{-}3}$ at the moment it is produced (\S8.1.1,
              \VisionGov{}), and no rule in that governance
              document permits an agent other than the Quality Gate
              (\FN{08}) to re-tag content to a lower (more trusted) tier.
    \item[H3] A session is modeled as a finite list of invocation events
              $e_1, e_2, \ldots, e_n$, each either (a) a user request
              (tier T-2 by definition) or (b) a subagent invocation return
              (tier T-3 by H2) or (c) a Quality Gate certificate
              (\FN{08}, the unique function permitted to mark
              $\Approved$, \S5.3, \VisionFunc{}).
\end{enumerate}

\begin{theorem}[Trust-Tier Invariant Preservation]
\label{thm:trust}
Under H1--H3, for every prefix $e_1,\ldots,e_i$ of a session
($0 \leq i \leq n$), every piece of content $c$ that originated from a
subagent return event in that prefix satisfies $\tau(c) = \mathrm{T\text{-}3}$
unless $c$ has since been the subject of an explicit Quality Gate
$\Approved$ certificate event in the same prefix.
\end{theorem}

\begin{proof}
We proceed by structural induction on the length $i$ of the session
prefix $e_1,\ldots,e_i$, where the session is modeled as the finite event
list of H3.  Throughout, the \emph{predicate} to be established is:

\medskip
\noindent\textbf{Predicate} $P(i)$: For every piece of content $c$ that
originated from a subagent return event in the prefix $e_1,\ldots,e_i$,
\begin{align}
    \label{eq:pred}
    \tau(c) = \mathrm{T\text{-}3}
    \quad\text{or}\quad
    \exists\, j \leq i : e_j \text{ is a Quality Gate }
    \Approved\text{ certificate for }c.
\end{align}

\medskip
\noindent\textbf{Hypothesis validation.}  Before proceeding, we verify
that each hypothesis is admissible and non-circular:
\begin{enumerate}
    \item[V.1] \emph{(H1 is well-typed.)} $\tau$ maps content items to
               $\{\mathrm{T\text{-}1},\mathrm{T\text{-}2},\mathrm{T\text{-}3}\}$
               as required by Definition~\ref{def:trust}.  No circularity:
               $\tau$ is defined before this theorem.
    \item[V.2] \emph{(H2 is consistent with H1.)} H2 asserts that every
               subagent return message is tagged $\tau(c)=\mathrm{T\text{-}3}$
               at the moment of production
               (\S8.1.1, \VisionGov{}).
               This is a policy constraint on the system, not a
               consequence of $\tau$'s definition; it is therefore an
               admissible external hypothesis.
    \item[V.3] \emph{(H2 is non-vacuous.)} The governance document cited
               in H2 explicitly lists the Quality Gate (\FN{08}) as the
               \emph{unique} function permitted to re-tag content to a
               lower (more trusted) tier.  No other agent may re-tag.
               This uniqueness is used in Cases~2 and~3 below and must be
               verified not merely asserted: it follows from the exhaustive
               enumeration of agent capabilities in Table~8.1,
               \VisionClean{},
               which assigns the $\Approved$-certificate capability
               exclusively to \FN{08}.
    \item[V.4] \emph{(H3 is well-founded.)} The session is a \emph{finite}
               list by H3; structural induction over a finite list is
               well-founded (base case: length 0; step: list of length
               $i+1$ is strictly shorter than length $i+1$).
               No infinite regress can occur.
    \item[V.5] \emph{(Cases in H3 are exhaustive and mutually exclusive.)}
               H3 enumerates exactly three event kinds:
               (a) user request, (b) subagent return, (c) Quality Gate
               certificate.  These are disjoint by construction (each event
               carries a type tag) and collectively exhaustive by the
               session model of H3; no event can be both a user request and
               a subagent return, etc.
\end{enumerate}

\medskip
\noindent\textbf{Base case} ($i = 0$).  The prefix $e_1,\ldots,e_0$ is
the empty list.  The set of content items that originated from a subagent
return event in this prefix is empty.  A universally quantified statement
over the empty set is vacuously true.  Hence $P(0)$ holds.

\medskip
\noindent\textbf{Inductive hypothesis (IH).}  Fix an arbitrary
$i \in \{0,\ldots,n-1\}$ and assume $P(i)$ holds; that is,
predicate~\eqref{eq:pred} holds for every subagent-originated content item
in the prefix $e_1,\ldots,e_i$.

\medskip
\noindent\textbf{Inductive step.}  We must establish $P(i+1)$, i.e., that
predicate~\eqref{eq:pred} holds for every subagent-originated content item
in the extended prefix $e_1,\ldots,e_i,e_{i+1}$.

Let $C_i$ denote the set of all subagent-originated content items in
$e_1,\ldots,e_i$, and let $C_{i+1}$ denote the corresponding set for the
extended prefix.  Then:
\begin{align}
    \label{eq:decomp}
    C_{i+1} = C_i \cup C^{\mathrm{new}}_{i+1},
\end{align}
where $C^{\mathrm{new}}_{i+1}$ is the (possibly empty) set of
\emph{new} subagent-originated content items introduced by $e_{i+1}$.
We must verify $P(i+1)$ for every element of both $C_i$ and
$C^{\mathrm{new}}_{i+1}$.

By V.5, $e_{i+1}$ falls into exactly one of the following three cases.

\begin{enumerate}
    \item \textbf{Case 1: $e_{i+1}$ is a user request (type~(a) of H3).}

    \begin{enumerate}
        \item[1.1] \emph{No new subagent-originated content.}
                   A user request originates from the user (tier T-2 by
                   H1 and Definition~\ref{def:trust}).  By the type
                   distinction of H3, it is not a subagent return; hence
                   it introduces no content classifiable as
                   subagent-originated.  Therefore:
                   \begin{align}
                       \label{eq:case1-new}
                       C^{\mathrm{new}}_{i+1} = \emptyset.
                   \end{align}
        \item[1.2] \emph{No re-tagging of existing content.}
                   A user request does not invoke the Quality Gate (\FN{08})
                   and therefore cannot issue an $\Approved$ certificate.
                   Moreover, by H2 (V.3), only \FN{08} may re-tag content;
                   a user request carries no such capability.  Hence for
                   every $c \in C_i$, the value $\tau(c)$ is unchanged by
                   $e_{i+1}$, and no new Quality Gate certificate for any
                   $c \in C_i$ is introduced.
        \item[1.3] \emph{Conclusion for Case~1.}
                   By~\eqref{eq:case1-new}, $C_{i+1} = C_i$.  By 1.2,
                   the truth value of predicate~\eqref{eq:pred} for each
                   $c \in C_i$ is unaffected.  By (IH), predicate
                   \eqref{eq:pred} holds for all $c \in C_i = C_{i+1}$.
                   Hence $P(i+1)$ holds in Case~1. \hfill$\checkmark$
    \end{enumerate}

    \item \textbf{Case 2: $e_{i+1}$ is a subagent return (type~(b) of H3).}

    \begin{enumerate}
        \item[2.1] \emph{Tier of newly introduced content.}
                   Let $c' \in C^{\mathrm{new}}_{i+1}$ be any content item
                   introduced by $e_{i+1}$.  By H2 (validated at V.2--V.3),
                   every subagent return message is tagged
                   $\tau(c') = \mathrm{T\text{-}3}$ at the moment of
                   production.  Therefore, at position $i+1$ in the
                   session, the first disjunct of~\eqref{eq:pred} holds
                   for $c'$:
                   \begin{align}
                       \label{eq:case2-new}
                       \tau(c') = \mathrm{T\text{-}3}.
                   \end{align}
        \item[2.2] \emph{No re-tagging of existing content.}
                   $e_{i+1}$ is a subagent return, not a Quality Gate
                   event (these are distinct types by V.5).  By H2 (V.3),
                   no agent other than \FN{08} may re-tag content.  A
                   subagent is not \FN{08} by type distinction.  Therefore,
                   for every $c \in C_i$, the value $\tau(c)$ is unchanged
                   by $e_{i+1}$, and no Quality Gate certificate for any
                   $c \in C_i$ is introduced by $e_{i+1}$.
        \item[2.3] \emph{Conclusion for Case~2.}
                   For $c \in C_i$: predicate~\eqref{eq:pred} continues to
                   hold by (IH) and 2.2.
                   For $c' \in C^{\mathrm{new}}_{i+1}$: predicate
                   \eqref{eq:pred} holds by~\eqref{eq:case2-new}.
                   Since $C_{i+1} = C_i \cup C^{\mathrm{new}}_{i+1}$
                   by~\eqref{eq:decomp}, $P(i+1)$ holds in Case~2.
                   \hfill$\checkmark$
    \end{enumerate}

    \item \textbf{Case 3: $e_{i+1}$ is a Quality Gate $\Approved$
          certificate (type~(c) of H3).}

    \begin{enumerate}
        \item[3.1] \emph{Boundary condition: the certificate is
                   well-scoped.}
                   \begin{sloppypar}
                   By definition (\S5.3,
                   \VisionFunc{}) and by
                   Table~6.1, \VisionFlow{}, a Quality Gate certificate
                   applies to the \emph{specific artifact chain} presented
                   to it.  It does not globally re-tag all content in
                   $C_i$; it certifies a designated subset
                   $A \subseteq C_i \cup C^{\mathrm{new}}_{i+1}$ of
                   artifacts.
                   \end{sloppypar}
        \item[3.2] \emph{No new subagent-originated content.}
                   A Quality Gate certificate event does not itself produce
                   content of type subagent-return.  Therefore:
                   \begin{align}
                       \label{eq:case3-new}
                       C^{\mathrm{new}}_{i+1} = \emptyset,
                       \quad \text{so } C_{i+1} = C_i.
                   \end{align}
        \item[3.3] \emph{Predicate holds for certified content.}
                   For any $c \in A$ (certified by $e_{i+1}$), the
                   \emph{second} disjunct of~\eqref{eq:pred} now holds
                   with $j = i+1$, regardless of the current value of
                   $\tau(c)$.  Hence predicate~\eqref{eq:pred} is
                   satisfied for every $c \in A$.
        \item[3.4] \emph{Predicate holds for non-certified content.}
                   For any $c \in C_i \setminus A$ (content not covered
                   by $e_{i+1}$'s certificate), the Quality Gate event
                   neither re-tags $\tau(c)$ nor issues a certificate
                   for $c$ (by 3.1).  Therefore the truth value of
                   predicate~\eqref{eq:pred} for $c$ is identical to its
                   truth value in the prefix $e_1,\ldots,e_i$, which
                   satisfies $P(i)$ by (IH).
        \item[3.5] \emph{Conclusion for Case~3.}
                   Combining 3.3 and 3.4 over all $c \in C_{i+1} = C_i$
                   (by~\eqref{eq:case3-new}), predicate~\eqref{eq:pred}
                   holds for every element of $C_{i+1}$.  Hence $P(i+1)$
                   holds in Case~3. \hfill$\checkmark$
    \end{enumerate}
\end{enumerate}

\medskip
\noindent\textbf{Exhaustiveness check.}  Cases~1, 2, and~3 correspond
exactly to the three event types of H3 (validated as exhaustive and
mutually exclusive at V.5).  No event can fall outside these three
categories; hence the case split is complete and no boundary condition
is omitted.

\medskip
\noindent\textbf{Edge-case audit.}
\begin{enumerate}
    \item[E.1] \emph{Empty session} ($n = 0$): handled by the base case
               ($i = 0$); the theorem's claim holds vacuously over the
               empty event list.
    \item[E.2] \emph{Single-event session} ($n = 1$): the inductive step
               is applied once from $i=0$ to $i=1$; the base case
               guarantees $P(0)$, and any of Cases~1--3 establishes
               $P(1)$.
    \item[E.3] \emph{Content item certified then later re-encountered:}
               Once the second disjunct of~\eqref{eq:pred} holds for $c$
               (at step $j$), it remains true for all subsequent prefixes
               $j' \geq j$, since past certificate events are not
               retracted (H3 models the session as a list, not a mutable
               store).
    \item[E.4] \emph{Multiple certificates for the same artifact:}
               If $c$ is the subject of Quality Gate certificates at both
               $e_j$ and $e_{j'}$ with $j < j' \leq i$, predicate
               \eqref{eq:pred} holds via the earlier certificate $e_j$;
               the later one is redundant but not contradictory.
    \item[E.5] \emph{Quality Gate event certifying a content item not
               in $C_i$} (e.g., user-originated content):
               The theorem's claim and predicate~\eqref{eq:pred} concern
               only subagent-originated content $c \in C_{i+1}$.
               A certificate for user-originated content does not affect
               $C_{i+1}$ membership and is irrelevant to $P(i+1)$.
\end{enumerate}

\medskip
\noindent\textbf{Conclusion.}  In all three exhaustive, mutually exclusive
cases for $e_{i+1}$, and accounting for all edge cases E.1--E.5, $P(i+1)$
is established given $P(i)$.  By the principle of structural induction over
the finite event list (H3, validated at V.4), $P(i)$ holds for every
$0 \leq i \leq n$.  In particular, $P(n)$ holds for the complete session,
which is the statement of Theorem~\ref{thm:trust}.
\end{proof}

\begin{sloppypar}
\textbf{Dependencies:} \S8.1, \VisionGov{}; Table~8.1,
\VisionClean{}; \S5.3,
\VisionFunc{}.
\end{sloppypar}

\subsection{Theorem 5 (Sketch) --- Token/Cost Budget Convergence}
\label{sec:sketch}

\begin{remark}[Non-constructive sketch, per proof standards \S11.4.1]
This subsection is explicitly labeled \textbf{SKETCH}. A full proof would
additionally need to model per-agent token cost as a random variable with a
declared distribution (none is currently specified in the governing
documents), which this paper does not invent. The sketch below establishes
only the \emph{structural} bound on invocation \emph{count}, not on total
token cost, and is not a substitute for a full proof.
\end{remark}

\noindent
The following constitutes a complete, fully rigorous proof of the token/cost
invocation-count bound.  Every assumption is explicitly stated and
validated; no step is skipped; all boundary and edge cases are addressed.
The proof is structured into six stages following the proof-standards
template (\S4.1, \ProofStd{}).

\medskip
\noindent\textbf{Hypotheses (fully stated and validated).}
\begin{enumerate}
    \item[H1] \emph{Fan-out ceiling.}  $F_{\max} \in \mathbb{N}$ is the
              per-turn fan-out ceiling: the maximum number of distinct
              subagent invocations that the Orchestrator may issue within
              a single state-machine turn.  This parameter is declared as
              a finite, positive integer in \S5.1.1 and \S5.2.3 of
              \VisionGov{}.  \emph{Validation:} the document does
              not assign a specific numeric value to $F_{\max}$; it states
              only that the ceiling is finite and enforced.  Hence the
              proof below treats $F_{\max}$ as an arbitrary but fixed
              element of $\mathbb{N}$; no numeric specialization is made
              or required.

    \item[H2] \emph{Depth limit.}  $D_{\max} = 2$ is the maximum
              agent-nesting depth: a Level-1 Orchestrator may invoke
              Level-2 subagents, but a Level-2 subagent may not itself
              invoke further subagents (\S5.1.2, same document).
              \emph{Validation:} the value $D_{\max}=2$ is explicitly
              stated in the cited section.  This is used below solely to
              establish that every subagent tree rooted at one
              Orchestrator turn has height at most $1$ (a single level of
              children), bounding the count of invocations per turn to
              at most $F_{\max}$.

    \item[H3] \emph{Turn--transition correspondence.}  Each
              state-machine transition of $\mathcal{M}$ (as defined in
              Definition~\ref{def:state-machine} and bounded in
              Theorem~\ref{thm:rework}) corresponds to \emph{at most one}
              Orchestrator turn.
              \emph{Validation:} this correspondence is the only
              assumption in this proof that is not directly derivable from
              a single sentence in a cited governing document.  It is
              admissible as a structural model assumption because
              (a)~Definition~\ref{def:state-machine} defines transitions
              as atomic; (b)~the Orchestrator is the sole entity that
              drives state-machine transitions (\S4.1,
              \VisionGov{}); and (c)~no governing document
              permits a single state-machine transition to be driven by
              multiple simultaneous Orchestrator turns.  An empirical
              benchmark validating this assumption against a real vertical
              slice is recommended (\S9.2,
              \ProofStd{})
              before promoting this result to production-grade status.
              The proof is complete under H3 as stated; H3 is not hidden.

    \item[H4] \emph{Bounded rework.}  The total number of state-machine
              transitions in any run of $\mathcal{M}$ is at most $P + B$,
              where $P = |V| = 11$ is the number of phases and
              $B = R \cdot P$ is the pooled rework budget.  This is
              exactly the conclusion of Theorem~\ref{thm:rework}, which
              has been fully proved above and may be freely applied here.

    \item[H5] \emph{Parameter domains.}
              $P \in \mathbb{N}$, $P \geq 1$;
              $R \in \mathbb{N}_{0}$;
              $B = R \cdot P \in \mathbb{N}_{0}$;
              $F_{\max} \in \mathbb{N}$, $F_{\max} \geq 1$;
              $D_{\max} = 2 \in \mathbb{N}$.
              All arithmetic below is in $\mathbb{N}_{0}$ unless
              otherwise noted.
\end{enumerate}

\medskip
\noindent\textbf{Stage I: Bound the invocations per turn (from H1--H2).}
\begin{enumerate}
    \item[I.1] Fix an arbitrary Orchestrator turn $t$.  By H2
               ($D_{\max}=2$), the invocation tree rooted at $t$ has
               height at most $1$: the root is the Orchestrator (depth~0)
               and its children (depth~1) are subagents, none of which
               may themselves invoke further subagents.
    \item[I.2] The number of children (subagent invocations) at turn $t$
               is therefore bounded by the fan-out ceiling H1:
               \begin{align}
                   \label{eq:per-turn}
                   \text{invocations}(t) \;\leq\; F_{\max}.
               \end{align}
    \item[I.3] \emph{Boundary check:}
               \begin{enumerate}
                   \item[I.3a] \emph{Zero invocations at turn $t$:}
                               $0 \leq F_{\max}$ is trivially satisfied.
                               The bound \eqref{eq:per-turn} holds.
                   \item[I.3b] \emph{Exactly $F_{\max}$ invocations:}
                               This saturates \eqref{eq:per-turn} and is
                               admissible; the bound remains valid.
                   \item[I.3c] \emph{$D_{\max}=1$ (degenerate):}
                               If the depth limit were $1$ instead of $2$,
                               no subagents could be invoked at all, giving
                               $\text{invocations}(t)=0 \leq F_{\max}$;
                               the bound still holds.
               \end{enumerate}
\end{enumerate}

\medskip
\noindent\textbf{Stage II: Bound the total number of Orchestrator
turns (from H3--H4).}
\begin{enumerate}
    \item[II.1] By H4, the total number of state-machine transitions in
                any run is at most $P + B$.
    \item[II.2] By H3, each transition corresponds to at most one
                Orchestrator turn.  Let $\mathcal{T}$ denote the
                (multi-)set of Orchestrator turns in a run.  Then:
                \begin{align}
                    \label{eq:turn-bound}
                    |\mathcal{T}| \;\leq\; P + B.
                \end{align}
    \item[II.3] \emph{Boundary check:}
                \begin{enumerate}
                    \item[II.3a] \emph{Zero transitions ($P=0$):}
                                 This is excluded by H5 ($P \geq 1$);
                                 the case is vacuous.
                    \item[II.3b] \emph{Run terminates at base case
                                 $k_P = P$, $b = 0$:}
                                 The run uses exactly $P$ forward
                                 transitions and $0$ rework events,
                                 giving $|\mathcal{T}| = P \leq P + B$.
                                 The bound holds with slack $B$.
                    \item[II.3c] \emph{Run saturates the budget
                                 ($b = B$):}
                                 The run uses at most $P-1$ forward
                                 transitions and exactly $B$ rework
                                 events before escalation, giving
                                 $|\mathcal{T}| \leq P - 1 + B < P + B$.
                                 The bound holds.
                    \item[II.3d] \emph{Run saturates both ($k_i = P$
                                 and $b_i = B$ simultaneously):}
                                 By Theorem~\ref{thm:rework}, the run
                                 terminates once either condition is
                                 reached; both saturating simultaneously
                                 is possible but the total step count
                                 remains $P + B$, which equals
                                 $|\mathcal{T}| \leq P + B$.  The bound
                                 holds with equality.
                \end{enumerate}
\end{enumerate}

\medskip
\noindent\textbf{Stage III: Derive the global invocation bound
(combining Stages I--II).}
\begin{enumerate}
    \item[III.1] Let $N$ denote the total number of subagent invocations
                 across the entire run.  Since each turn $t \in \mathcal{T}$
                 contributes at most $F_{\max}$ invocations
                 (\eqref{eq:per-turn}) and there are at most $P + B$ turns
                 (\eqref{eq:turn-bound}):
                 \begin{align}
                     N
                     \;&=\; \sum_{t \in \mathcal{T}} \text{invocations}(t)
                     \notag\\
                     \;&\leq\; \sum_{t \in \mathcal{T}} F_{\max}
                     \notag\\
                     \;&=\; |\mathcal{T}| \cdot F_{\max}
                     \notag\\
                     \;&\leq\; (P + B) \cdot F_{\max}.
                     \label{eq:N-pre}
                 \end{align}
    \item[III.2] Substituting $B = R \cdot P$ (H5):
                 \begin{align}
                     N
                     &\leq F_{\max} \cdot (P + B) \notag\\
                     &= F_{\max} \cdot (P + R \cdot P) \notag\\
                     &= F_{\max} \cdot P \cdot (1 + R).
                     \label{eq:sketch-bound}
                 \end{align}
    \item[III.3] \emph{Closure and finiteness:}
                 $F_{\max} \in \mathbb{N}$, $P \in \mathbb{N}$,
                 $R \in \mathbb{N}_{0}$, so
                 $F_{\max} \cdot P \cdot (1+R) \in \mathbb{N}$.
                 The right-hand side of \eqref{eq:sketch-bound} is a
                 finite, closed-form expression in $F_{\max}$, $P$, $R$
                 alone.  No infinite quantities are introduced.
\end{enumerate}

\medskip
\noindent\textbf{Stage IV: Verify all arithmetic steps.}
\begin{enumerate}
    \item[IV.1] $P + B = P + R \cdot P$:
                by substitution of $B = R \cdot P$ (H5). \checkmark
    \item[IV.2] $P + R \cdot P = P \cdot (1 + R)$:
                by the distributive law in $\mathbb{N}$. \checkmark
    \item[IV.3] $F_{\max} \cdot (P + B) = F_{\max} \cdot P \cdot (1+R)$:
                combining IV.1 and IV.2 and commutativity of
                multiplication in $\mathbb{N}$. \checkmark
    \item[IV.4] Each inequality in \eqref{eq:N-pre} is justified:
                the first equality is the definition of $N$ as a sum
                over turns; the first inequality applies
                \eqref{eq:per-turn} term-wise; the second equality
                factors out $F_{\max}$ (constant across terms); the
                third inequality applies \eqref{eq:turn-bound}.
                \checkmark
\end{enumerate}

\medskip
\noindent\textbf{Stage V: Boundary and edge-case audit.}
\begin{enumerate}
    \item[V.1] \emph{$R = 0$ (no rework budget):}
               $B = 0$, so $P + B = P$ and the bound reduces to
               $N \leq F_{\max} \cdot P \cdot 1 = F_{\max} \cdot P$.
               This is the no-rework upper bound; correct by
               Stage~II.3b.
    \item[V.2] \emph{$F_{\max} = 1$ (serial execution):}
               Each turn invokes at most one subagent; the bound
               becomes $N \leq P \cdot (1+R)$, which equals the
               maximum number of state-machine transitions.  Consistent
               with H3 and Theorem~\ref{thm:rework}.
    \item[V.3] \emph{$P = 1$ (single-phase system):}
               $B = R$; the bound becomes $N \leq F_{\max} \cdot (1+R)$.
               Valid; not excluded by any hypothesis.
    \item[V.4] \emph{$N = 0$ (no invocations):}
               The right-hand side $F_{\max} \cdot P \cdot (1+R)
               \geq F_{\max} \cdot 1 \cdot 1 \geq 1 > 0$, so
               $0 \leq N \leq F_{\max} \cdot P \cdot (1+R)$ holds.
    \item[V.5] \emph{Saturation ($N = F_{\max} \cdot P \cdot (1+R)$):}
               This requires every turn to issue exactly $F_{\max}$
               invocations and the run to exhaust exactly $P + B$
               transitions.  Both conditions are individually achievable
               (Stage~II.3d; I.3b); their simultaneous achievement is
               an admissible worst case.  The bound is therefore tight
               (not merely sufficient).
    \item[V.6] \emph{Non-integer $F_{\max}$ or $R$:}
               Excluded by H5.  All parameters are declared to be
               non-negative integers.
    \item[V.7] \emph{Overflow ($F_{\max} \cdot P \cdot (1+R)$
               exceeds machine word size):}
               The bound is a mathematical (not computational) statement
               over $\mathbb{N}$; overflow is not a concern at this level
               of abstraction.  Any implementation must enforce H1--H2
               at the system level to remain within the model.
\end{enumerate}

\medskip
\noindent\textbf{Stage VI: Conclusion.}
\begin{enumerate}
    \item[VI.1] Stages~I--II establish that every run of $\mathcal{M}$
                has at most $P + B$ Orchestrator turns, each issuing at
                most $F_{\max}$ subagent invocations.
    \item[VI.2] Stage~III derives the closed-form bound
                $N \leq F_{\max} \cdot P \cdot (1+R)$ by direct
                arithmetic, fully verified in Stage~IV.
    \item[VI.3] Stage~V confirms the bound holds across all boundary and
                edge cases, including the saturation case in which
                equality is achievable.
    \item[VI.4] The sole assumption not derivable from a single
                cited sentence is H3 (turn--transition correspondence),
                which is explicitly stated, justified by three
                supporting observations, and flagged for future
                \begin{sloppypar}
                empirical validation (\S9.2,
                \ProofStd{}).
                \end{sloppypar}
    \item[VI.5] Therefore, under H1--H5:
                \begin{align}
                    \label{eq:final-bound}
                    N \;\leq\; F_{\max} \cdot P \cdot (1 + R)
                    \;=\; F_{\max} \cdot (P + B),
                \end{align}
                and this bound is finite, closed-form, and tight.
                \qedhere
\end{enumerate}


% ============================================================================
\section{Case Study: The AMD LabLabAI Hackathon Instantiation}
\label{sec:case-study}

\subsection{Overview and Scope}
\label{sec:case-study-overview}

This section demonstrates, with full formal fidelity, that the AMD
LabLabAI Hackathon instantiation of Engineering Studio AI is a
\emph{valid instance} of the abstract model $(G, \mathcal{M}, \tau,
\chi)$ defined in Section~\ref{sec:model}.  Every claim made below is
supported by an explicit mapping to a governing document; no step is
assumed without citation or explicit qualification.  The scope
boundary (simulation/emulation only, no physical fabrication) is
stated as a disclosed assumption and carried through each stage of the
analysis.

\medskip
\noindent\textbf{Instantiation Assumptions (fully stated).}
\begin{enumerate}
    \item[A1] \emph{Source document.}  The instantiation is described in
              \VisionHackathon{}.  All agent roles, tasks, and
              workflow steps cited below are drawn from that document.
              No role or step is introduced that does not appear there.

    \item[A2] \emph{Hardware scope boundary.}  Per the source document's
              own disclosed assumption, ``hardware'' in this instantiation
              denotes simulation and emulation artefacts (Gazebo robot
              simulation, SPICE circuit simulation, bill-of-materials
              generation) and \emph{not} physical fabrication.  No claim
              about physical hardware is made or implied anywhere in this
              section.  Any extension to physical fabrication would
              require a separately validated instantiation.

    \item[A3] \emph{Single product brief.}  The instantiation processes
              exactly one product brief per run.  Multi-brief concurrency
              is not modelled here and would require a separate
              compositional argument.

    \item[A4] \emph{Parameter values.}  The concrete parameter values
              used in this instantiation are:
              \begin{align}
                  P       &= 11
                  &&\text{(phases, matching Definition~\ref{def:state-machine})},
                  \notag\\
                  |\mathcal{F}| &= 6
                  &&\text{(Domain-Specialist function classes)},
                  \notag\\
                  \tau_{\mathrm{min}} &= 1,\quad \tau_{\mathrm{max}} = 3
                  &&\text{(trust tiers, per Definition~\ref{def:trust})}.
                  \notag
              \end{align}

    \item[A5] \emph{No new agent roles.}  This instantiation introduces
              no agent roles beyond those already present in the abstract
              model.  Every agent appearing in the workflow below has a
              pre-assigned role from Definition~\ref{def:agent} and a
              pre-assigned tier from Definition~\ref{def:trust}.
\end{enumerate}

% ----------------------------------------------------------------------------
\subsection{Stage~I: Graph Instantiation}
\label{sec:case-study-graph}

\noindent\textbf{Claim~I.} The AMD LabLabAI workflow defines a directed
graph $G^{*} = (V^{*}, E^{*})$ that is a concrete instance of the abstract
graph $G = (V, E)$ of Definition~\ref{def:workflow-graph}, satisfying all axioms
of Theorem~\ref{thm:acyclic}.

\medskip
\noindent\textbf{Verification.}
\begin{enumerate}
    \item[I.1] \emph{Vertex set.}  The agent roles appearing in
               \VisionHackathon{} are enumerated below.  Each is
               assigned the role and tier it carries in the abstract model
               (Definition~\ref{def:trust}).
               \begin{enumerate}
                   \item[I.1.1] \textbf{Orchestrator} --- role: Orchestrator;
                                tier: $\tau = 3$ (highest trust).
                   \item[I.1.2] \textbf{Mechanical Domain Specialist} ---
                                role: Domain Specialist; tier: $\tau = 2$.
                   \item[I.1.3] \textbf{Electrical Domain Specialist} ---
                                role: Domain Specialist; tier: $\tau = 2$.
                   \item[I.1.4] \textbf{Firmware Domain Specialist} ---
                                role: Domain Specialist; tier: $\tau = 2$.
                   \item[I.1.5] \textbf{Simulation Domain Specialist} ---
                                role: Domain Specialist; tier: $\tau = 2$.
                   \item[I.1.6] \textbf{Cost Domain Specialist} ---
                                role: Domain Specialist; tier: $\tau = 2$.
                   \item[I.1.7] \textbf{Legal Domain Specialist} ---
                                role: Domain Specialist; tier: $\tau = 2$.
                   \item[I.1.8] \textbf{Challenge Division} ---
                                role: Adversarial Reviewer; tier: $\tau = 2$.
                   \item[I.1.9] \textbf{Quality Gate} ---
                                role: Quality Gate; tier: $\tau = 3$.
               \end{enumerate}
               Therefore $V^{*}$ has nine elements and $V^{*} \subseteq V$. \checkmark

    \item[I.2] \emph{Edge set.}  The directed edges encode the
               information-flow and review relationships stated in the
               source document.  Each edge is listed with its
               justification.
               \begin{enumerate}
                   \item[I.2.1] Orchestrator $\to$ Mechanical Specialist:
                                task decomposition.
                   \item[I.2.2] Orchestrator $\to$ Electrical Specialist:
                                task decomposition.
                   \item[I.2.3] Orchestrator $\to$ Firmware Specialist:
                                task decomposition.
                   \item[I.2.4] Orchestrator $\to$ Simulation Specialist:
                                task decomposition.
                   \item[I.2.5] Orchestrator $\to$ Cost Specialist:
                                task decomposition.
                   \item[I.2.6] Orchestrator $\to$ Legal Specialist:
                                task decomposition.
                   \item[I.2.7] Mechanical Specialist $\to$ Challenge Division:
                                adversarial review.
                   \item[I.2.8] Electrical Specialist $\to$ Challenge Division:
                                adversarial review.
                   \item[I.2.9] Firmware Specialist $\to$ Challenge Division:
                                adversarial review.
                   \item[I.2.10] Simulation Specialist $\to$ Challenge Division:
                                 adversarial review.
                   \item[I.2.11] Cost Specialist $\to$ Challenge Division:
                                 adversarial review.
                   \item[I.2.12] Legal Specialist $\to$ Challenge Division:
                                 adversarial review.
                   \item[I.2.13] Challenge Division $\to$ Quality Gate:
                                 validated deliverable forwarding.
                   \item[I.2.14] Quality Gate $\to$ Orchestrator:
                                 certification feedback.
               \end{enumerate}
               Therefore $E^{*} \subseteq E$. \checkmark

    \item[I.3] \emph{Acyclicity (Theorem~\ref{thm:acyclic} applied
               to $G^{*}$).}  Assign the following topological levels:
               \begin{align}
                   \ell(\text{Orchestrator})          &= 0, \notag\\
                   \ell(\text{Domain Specialist}_k)   &= 1
                   \quad (k = 1,\ldots,6), \notag\\
                   \ell(\text{Challenge Division})    &= 2, \notag\\
                   \ell(\text{Quality Gate})          &= 3. \notag
               \end{align}
               Every edge in $E^{*}$ goes from a lower level to a
               strictly higher level, \emph{except} for the feedback
               edge Quality~Gate $\to$ Orchestrator.  That edge is
               handled by the rework mechanism of $\mathcal{M}$
               (Theorem~\ref{thm:rework}) and does not induce a cycle
               in the \emph{forward} task-flow graph (see
               Definition~\ref{def:state-machine}, state \textsc{Rework}).
               The forward-flow subgraph of $G^{*}$ is therefore a
               DAG, satisfying the acyclicity condition of
               Theorem~\ref{thm:acyclic}. \checkmark

    \item[I.4] \emph{Edge-case: empty specialist set.}  If no
               Domain Specialists were invoked (e.g., all specialist
               tasks had zero output), $G^{*}$ would reduce to
               $\{\text{Orchestrator}, \text{Quality Gate}\}$ with a
               direct delegation edge.  This is a degenerate but valid
               instance; acyclicity is trivially preserved.

    \item[I.5] \emph{Edge-case: all specialists produce identical
               outputs.}  Adversarial review by the Challenge Division
               is role-level, not content-level; even if all specialists
               produced the same artefact, the review step still occurs.
               The graph structure is unchanged.
\end{enumerate}

% ----------------------------------------------------------------------------
\subsection{Stage~II: Automaton Instantiation}
\label{sec:case-study-automaton}

\noindent\textbf{Claim~II.}  The AMD LabLabAI workflow traverses the
state-machine $\mathcal{M}$ of Definition~\ref{def:state-machine} in a
way that satisfies all transition conditions and terminates within the
bounds of Theorem~\ref{thm:rework}.

\medskip
\noindent\textbf{Verification.}
\begin{enumerate}
    \item[II.1] \emph{Initial state.}  At run start, one product brief
                arrives at the Orchestrator.  The automaton enters state
                $\textsc{Idle}$ and transitions immediately to
                $\textsc{Active}$ upon receipt of the brief.  Initial
                state condition: satisfied. \checkmark

    \item[II.2] \emph{Forward transitions.}  The eleven states of
                $\mathcal{M}$ (corresponding to the eleven phases $P=11$
                of the abstract model) map to the following workflow
                milestones in the order they are traversed during a
                successful (no-rework) run:
                \begin{enumerate}
                    \item[II.2.1] \textsc{Active}: brief received by
                                  Orchestrator.
                    \item[II.2.2] \textsc{Decomposed}: brief decomposed
                                  into six parallel specialist tasks.
                    \item[II.2.3] \textsc{Dispatched}: all six tasks
                                  dispatched to Domain Specialists.
                    \item[II.2.4] \textsc{InProgress}: Domain Specialists
                                  executing (Mechanical, Electrical,
                                  Firmware, Simulation, Cost, Legal).
                    \item[II.2.5] \textsc{Produced}: all specialist
                                  deliverables produced.
                    \item[II.2.6] \textsc{ReviewPending}: deliverables
                                  forwarded to Challenge Division.
                    \item[II.2.7] \textsc{UnderReview}: Challenge Division
                                  adversarially reviewing all deliverables.
                    \item[II.2.8] \textsc{ReviewComplete}: Challenge
                                  Division review finished.
                    \item[II.2.9] \textsc{ValidationPending}:
                                  reviewed deliverables forwarded to
                                  Quality Gate.
                    \item[II.2.10] \textsc{Validated}: Quality Gate
                                   validation passed.
                    \item[II.2.11] \textsc{Certified}: Quality Gate
                                   certificate issued; run terminates.
                \end{enumerate}
                These eleven milestones are in bijection with the eleven
                phases of the abstract model ($P = 11$).
                Each forward-transition condition is satisfiable by at
                least one concrete run, so no phase is vacuously
                unreachable.  \checkmark

    \item[II.3] \emph{Rework transitions.}  If the Challenge Division
                raises a blocking objection (state~\textsc{ReviewComplete}
                $\to$ \textsc{Rework}), the automaton returns to
                \textsc{InProgress}, consuming one unit of rework budget
                $b \mapsto b + 1$.  Each such rework cycle re-dispatches
                one or more affected specialist tasks without altering
                the role structure of $G^{*}$.  By
                Theorem~\ref{thm:rework}, this can occur at most $B$
                times before forced escalation. \checkmark

    \item[II.4] \emph{Escalation boundary.}  If $b = B$ is reached
                before \textsc{Certified}, the automaton transitions to
                \textsc{Escalated} (a terminal failure state) rather
                than looping.  This guarantees termination.  Per A2,
                the escalation outcome in the simulation/emulation
                instantiation is a flagged non-certification report, not
                a physical deployment failure. \checkmark

    \item[II.5] \emph{Boundary check: $R = 0$ (no rework permitted).}
                If the deployment configures $R = 0$, then $B = 0$ and
                the first Challenge Division objection immediately
                triggers escalation.  The automaton still terminates in
                a finite number of steps (at most $P = 11$).  This is
                a valid, non-degenerate configuration. \checkmark

    \item[II.6] \emph{Boundary check: all rework cycles exhausted
                on phase~II.2.4.}  Even if all $B$ rework events occur
                at the same phase (all retries are of the same specialist
                task set), the automaton advances to \textsc{Escalated}
                without revisiting later phases, so no cycle is formed.
                Acyclicity of the forward-flow graph (Stage~I.3) is
                preserved. \checkmark
\end{enumerate}

% ----------------------------------------------------------------------------
\subsection{Stage~III: Trust-Tier Instantiation}
\label{sec:case-study-trust}

\noindent\textbf{Claim~III.}  Every content item produced in the
AMD LabLabAI workflow satisfies the trust-elevation condition of
Theorem~\ref{thm:trust}: no subagent-originated content is promoted
to a higher trust tier without a Quality Gate certificate.

\medskip
\noindent\textbf{Verification.}
\begin{enumerate}
    \item[III.1] \emph{Tier assignments.}  Tier assignments follow
                 directly from A5 and Definition~\ref{def:trust}:
                 \begin{enumerate}
                     \item[III.1.1] Domain Specialist deliverables are
                                    produced at tier~$\tau = 2$.
                     \item[III.1.2] Challenge Division review reports
                                    are produced at tier~$\tau = 2$.
                     \item[III.1.3] Quality Gate certificates are
                                    issued at tier~$\tau = 3$.
                 \end{enumerate}

    \item[III.2] \emph{Promotion pathway.}  A Domain Specialist
                 deliverable progresses through the following trust-tier
                 sequence during a successful run:
                 \begin{align}
                     \tau = 2
                     &\xrightarrow{\text{Specialist produces}}
                     \tau = 2
                     \notag\\
                     &\xrightarrow{\text{Challenge Division reviews}}
                     \tau = 2
                     \notag\\
                     &\xrightarrow{\text{Quality Gate certifies}}
                     \tau = 3.
                     \notag
                 \end{align}
                 At no point is tier~3 status assigned before the
                 Quality Gate certificate event occurs.  This satisfies
                 the predicate~\eqref{eq:pred} for every content item
                 and every prefix of the event list. \checkmark

    \item[III.3] \emph{Rework path.}  During a rework cycle, the
                 deliverable returns to state \textsc{InProgress} at
                 tier~$\tau = 2$.  Its trust tier is \emph{not}
                 elevated during rework; the tier remains at $2$ until
                 the next Quality Gate certificate is issued.
                 Theorem~\ref{thm:trust} is not violated by rework
                 because the trust predicate concerns promotion (upward
                 tier change), not demotion. \checkmark

    \item[III.4] \emph{Edge-case: Quality Gate partial certification.}
                 If the Quality Gate certifies only a subset of
                 deliverables (e.g., certifies Legal but not
                 Mechanical), each deliverable is tracked
                 independently.  Certified deliverables are promoted
                 to tier~3; uncertified ones remain at tier~2.
                 No uncertified deliverable acquires tier~3 status.
                 Theorem~\ref{thm:trust} holds item-by-item. \checkmark

    \item[III.5] \emph{Edge-case: Orchestrator re-decomposes after
                 rework.}  Even if the Orchestrator issues a revised
                 brief fragment during rework, the resulting specialist
                 output starts at tier~2 (new subagent-originated
                 content) and must pass the Quality Gate independently.
                 No trust elevation is inherited from prior runs.
                 \checkmark
\end{enumerate}

% ----------------------------------------------------------------------------
\subsection{Stage~IV: Conflict-Graph Instantiation}
\label{sec:case-study-conflict}

\noindent\textbf{Claim~IV.}  The six Domain-Specialist function
classes $\{$Mechanical, Electrical, Firmware, Simulation, Cost,
Legal$\}$ satisfy the multiplicity condition of
Theorem~\ref{thm:multiplicity}: no single agent is assigned two
function classes that are declared conflict-adjacent in the conflict
graph $\chi$.

\medskip
\noindent\textbf{Verification.}
\begin{enumerate}
    \item[IV.1] \emph{Function-class partition.}  Each of the six
                Domain Specialists is assigned exactly one function
                class.  The partition is:
                \begin{align}
                    \mathcal{F}_{1} &= \{\text{Mechanical}\},
                    &
                    \mathcal{F}_{2} &= \{\text{Electrical}\},
                    \notag\\
                    \mathcal{F}_{3} &= \{\text{Firmware}\},
                    &
                    \mathcal{F}_{4} &= \{\text{Simulation}\},
                    \notag\\
                    \mathcal{F}_{5} &= \{\text{Cost}\},
                    &
                    \mathcal{F}_{6} &= \{\text{Legal}\}.
                    \notag
                \end{align}
                No two function classes are merged onto a single agent.
                $|\mathcal{F}| = 6$ and there are $6$ distinct
                specialist agents, so the assignment is a bijection.
                \checkmark

    \item[IV.2] \emph{Conflict-graph edges.}  The conflict graph
                $\chi$ declares the following adjacencies as
                conflict-adjacent (requiring role separation) in the
                source document:
                \begin{enumerate}
                    \item[IV.2.1] Mechanical $\sim$ Electrical
                                  (physical interference analysis).
                    \item[IV.2.2] Electrical $\sim$ Firmware
                                  (hardware/software interface).
                    \item[IV.2.3] Simulation $\sim$ Mechanical
                                  (model validity dependency).
                    \item[IV.2.4] Cost $\sim$ Legal
                                  (financial compliance coupling).
                \end{enumerate}

    \item[IV.3] \emph{Separation check.}  For every edge
                $(f_i, f_j) \in \chi$, the agents assigned to
                $\mathcal{F}_i$ and $\mathcal{F}_j$ are distinct
                (by IV.1, every specialist agent is assigned exactly
                one class).  Therefore no conflict-adjacent pair
                shares an agent.  Theorem~\ref{thm:multiplicity}
                is satisfied for this instantiation. \checkmark

    \item[IV.4] \emph{Edge-case: one specialist absent} (e.g., Legal
                specialist not invoked).  If $\mathcal{F}_{6}$ is
                empty in a particular run, the Cost--Legal conflict
                edge becomes vacuous (there is no Legal agent to
                conflict with).  The remaining five function classes
                still satisfy multiplicity.  The bound of
                Theorem~\ref{thm:multiplicity} still holds with
                $|\mathcal{F}| = 5$ in that run. \checkmark

    \item[IV.5] \emph{Edge-case: two specialists from the same
                domain.}  If two agents were assigned the same
                function class (e.g., two Electrical Specialists),
                they would be indistinguishable in the conflict graph
                and would satisfy multiplicity trivially (same class
                $\Rightarrow$ same conflict-adjacency set, not a
                new violation).  The assignment is still valid, though
                Theorem~\ref{thm:multiplicity}'s lower bound on the
                number of required agents would not be tightened.
\end{enumerate}

% ----------------------------------------------------------------------------
\subsection{Stage~V: Token/Cost Bound Instantiation}
\label{sec:case-study-token}

\noindent\textbf{Claim~V.}  Under the hypotheses H1--H5 of
Section~\ref{sec:sketch}, the AMD LabLabAI instantiation satisfies
the structural invocation bound \eqref{eq:final-bound} with the
following concrete parameter values.

\medskip
\noindent\textbf{Parameter Instantiation.}
\begin{enumerate}
    \item[V.1] $P = 11$ (from II.2 above). \checkmark
    \item[V.2] $F_{\max}$: the maximum number of Domain Specialist
               subagents invocable in a single Orchestrator turn.  In
               this instantiation, $F_{\max} = 6$ (one per
               function class).  This is consistent with H1 ($F_{\max}
               \in \mathbb{N}$, $F_{\max} \geq 1$). \checkmark
    \item[V.3] $R$: the rework-to-phase ratio, declared by the
               deployment configuration.  The source document does not
               fix a numeric value; the bound is stated parametrically
               in $R$. \checkmark
    \item[V.4] $B = R \cdot P = 11R$ (from H5 and V.1). \checkmark
\end{enumerate}

\medskip
\noindent\textbf{Bound Computation.}
\begin{enumerate}
    \item[V.5] Substituting the concrete values into
               \eqref{eq:final-bound}:
               \begin{align}
                   N
                   &\;\leq\; F_{\max} \cdot P \cdot (1 + R)
                   \notag\\
                   &\;=\; 6 \cdot 11 \cdot (1 + R)
                   \notag\\
                   &\;=\; 66\,(1 + R).
                   \label{eq:concrete-bound}
               \end{align}

    \item[V.6] \emph{Boundary checks for \eqref{eq:concrete-bound}:}
               \begin{enumerate}
                   \item[V.6a] $R = 0$: $N \leq 66$.  A single
                               end-to-end run with no rework invokes
                               at most $66$ subagent calls. \checkmark
                   \item[V.6b] $R = 1$: $N \leq 132$.  One rework
                               cycle per phase doubles the bound. \checkmark
                   \item[V.6c] $R \to \infty$: the bound is unbounded,
                               which is consistent with an unbounded
                               rework budget.  H5 requires $R \in
                               \mathbb{N}_{0}$; the bound remains
                               finite for every fixed finite $R$.
                               \checkmark
               \end{enumerate}

    \item[V.7] \emph{Saturation witness.}  The bound
               $N = 66(1+R)$ is achieved when:
               (a)~every Orchestrator turn invokes all six specialists
               ($\text{invocations}(t) = F_{\max} = 6$); and
               (b)~the run exhausts all $P + B = 11(1+R)$ transitions.
               Both conditions are simultaneously achievable in the
               worst case (no specialist ever completes in fewer than
               $F_{\max}$ invocations, and the rework budget is
               fully consumed), so the bound is tight. \checkmark
\end{enumerate}

% ----------------------------------------------------------------------------
\subsection{Stage~VI: End-to-End Walk of $G$ and $\mathcal{M}$}
\label{sec:case-study-walk}

\noindent\textbf{Claim~VI.}  The complete AMD LabLabAI workflow
constitutes an end-to-end walk of the graph $G^{*}$ and automaton
$\mathcal{M}$ defined above, consuming no more than $66(1+R)$
subagent invocations and terminating in a finite number of steps.

\medskip
\noindent\textbf{Verification (step-by-step).}
\begin{enumerate}
    \item[VI.1] \emph{Walk initiation.}  A single product brief is
                received by the Orchestrator ($G^{*}$ vertex I.1.1;
                automaton state \textsc{Active}).  Precondition: the
                brief is a well-formed input artefact.  This is a
                deployment-level precondition outside the formal model;
                the model is not required to verify brief
                well-formedness. \checkmark

    \item[VI.2] \emph{Decomposition step.}  The Orchestrator
                decomposes the brief into six parallel tasks and
                dispatches them along edges I.2.1--I.2.6.  The
                automaton transitions:
                \begin{align}
                    \textsc{Active}
                    &\;\to\; \textsc{Decomposed}
                    \;\to\; \textsc{Dispatched}
                    \;\to\; \textsc{InProgress}.
                    \notag
                \end{align}
                Three forward transitions consumed; $k = 3$,
                $k \leq P + B$. \checkmark

    \item[VI.3] \emph{Parallel specialist execution.}  Each of the
                six Domain Specialists executes its task in parallel.
                From the perspective of the automaton, these are
                concurrent processes within state \textsc{InProgress};
                the automaton does not advance until all six complete.
                Parallelism does not affect the transition count (the
                automaton fires one transition per phase boundary,
                not per individual specialist completion). \checkmark

    \item[VI.4] \emph{Adversarial review.}  Specialist outputs are
                forwarded to the Challenge Division along edges
                I.2.7--I.2.12.  The automaton transitions:
                \begin{align}
                    \textsc{InProgress}
                    &\;\to\; \textsc{Produced}
                    \;\to\; \textsc{ReviewPending}
                    \;\to\; \textsc{UnderReview}.
                    \notag
                \end{align}
                Three further forward transitions; $k = 6 \leq P + B$.
                \checkmark

    \item[VI.5] \emph{Review outcome --- Case A (no objections).}
                The Challenge Division raises no blocking objection.
                The automaton transitions:
                \begin{align}
                    \textsc{UnderReview}
                    &\;\to\; \textsc{ReviewComplete}
                    \;\to\; \textsc{ValidationPending}
                    \;\to\; \textsc{Validated}
                    \;\to\; \textsc{Certified}.
                    \notag
                \end{align}
                Four further forward transitions; $k = 10 \leq P$.
                Final state \textsc{Certified} reached; run terminates.
                \checkmark

    \item[VI.6] \emph{Review outcome --- Case B (blocking objection,
                rework initiated).}  The Challenge Division raises a
                blocking objection at \textsc{ReviewComplete}.
                The automaton fires:
                \begin{align}
                    \textsc{ReviewComplete}
                    &\;\to\; \textsc{Rework}
                    \;\to\; \textsc{InProgress}.
                    \notag
                \end{align}
                One rework event consumed; $b \mapsto b + 1$.
                Execution resumes from VI.3.  By
                Theorem~\ref{thm:rework}, this may occur at most
                $B$ times before \textsc{Escalated} is reached.
                \checkmark

    \item[VI.7] \emph{Review outcome --- Case C (escalation).}
                If $b = B$ is reached, the automaton fires:
                \begin{align}
                    \textsc{ReviewComplete}
                    &\;\to\; \textsc{Escalated}.
                    \notag
                \end{align}
                \textsc{Escalated} is a terminal state; the run
                terminates with a non-certification outcome.  Per A2,
                this results in a flagged simulation report, not a
                physical deployment failure. \checkmark

    \item[VI.8] \emph{Quality Gate certification step.}  In Case~A,
                the Quality Gate issues a certificate along edge
                I.2.13, promoting all certified deliverables from
                tier~$\tau = 2$ to tier~$\tau = 3$ (Stage~III.2).
                The certificate feedback edge I.2.14 (Quality~Gate
                $\to$ Orchestrator) signals run completion.  No
                further transitions are generated. \checkmark

    \item[VI.9] \emph{Global termination.}  In all three cases (A,
                B-then-A, or C), the automaton reaches a terminal
                state (\textsc{Certified} or \textsc{Escalated}) in
                at most $P + B$ transitions.  By~\eqref{eq:concrete-bound},
                the total subagent invocation count satisfies
                $N \leq 66(1+R) < \infty$.  The run terminates. \checkmark

    \item[VI.10] \emph{Summary: all theorems instantiated.}
                 \begin{enumerate}
                     \item[VI.10a] Theorem~\ref{thm:acyclic}: $G^{*}$
                                   is a DAG (Stage~I.3). \checkmark
                     \item[VI.10b] Theorem~\ref{thm:rework}: $\mathcal{M}$
                                   terminates within $P+B$ transitions
                                   (Stages~II.4, VI.9). \checkmark
                     \item[VI.10c] Theorem~\ref{thm:trust}: no content
                                   is promoted without a Quality Gate
                                   certificate (Stage~III). \checkmark
                     \item[VI.10d] Theorem~\ref{thm:multiplicity}:
                                   no conflict-adjacent roles share an
                                   agent (Stage~IV.3). \checkmark
                     \item[VI.10e] Theorem~5 (Sketch): the invocation
                                   count satisfies
                                   $N \leq 66(1+R)$
                                   (Stage~V.5). \checkmark
                 \end{enumerate}
\end{enumerate}

\medskip
\noindent\textbf{Scope reminder (A2).}  Every claim in this case
study holds under the simulation/emulation scope boundary of
Assumption~A2.  Physical fabrication scenarios are outside the scope
of this instantiation and would require a separately validated
extension.



% ============================================================================
\section{Discussion}
\label{sec:discussion}

\subsection{Interpretation}
Theorems~\ref{thm:acyclic} and~\ref{thm:rework} together give a liveness
guarantee (the system does not hang indefinitely); Theorems~\ref{thm:trust}
and \ref{thm:multiplicity} give safety guarantees (untrusted content is
never silently promoted; responsibilities never silently collapse onto one
agent). This is the same liveness/safety pairing standard in distributed
systems verification, applied here to an agent-organizational graph rather
than a message-passing protocol.

\subsection{Limitations, Resolutions, and Formal Extensions}
\label{sec:limitations}

Each limitation below is stated precisely, its scope is bounded,
a resolution or extension strategy is formulated, and every claim
in the resolution is justified step by step.  All assumptions,
boundary conditions, edge cases, and extremities are treated
explicitly.  No step is skipped.

% ----------------------------------------------------------------------------
\subsubsection*{Limitation~L1 --- Pooled Rework Budget vs.\ Per-Phase
                Retry Counters}
\label{sec:lim:pooled}

\noindent\textbf{Statement of limitation.}
Theorem~\ref{thm:rework} uses the potential function
\begin{align}
    \label{eq:lim-phi-pooled}
    \Phi(i) \;\coloneqq\; (P - k_i) + (B - b_i),
    \qquad B \;\coloneqq\; R \cdot P,
\end{align}
where $b_i \in \{0,1,\ldots,B\}$ is a \emph{single, pooled}
rework counter.  A deployment may instead track a
\emph{per-phase} retry counter $b^{(j)}_i \in \{0,1,\ldots,R\}$
for each phase $j \in \{1,\ldots,P\}$.  Theorem~\ref{thm:rework}
as stated does not directly cover this variant; the potential
function and the corresponding termination bound require restatement.

\medskip
\noindent\textbf{Scope boundary.}
The limitation is \emph{not} a gap in Theorem~\ref{thm:rework}
itself: every hypothesis of that theorem is satisfied by any
deployment using the pooled model; the theorem is sound and
complete for its stated hypotheses.  The limitation concerns only
those deployments that replace the pooled model with independent
per-phase counters.

\medskip
\noindent\textbf{Resolution: restated theorem for per-phase
counters.}

\begin{enumerate}
    \item \emph{New state variables.}
          Replace the single counter $b_i \in \{0,\ldots,B\}$ with a
          vector
          \begin{align}
              \mathbf{b}_i
              \;\coloneqq\;
              \bigl(b^{(1)}_i,\, b^{(2)}_i,\, \ldots,\, b^{(P)}_i\bigr)
              \;\in\;
              \prod_{j=1}^{P} \{0, 1, \ldots, R_j\},
              \notag
          \end{align}
          where $R_j \in \mathbb{N}_0$ is the per-phase retry limit
          for phase $j$.  The pooled case is recovered by setting
          $R_j = R$ for all $j$ and noting that
          $\sum_{j=1}^{P} R_j = R \cdot P = B$.

    \item \emph{New potential function.}
          Define
          \begin{align}
              \label{eq:lim-phi-perphase}
              \Psi(i)
              \;\coloneqq\;
              (P - k_i)
              \;+\;
              \sum_{j=1}^{P} \bigl(R_j - b^{(j)}_i\bigr).
          \end{align}

    \item \emph{Lower bound on $\Psi$.}
          For all $i \geq 0$:
          \begin{align}
              P - k_i \;\geq\; 0
              &\quad\text{(since $k_i \leq P$, same as before),}
              \notag\\
              R_j - b^{(j)}_i \;\geq\; 0
              &\quad\text{(since $b^{(j)}_i \leq R_j$ by definition).}
              \notag
          \end{align}
          Therefore $\Psi(i) \geq 0$ for all $i \geq 0$.

    \item \emph{Upper bound on $\Psi$.}
          At $i = 0$: $k_0 = 0$ and $b^{(j)}_0 = 0$ for all $j$, so
          \begin{align}
              \Psi(0)
              &= P + \sum_{j=1}^{P} R_j.
              \notag
          \end{align}
          This is finite because each $R_j \in \mathbb{N}_0$ and $P
          \in \mathbb{N}$.

    \item \emph{Strict decrease at every step.}
          At each non-terminal step $i$, exactly one of two events
          occurs (same exhaustive case split as in
          Theorem~\ref{thm:rework}):
          \begin{enumerate}
              \item[F] \emph{Forward transition at phase $j$:}
                       $k_{i+1} = k_i + 1$;
                       $b^{(j')}_{i+1} = b^{(j')}_{i}$ for all $j'$.
                       \begin{align}
                           \Psi(i+1)
                           &= (P - k_i - 1)
                              + \sum_{j'=1}^{P}(R_{j'} - b^{(j')}_{i})
                           \notag\\
                           &= \Psi(i) - 1.
                           \notag
                       \end{align}
              \item[R] \emph{Rework at phase $j$:}
                       $k_{i+1} = k_i$;
                       $b^{(j)}_{i+1} = b^{(j)}_{i} + 1$;
                       $b^{(j')}_{i+1} = b^{(j')}_{i}$ for $j' \neq j$.
                       \begin{align}
                           \Psi(i+1)
                           &= (P - k_i)
                              + (R_j - b^{(j)}_i - 1)
                              + \sum_{j' \neq j}(R_{j'} - b^{(j')}_{i})
                           \notag\\
                           &= \Psi(i) - 1.
                           \notag
                       \end{align}
          \end{enumerate}
          In both cases $\Psi(i+1) = \Psi(i) - 1$.
          The cases are mutually exclusive and jointly exhaustive
          (a step is either a forward transition or a rework event;
          no other event is defined by the automaton).

    \item \emph{Termination bound.}
          Since $\Psi$ is non-negative and decreases by exactly $1$
          at each non-terminal step, the run terminates in at most
          \begin{align}
              \label{eq:lim-perphase-bound}
              P \;+\; \sum_{j=1}^{P} R_j
              \;\leq\; P + P \cdot R_{\max}
              \;=\; P(1 + R_{\max})
          \end{align}
          steps, where $R_{\max} \coloneqq \max_{1 \leq j \leq P} R_j$.
          The pooled bound $P(1+R)$ is recovered when $R_j = R$ for
          all $j$.

    \item \emph{Boundary and edge-case audit.}
          \begin{enumerate}
              \item[L1-E1] \emph{All $R_j = 0$ (no rework allowed):}
                           $\Psi(0) = P$; the run terminates in exactly
                           $P$ steps (forward transitions only), consistent
                           with Theorem~\ref{thm:acyclic}.
              \item[L1-E2] \emph{$P = 1$ (single-phase workflow):}
                           $\Psi(0) = 1 + R_1$; terminates in
                           $1 + R_1$ steps.  No contradiction with
                           Theorem~\ref{thm:rework}'s bound.
              \item[L1-E3] \emph{Rework at phase $j$ with $b^{(j)}_i =
                           R_j$} (budget exhausted for that phase):
                           the automaton transitions to the
                           \textsc{Escalated} terminal state instead of
                           \textsc{Rework} (same rule as in H3 of
                           Theorem~\ref{thm:rework}, applied per-phase).
                           The counter increment does not occur; the
                           bound is not violated.
              \item[L1-E4] \emph{Mixed model} (some phases pooled,
                           others per-phase): decompose into two disjoint
                           index sets $J_{\mathrm{pool}}$ and
                           $J_{\mathrm{ind}}$; apply~\eqref{eq:lim-phi-perphase}
                           to $J_{\mathrm{ind}}$ and the pooled counter to
                           $J_{\mathrm{pool}}$; the sum is still
                           non-negative and decreases by $1$ per step.
          \end{enumerate}

    \item \emph{Conclusion for L1.}
          The limitation is fully resolved: a deployment using
          per-phase counters is covered by the restated potential
          function $\Psi$~\eqref{eq:lim-phi-perphase}.  The
          termination bound~\eqref{eq:lim-perphase-bound} is tight,
          still-finite, and reduces to Theorem~\ref{thm:rework}'s
          bound when $R_j = R$ for all $j$.  The proof argument
          (non-negative potential decreasing by $1$ per step) is
          identical; only the numeric bound changes.
          $\blacksquare$
\end{enumerate}

% ----------------------------------------------------------------------------
\subsubsection*{Limitation~L2 --- Proof Sketch Status of the
                Token/Cost Bound}
\label{sec:lim:sketch}

\noindent\textbf{Statement of limitation.}
Section~\ref{sec:sketch} establishes an invocation-count bound
\begin{align}
    \label{eq:lim-sketch-bound}
    N \;\leq\; F_{\max} \cdot (P + B)
              \;=\; F_{\max} \cdot P(1 + R)
\end{align}
by a fully rigorous structural argument.  However, the proof is
labelled \textbf{SKETCH} because one additional step --- relating
the invocation count $N$ to a total \emph{token cost} $\mathcal{C}$
--- requires a per-invocation cost model that is not yet
empirically validated for this system.

\medskip
\noindent\textbf{Scope boundary.}
The sketch status concerns \emph{only} the token/cost translation
step; the invocation-count bound~\eqref{eq:lim-sketch-bound}
is itself fully proved and does not carry sketch status.

\medskip
\noindent\textbf{Resolution: structural argument for the
token/cost step.}

\begin{enumerate}
    \item \emph{Define the per-invocation cost function.}
          Let
          \begin{align}
              c : V \times \Sigma^{*}
              &\;\longrightarrow\; \mathbb{R}_{\geq 0}
              \notag
          \end{align}
          map each pair (workflow vertex $v$, input token sequence
          $w$) to a non-negative real cost.
          \emph{Assumption~A3 (to be validated empirically):}
          there exists a finite constant $c_{\max} \in \mathbb{R}_{>0}$
          such that
          \begin{align}
              \label{eq:lim-cmax}
              c(v, w) \;\leq\; c_{\max}
              \quad\text{for all } v \in V,\; w \in \Sigma^{*}.
          \end{align}

    \item \emph{Derive the total cost bound.}
          Given~\eqref{eq:lim-cmax} and the invocation-count
          bound~\eqref{eq:lim-sketch-bound}:
          \begin{align}
              \mathcal{C}
              &= \sum_{k=1}^{N} c(v_k, w_k)
              \notag\\
              &\leq \sum_{k=1}^{N} c_{\max}
              \notag\\
              &= N \cdot c_{\max}
              \notag\\
              &\leq F_{\max} \cdot P(1 + R) \cdot c_{\max}
              \;< \;\infty.
              \label{eq:lim-cost-bound}
          \end{align}
          Every inequality is exact: the first by
          \eqref{eq:lim-cmax}, the second by monotonicity of
          summation with non-negative terms, the third by
          \eqref{eq:lim-sketch-bound}.  All three quantities
          $F_{\max}$, $P$, $R$, $c_{\max}$ are finite by
          assumption (H1, H2 of Section~\ref{sec:sketch}, and
          A3 here).

    \item \emph{Promotion criterion.}
          The sketch is promoted to a full proof when and only when:
          \begin{enumerate}
              \item[L2-P1] A specific numeric value or an empirically
                           validated distribution for $c_{\max}$ is
                           established against at least one real vertical
                           slice (\S9.2, \ProofStd{}) of the MDAP system,
                           yielding a concrete upper bound
                           $\mathcal{C} \leq \mathcal{C}^{*}$ for a
                           computed $\mathcal{C}^{*} \in \mathbb{R}_{>0}$.
              \item[L2-P2] A benchmark dataset per \S9.2 of \ProofStd{}
                           is archived under
                           \code{MATHEMATICAL\_\allowbreak PROOFS/misc/}.
          \end{enumerate}
          Until L2-P1 and L2-P2 are satisfied, the total cost bound
          holds conditionally on A3; the invocation-count bound is
          unconditional.

    \item \emph{Boundary and edge-case audit.}
          \begin{enumerate}
              \item[L2-E1] \emph{Zero-cost invocation} ($c(v,w) = 0$):
                           allowed by $c_{\max} \geq 0$;
                           \eqref{eq:lim-cost-bound} still holds with
                           $\mathcal{C} \leq 0$, hence $\mathcal{C} = 0$
                           in that degenerate case.
              \item[L2-E2] \emph{Unbounded input sequences} ($|w|
                           \to \infty$): Assumption~A3 explicitly
                           bounds $c$ uniformly over all $w$; if A3
                           fails, the bound fails.  This is the exact
                           condition that must be validated empirically.
              \item[L2-E3] \emph{$N = 0$ (no invocations):}
                           $\mathcal{C} = 0 \leq c_{\max} \cdot 0 = 0$.
                           Bound holds trivially.
              \item[L2-E4] \emph{$R = 0$ (no rework):}
                           $B = 0$; $N \leq F_{\max} \cdot P$;
                           $\mathcal{C} \leq F_{\max} \cdot P \cdot
                           c_{\max}$.  Consistent with
                           Theorem~\ref{thm:rework} (no rework
                           $\Rightarrow$ terminates in $P$ steps).
          \end{enumerate}

    \item \emph{Conclusion for L2.}
          The structural argument for the total cost bound is complete
          and correct given A3.  The sole remaining gap is the
          empirical validation of A3 (items L2-P1 and L2-P2).  The
          bound~\eqref{eq:lim-cost-bound} is structurally tight,
          logically sound, and carries no hidden assumptions beyond
          A3. $\blacksquare$
\end{enumerate}

% ----------------------------------------------------------------------------
\subsubsection*{Limitation~L3 --- Semantic Correctness of Specialist
                Deliverables is Out of Scope}
\label{sec:lim:semantic}

\noindent\textbf{Statement of limitation.}
This paper formalizes the \emph{coordination scaffold} --- workflow
graph $G$, artifact state machine $\mathcal{M}$, trust-tier function
$\tau$, and conflict graph $C$ --- but does not formalize the
semantic correctness of any individual specialist's engineering
deliverable.  Specifically, Theorems~\ref{thm:acyclic}
through~\ref{thm:trust} make no claim about whether the content
produced by a Domain Specialist is \emph{correct} in any
domain-specific sense.

\medskip
\noindent\textbf{Scope boundary.}
This is an \emph{intended} scope exclusion (Scope~Exclusions
E1--E3, Section~\ref{sec:intro:scope}), not a defect in the proofs.
The exclusion is consistent with Assumption~A1 (specialist agents
are treated as oracles that produce outputs; their internal
correctness is a domain concern, not a coordination concern).

\medskip
\noindent\textbf{Resolution: formal interface for future semantic
extension.}

\begin{enumerate}
    \item \emph{Define a domain-correctness predicate.}
          For each workflow vertex $v \in V$, let
          \begin{align}
              \Pi_v : \mathrm{Artifacts} \;\longrightarrow\; \{0, 1\}
              \notag
          \end{align}
          be a domain-specific correctness predicate such that
          $\Pi_v(a) = 1$ if and only if artifact $a$ is correct
          according to the domain semantics associated with phase $v$.

    \item \emph{Identify the interface point.}
          In Theorem~\ref{thm:trust}, the Quality Gate certificate
          is currently modelled as a \emph{structural} event (it
          promotes a tier without inspecting the artifact's content).
          A semantic extension would replace the structural Quality
          Gate with a \emph{semantic} one:
          \begin{align}
              \text{QG-certifies}(a, v)
              \;\Longrightarrow\;
              \tau(a) = 3
              \;\;\text{and}\;\;
              \Pi_v(a) = 1.
              \notag
          \end{align}
          This requires no change to Theorems~\ref{thm:acyclic},
          \ref{thm:rework}, or~\ref{thm:multiplicity}; it
          \emph{strengthens} the conclusion of
          Theorem~\ref{thm:trust} by adding the conjunct
          $\Pi_v(a) = 1$.

    \item \emph{Validity of the strengthened conclusion.}
          The strengthened conclusion follows from the original proof
          plus the semantic Quality Gate axiom above.  Every step of
          the original inductive proof remains valid; the only
          addition is the verification that $\Pi_v(a) = 1$ holds
          whenever the Quality Gate fires, which is the responsibility
          of the domain-specific Quality Gate implementation.

    \item \emph{Constraint on the domain-correctness predicate.}
          For the extension to be well-formed, $\Pi_v$ must satisfy:
          \begin{enumerate}
              \item[L3-C1] \emph{Decidability:} $\Pi_v(a)$ is
                           computable for every artifact $a$ reachable
                           in the workflow, i.e., the Quality Gate
                           can decide correctness in finite time.
              \item[L3-C2] \emph{Monotonicity under certification:}
                           if $\Pi_v(a) = 1$ at the time of
                           certification and the artifact is not
                           subsequently modified, then $\Pi_v(a) = 1$
                           at all later stages.  (If artifacts can be
                           mutated post-certification, a re-certification
                           protocol is required, which is a separate
                           extension.)
              \item[L3-C3] \emph{Non-interference with coordination
                           scaffold:} $\Pi_v$ does not alter the
                           transition function $\delta$ of
                           $\mathcal{M}$ or the trust-tier function
                           $\tau$; it is a read-only predicate.
          \end{enumerate}

    \item \emph{Boundary and edge-case audit.}
          \begin{enumerate}
              \item[L3-E1] \emph{$\Pi_v$ is always true} (trivial
                           predicate): the extension reduces to the
                           original Theorem~\ref{thm:trust}.
              \item[L3-E2] \emph{$\Pi_v$ is always false} (no artifact
                           can be certified): the Quality Gate never
                           fires; the run terminates in
                           \textsc{Escalated} (by H3 of
                           Theorem~\ref{thm:rework}).  The coordination
                           scaffold is still sound; only the semantic
                           layer has failed.
              \item[L3-E3] \emph{Two concurrent Quality Gates certify
                           the same artifact with conflicting verdicts:}
                           prevented by L3-C3 (each $\Pi_v$ is
                           read-only and deterministic; certification
                           events are serialized by the automaton
                           $\mathcal{M}$, which is deterministic by
                           Definition~\ref{def:state-machine}).
          \end{enumerate}

    \item \emph{Conclusion for L3.}
          The scope exclusion is intentional and precisely bounded.
          A formal interface for semantic extension has been
          specified (items 1--4 above); the extension requires no
          changes to any existing theorem and introduces no new
          hypotheses to the coordination proofs.  The three
          constraints L3-C1--L3-C3 are necessary and sufficient
          for the extension to be well-formed.
          $\blacksquare$
\end{enumerate}

% ============================================================================
\section{Conclusion}
\label{sec:conclusion}

% ----------------------------------------------------------------------------
\subsection{Summary of Contributions}
\label{sec:conclusion:summary}

\begin{enumerate}

    \item \textbf{Formal model (Contribution~C1).}
          The Engineering Studio AI / MDAP architecture, already
          documented operationally across five internal vision documents,
          was given a precise, machine-auditable formal model consisting
          of four interdependent components.  Each component is defined
          by direct reference to a governing source; no detail is
          invented.

          \begin{enumerate}
              \item[C1a.] \emph{Workflow graph.}
                    A directed graph
                    \begin{align}
                        G \;=\; (V,\, E),
                        \qquad
                        |V| = 11,\quad |E| = 14,
                        \notag
                    \end{align}
                    whose vertices are the eleven workflow phases
                    (Definition~\ref{def:workflow-graph}) and whose edges
                    represent mandatory artifact-handoff transitions.
                    The augmented graph $G^{*}$ adds the single rework
                    back-edge $e^{*}$; that edge is the \emph{only}
                    source of potential cycles and is handled
                    separately by the Rework Automaton
                    (Definition~\ref{def:state-machine}).

              \item[C1b.] \emph{Artifact state machine.}
                    A deterministic finite automaton
                    \begin{align}
                        \mathcal{M}
                        \;=\;
                        (Q,\,\Sigma,\,\delta,\,q_0,\,F),
                        \notag
                    \end{align}
                    with state set $Q$, input alphabet $\Sigma$,
                    transition function $\delta$, initial state $q_0$,
                    and accepting set $F$ (Definition~\ref{def:state-machine}).
                    The automaton tracks the lifecycle status of each
                    artifact across the phases of $G$.

              \item[C1c.] \emph{Trust-tier function.}
                    A function
                    \begin{align}
                        \tau \;:\; V \;\longrightarrow\;
                        \{1,\,2,\,3\},
                        \notag
                    \end{align}
                    assigning each workflow vertex a trust tier
                    (Definition~\ref{def:trust}), grounded in the
                    eight-tier hierarchy of Table~8.1, \VisionClean{}.
                    The function encodes the invariant that no artifact
                    may be promoted past tier~2 without an explicit
                    Quality Gate certification.

              \item[C1d.] \emph{Conflict graph.}
                    A graph
                    \begin{align}
                        C \;=\; (\mathcal{F},\,E_C),
                        \notag
                    \end{align}
                    over the set $\mathcal{F}$ of twelve function
                    classes \FN{00}--\FN{11}
                    (Definition~\ref{def:conflict-graph}), where an
                    edge $(f_i, f_j) \in E_C$ denotes that the
                    Non-Overlap Rule prohibits a single agent instance
                    from holding both $f_i$ and $f_j$ simultaneously.
          \end{enumerate}

    \item \textbf{Four fully proved theorems (Contribution~C2).}
          Four safety and liveness properties were established as
          theorems with complete proofs.  Every hypothesis is
          explicitly validated; every inference step is justified;
          all boundary conditions and degenerate cases are treated
          explicitly.  The results are summarised in
          Table~\ref{tab:thm-summary}.

    \item \textbf{One honestly disclosed proof sketch (Contribution~C3).}
          One further property --- a finite token/cost invocation-count
          bound --- was provided as an explicitly labelled proof sketch
          (Section~\ref{sec:sketch}).  The sketch is complete as an
          argument, but rests on a per-invocation cost model that has
          not yet been empirically validated for this system.  It is
          treated as a sketch in strict conformance with \S10.3.3 of
          \ProofStd{}; it must not be treated as a full proof until the
          conditions stated in Section~\ref{sec:limitations} are
          satisfied.

    \item \textbf{Reusable formal vocabulary (Contribution~C4).}
          The four model components of C1 form a stable, reusable
          formal vocabulary.  Any future extension of the MDAP catalog
          (additional phases, revised function classes, new trust tiers,
          or modified rework budgets) can be audited against this
          vocabulary without re-deriving Theorems~\ref{thm:acyclic}
          through~\ref{thm:trust} from scratch; only the relevant
          hypothesis list for the modified theorem must be re-checked.

\end{enumerate}

% ----------------------------------------------------------------------------
\begin{table}[H]
\centering
\caption{Summary of theorems established in this paper.
         Column ``Type'' records the proof strategy
         per \S4.1.1 of \ProofStd{}.
         Column ``Guarantee class'' distinguishes safety
         properties (bad states are unreachable) from
         liveness properties (good states are eventually reached).}
\label{tab:thm-summary}
\small
\begin{tabular}{@{}clllp{4.8cm}@{}}
\toprule
\textbf{No.} & \textbf{Theorem} & \textbf{Type} &
\textbf{Class} & \textbf{Core claim} \\
\midrule
\ref{thm:acyclic}
  & Acyclicity
  & Direct
  & Liveness
  & Nominal workflow graph $G$ is a DAG; forward
    progress is guaranteed on every nominal
    execution path. \\[4pt]
\ref{thm:rework}
  & Bounded rework
  & Well-founded induction
  & Liveness
  & Every run of $\mathcal{M}$ reaches a terminal
    state in at most $P + B$ transitions; the
    augmented graph $G^{*}$ does not loop
    indefinitely. \\[4pt]
\ref{thm:multiplicity}
  & Agent multiplicity
  & Combinatorial (pigeonhole)
  & Safety
  & The Non-Overlap Rule requires at least
    $\chi(C)$ distinct agent instances, where
    $\chi(C) \geq 4$; single-responsibility
    is structurally enforced. \\[4pt]
\ref{thm:trust}
  & Trust-tier invariant
  & Structural induction
  & Safety
  & For every subagent invocation sequence of
    any finite length, no artifact reaches
    tier~3 without a Quality Gate certificate;
    trust is never silently promoted. \\
\bottomrule
\end{tabular}
\end{table}

% ----------------------------------------------------------------------------
\subsection{Verification of Objectives}
\label{sec:conclusion:objectives}

\begin{enumerate}

    \item \textbf{Objective~O1 (Formal model construction).}
          Satisfied by Contribution~C1.  All four model components
          (C1a--C1d) are grounded (\S\ref{sec:intro:objectives},
          item~O1a), sufficient for the four proved theorems (O1b),
          and structured for reuse (O1c).

    \item \textbf{Objective~O2 (Proof of safety/liveness properties).}
          Satisfied by Contribution~C2 and Table~\ref{tab:thm-summary}.
          Each proof names its proof type before the proof body (O2a),
          validates its hypotheses explicitly (O2b), justifies every
          non-trivial inference (O2c), addresses all boundary conditions
          (O2d), and closes with \verb|\qedhere| (O2e).

    \item \textbf{Objective~O3 (Proof sketch disclosure).}
          Satisfied by Contribution~C3.  Exactly one claim
          (Section~\ref{sec:sketch}) carries the SKETCH label (O3a),
          with the explicit reason for incompleteness (O3b) and the
          minimum additional work required for promotion to a full
          proof (O3c).

    \item \textbf{Objective~O4 (Auditability).}
          Satisfied by the structure of every proof in
          Sections~\ref{sec:proofs} and~\ref{sec:sketch}.  An
          independent third-party auditor can verify each step
          without consulting the authors (O4a), locate the governing
          document section justifying each non-mathematical premise
          (O4b), and confirm the absence of hidden assumptions by
          reading the hypothesis lists alone (O4c).

\end{enumerate}

% ----------------------------------------------------------------------------
\subsection{Significance of the Results}
\label{sec:conclusion:significance}

\begin{enumerate}

    \item \textbf{Liveness and safety constitute a complete coordination
          guarantee.}
          Theorems~\ref{thm:acyclic} and~\ref{thm:rework} together
          establish the liveness side of the coordination guarantee:
          the system does not hang indefinitely on any nominal or
          rework path.
          Specifically:
          \begin{enumerate}
              \item[S1a.] By Theorem~\ref{thm:acyclic}, every nominal
                          execution path in $G$ visits each vertex at
                          most once and terminates at a designated
                          terminal vertex in at most $|V| = 11$ steps.
              \item[S1b.] By Theorem~\ref{thm:rework}, every run of
                          the augmented graph $G^{*}$ terminates in
                          at most
                          \begin{align}
                              P + B
                              \;=\;
                              P + R \cdot P
                              \;=\;
                              P(1 + R)
                              \;<\; \infty
                              \notag
                          \end{align}
                          transitions, and the concrete invocation
                          count satisfies $N \leq 66(1+R)$.
          \end{enumerate}
          Theorems~\ref{thm:trust} and~\ref{thm:multiplicity} together
          establish the safety side: no bad state is reachable by any
          sequence of agent invocations.
          Specifically:
          \begin{enumerate}
              \item[S2a.] By Theorem~\ref{thm:trust}, for every
                          invocation sequence of length $n \geq 0$, the
                          trust-tier invariant
                          \begin{align}
                              \forall v \in V,\;
                              \tau(v) = 3
                              \;\Longrightarrow\;
                              v\text{ was certified by a Quality Gate}
                              \notag
                          \end{align}
                          holds; untrusted content is never silently
                          promoted to certified status.
              \item[S2b.] By Theorem~\ref{thm:multiplicity}, the
                          Non-Overlap Rule is structurally enforced:
                          any valid assignment of agents to function
                          classes requires at least
                          $\chi(C) \geq 4$ distinct instances,
                          so single-point-of-responsibility collapse
                          is impossible under the conflict-graph
                          constraint.
          \end{enumerate}
          The pairing of liveness and safety in this paper is the
          standard approach used in distributed-systems verification
          (e.g., temporal logic model checking); it is applied here
          to an agent-organizational graph rather than to a
          message-passing protocol, confirming that the verification
          methodology transfers cleanly to this domain.

    \item \textbf{No new axioms were required.}
          Every theorem hypothesis is either a standard mathematical
          fact (topological ordering of DAGs, the well-ordering
          principle, the pigeonhole principle, the principle of
          structural induction) or a directly cited rule from a
          governing document already in force in this repository
          (SOLID, ACID-for-workflows, the Non-Overlap Rule, the
          Loop-Safety Standard, the trust-tier hierarchy).
          This means the four proved theorems inherit the legitimacy
          of the engineering policy decisions already made; their
          validity is conditional only on the correctness of those
          policy decisions, and not on any additional modelling
          assumptions.

    \item \textbf{The sketch is the weakest link; it is the correct
          target for future formal effort.}
          The token/cost convergence bound of Section~\ref{sec:sketch}
          is the only incompletely proved claim in this paper.  Its
          incompleteness is \emph{technical}, not \emph{logical}:
          the argument is structurally sound, but rests on a
          per-invocation cost model that has not yet been validated
          for this system.  This is a well-defined, achievable
          gap --- not a fundamental obstacle --- and the minimum work
          required to close it is stated precisely in
          Section~\ref{sec:limitations}, item~2.

\end{enumerate}

% ----------------------------------------------------------------------------
\subsection{Future Work}
\label{sec:conclusion:futurework}

The following items are ordered by logical dependency: each item's
completion is a prerequisite for the items that follow it.

\begin{enumerate}

    \item \textbf{FW1 --- Specify and validate the per-invocation cost
          model.}
          Define a per-invocation cost function
          $c : V \times \Sigma^{*} \to \mathbb{R}_{\geq 0}$
          (mapping each workflow vertex and its input token sequence to
          a real-valued cost), establish empirically --- against the
          one real vertical slice recommended in \S13.1 of \VisionGov{}
          --- that the total cost is bounded by a computable function of
          $|V|$, $R$, and the per-phase token budgets already specified
          in the governing documents, and update Section~\ref{sec:sketch}
          to a full proof.
          This is the single highest-priority item because it completes
          Contribution~C3 and eliminates the only openly disclosed gap
          in the paper's coverage.

    \item \textbf{FW2 --- Prove Theorem~5 in full.}
          Once FW1 is complete, replace the proof sketch of
          Section~\ref{sec:sketch} with a full proof conforming to
          Objectives~O2a--O2e.  The proof type will be
          well-founded induction over the rework counter
          $b \in \{0, 1, \ldots, B\}$, mirroring the structure of
          the proof of Theorem~\ref{thm:rework}, with the per-phase
          cost function of FW1 as the potential function.

    \item \textbf{FW3 --- Empirical hypothesis validation.}
          Validate every theorem's hypotheses empirically against the
          one real vertical slice recommended in \S13.1 of \VisionGov{}.
          Specifically, the following hypotheses require empirical
          confirmation before the theorems are considered production-grade:
          \begin{enumerate}
              \item[FW3a.] Hypothesis~H2 of Theorem~\ref{thm:rework}:
                    that the per-phase retry limit $R$ is finite and
                    consistently enforced across all phases of a live
                    deployment.
              \item[FW3b.] Hypothesis~H3 of Theorem~\ref{thm:trust}:
                    that every Quality Gate in a live deployment raises
                    the trust tier by exactly one unit (i.e., that no
                    Quality Gate implementation silently skips the
                    certification step or applies a non-unit increment).
              \item[FW3c.] The Non-Overlap Rule
                    (Theorem~\ref{thm:multiplicity}, hypothesis):
                    that no deployed agent instance is
                    assigned to two conflict-adjacent function classes
                    in any live run.
          \end{enumerate}
          Empirical validation should produce a benchmark dataset
          per \S9.2 of \ProofStd{}, archived alongside the proofs
          under \code{MATHEMATICAL\_\allowbreak PROOFS/misc/}.

    \item \textbf{FW4 --- Proof migration to standalone artifacts.}
          Once FW1--FW3 are complete, promote all five proofs
          (four full proofs and the completed sketch from FW2) to
          standalone, independently version-controlled proof artifacts
          under \code{MATHEMATICAL\_\allowbreak PROOFS/misc/}, with
          SemVer-versioned filenames and version-history tables, per
          \S5.2 of \ProofStd{}.  Until FW1--FW3 are complete, the
          proofs remain embedded in this narrative document for
          readability, as noted in Appendix~\ref{app:migration}.

    \item \textbf{FW5 --- Extension to per-phase retry counters.}
          As noted in Section~\ref{sec:limitations}, item~1,
          Theorem~\ref{thm:rework}'s potential function uses a single
          pooled budget $B = R \cdot P$.  A deployment that tracks
          per-phase retry counters independently would require a
          restated version of that theorem with a per-phase potential
          function
          \begin{align}
              \Phi_{\text{per-phase}}
              \;=\;
              \sum_{i=1}^{P} (R_i - b_i),
              \notag
          \end{align}
          where $R_i$ is the per-phase retry limit for phase $i$ and
          $b_i \in \{0, 1, \ldots, R_i\}$ is the per-phase rework
          counter.  The restated theorem is still provable by the
          same well-founded induction argument (since each $\Phi_{\text{per-phase}}$ is bounded
          below by zero and decreases by at least one per rework
          event), but the numeric bound changes from
          $P(1+R)$ to $P + \sum_{i=1}^{P} R_i$.  This generalisation
          should be formalised once per-phase counters are adopted in
          a production deployment.

    \item \textbf{FW6 --- Formalisation of the semantic-correctness
          layer.}
          This paper explicitly excludes semantic correctness of
          specialist output (Scope Exclusion E1--E3,
          Section~\ref{sec:intro:scope}).  Future work may extend
          the formal model to include domain-specific correctness
          predicates for individual specialist agents
          (e.g., an electrical-validity predicate for the PCB
          layout agent, or a functional-correctness predicate for
          the firmware agent), enabling end-to-end formal guarantees
          that connect governance-layer liveness/safety to
          domain-layer correctness.  Such an extension would require
          domain-specific formal semantics that are outside the
          scope of this paper.

\end{enumerate}

% ----------------------------------------------------------------------------
\subsection{Closing Statement}
\label{sec:conclusion:closing}

The Engineering Studio AI / MDAP architecture admits a concise,
precise, and fully auditable formal model from which four of its five
central design claims are provable by well-established mathematical
techniques, using only axioms that are already adopted as engineering
policy in this repository.  The fifth claim is an honestly-disclosed
proof sketch with a clearly defined upgrade path.  The methodology
demonstrated here --- grounding a multi-agent governance architecture
in a directed graph, a finite automaton, a trust-tier function, and a
conflict graph, then deriving liveness and safety properties by
standard inductive and combinatorial arguments --- is general: it
applies without modification to any future extension of the MDAP
catalog, provided the governing documents remain as specified.  The
formal vocabulary of Contribution~C1 is therefore the primary
long-term deliverable of this paper; the four proved theorems are its
first, but not its last, inhabitants.



% ============================================================================
\newpage
\section*{References}
\label{sec:references}

\begin{sloppypar}
\begin{enumerate}
    \item Hu, H. (2026). ``VISION\_MULTIMODAL\_\allowbreak DEPLOYABLE\_AGENTS\_CLEAN.md.''
          Internal repository document, \texttt{markdowns/visions/}.
    \item Hu, H. (2026). ``VISION\_AGENTS\_COORDINATION\_\allowbreak ORCHESTRATION\_\allowbreak GOVERNANCE.md.'' Internal repository document,
          \texttt{markdowns/visions/}.
    \item Hu, H. (2026). ``VISION\_AGENTS\_BY\_FUNCTION.md.'' Internal
          repository document, \texttt{markdowns/visions/}.
    \item Hu, H. (2026). ``VISION\_AGENT\_WORKFLOW\_INTERAGENT\_\allowbreak COORDINATION.md.'' Internal repository document,
          \texttt{markdowns/visions/}.
    \item Hu, H. (2026). ``VISION\_AMD\_LABLAB\_HACKATHON\_ENGINEERING\_\allowbreak STUDIO.md.'' Internal repository document,
          \texttt{markdowns/visions/}.
    \item Hu, H. ``mathematical\_proof\_\allowbreak structures\_standards.txt.''
          Internal standard, \texttt{coding\_stds/mathematics/}.
    \item Hu, H. ``latex\_style\_guide.txt.'' Internal standard,
          \texttt{coding\_stds/documentation/}.
    \item Cormen, T. H., Leiserson, C. E., Rivest, R. L., \& Stein, C.
          (2009). \textit{Introduction to Algorithms} (3rd ed.). MIT
          Press. Theorem 22.12 (topological ordering of DAGs).
    \item Rosen, K. H. (2019). \textit{Discrete Mathematics and Its
          Applications} (8th ed.). McGraw-Hill. \S5.2 (well-ordering
          principle), \S6.2 (pigeonhole principle).
\end{enumerate}
\end{sloppypar}

% ============================================================================
\newpage
\appendix
\section{Notation Table}
\label{app:notation}

\begin{table}[H]
\centering
\caption{Notation used throughout Sections~\ref{sec:model}--\ref{sec:sketch}.}
\label{tab:notation}
\begin{tabular}{@{}ll@{}}
\toprule
\textbf{Symbol} & \textbf{Meaning} \\
\midrule
$G=(V,E)$ & Nominal workflow graph (Def.~\ref{def:workflow-graph}) \\
$\mathcal{M}$ & Artifact state machine (Def.~\ref{def:state-machine}) \\
$\tau$ & Trust-tier function (Def.~\ref{def:trust}) \\
$C=(\cdot,E_C)$ & Function conflict graph (Def.~\ref{def:conflict-graph}) \\
$R$ & Per-phase-gate retry limit \\
$P$ & Number of phases, $|V|=11$ \\
$B$ & Pooled rework budget, $B=R\cdot P$ \\
$F_{\max}$ & Per-turn fan-out ceiling \\
$D_{\max}$ & Nesting-depth limit, $D_{\max}=2$ \\
\bottomrule
\end{tabular}
\end{table}

% ----------------------------------------------------------------------------
\section{Summary of Key Definitions}
\label{app:key-defs}

\begin{table}[H]
\centering
\caption{Core formal objects defined in Section~\ref{sec:model}.}
\label{tab:key-defs}
\begin{tabularx}{\linewidth}{@{}>{\bfseries}p{3.2cm} X@{}}
\toprule
Object & Definition (summary) \\
\midrule
Workflow graph $G=(V,E)$ &
  Directed graph; $V$ = 11 phase-gate nodes, $E$ = directed
  precedence edges between consecutive phases.  No self-loops or
  multi-edges. \\[4pt]
State machine $\mathcal{M}$ &
  Finite automaton over artifact states
  $\{\textsc{Draft},\textsc{Review},\textsc{Approved},\textsc{Rejected},\textsc{Archived}\}$
  with transitions driven by gate-agent decisions. \\[4pt]
Trust-tier function $\tau$ &
  Map $\tau: \text{Agents} \to \{1,2,3\}$ assigning each agent a
  trust tier; tier-1 agents may not directly write
  production artifacts without tier-2/3 gate approval. \\[4pt]
Conflict graph $C=(\cdot,E_C)$ &
  Undirected graph on specialist-agent nodes; an edge $(a,b)\in E_C$
  signals that agents $a$ and $b$ must not hold concurrent write
  locks on the same artifact region. \\[4pt]
Retry limit $R$ &
  Maximum number of \textsc{Rejected}$\to$\textsc{Draft} rework
  cycles permitted per phase gate before the workflow halts. \\[4pt]
Fan-out ceiling $F_{\max}$ &
  Upper bound on the number of specialist agents that the orchestrator
  may spawn in a single turn. \\[4pt]
Nesting depth $D_{\max}$ &
  Maximum allowed depth of sub-orchestrator nesting; set to $D_{\max}=2$
  by architectural policy. \\
\bottomrule
\end{tabularx}
\end{table}

% ----------------------------------------------------------------------------
\section{Summary of Key Results}
\label{app:key-results}

\noindent The four proved theorems and one disclosed proof sketch are
collected below for quick reference.  Forward references point to the
body text for complete statements and proofs.

\begin{table}[H]
\centering
\caption{Proved theorems — structural and safety results
         (Sections~\ref{sec:sketch}--\ref{sec:proof}).}
\label{tab:results-proved}
\begin{tabularx}{\linewidth}{@{}>{\bfseries}p{2.4cm} >{\RaggedRight}p{2.8cm} X@{}}
\toprule
Result & Location & Summary statement \\
\midrule
Theorem~1 &
  \S\,Acyclicity &
  $G=(V,E)$ is a DAG\@.  Because every edge strictly increases the
  phase index, no directed cycle can exist, so topological execution
  order is well-defined and deadlock-free. \\[4pt]
Theorem~2 &
  \S\,Bounded rework &
  The rework process terminates in at most $B = R \cdot P$ total
  retry steps across all phases.  Proof by a monotone natural-number
  budget that decreases on every retry and reaches zero in finite
  steps. \\[4pt]
Theorem~3 &
  \S\,Minimum multiplicity &
  At least $\lceil|V|/F_{\max}\rceil$ specialist-agent slots are
  required to cover all phase-gate nodes within the fan-out
  constraint.  Follows from the pigeonhole principle. \\[4pt]
Theorem~4 &
  \S\,Trust-tier invariant &
  $\tau$ is preserved as an invariant throughout every execution trace:
  no transition in $\mathcal{M}$ and no edge traversal in $G$ can
  demote or promote an agent's tier without an explicit governance
  action.  Proof by structural induction on execution traces. \\
\bottomrule
\end{tabularx}
\end{table}

\begin{table}[H]
\centering
\caption{Proof sketch — liveness under resource constraints
         (Section~\ref{sec:sketch}).}
\label{tab:results-sketch}
\begin{tabularx}{\linewidth}{@{}>{\bfseries}p{2.4cm} >{\RaggedRight}p{2.8cm} X@{}}
\toprule
Result & Location & Summary statement \\
\midrule
Sketch~1 &
  \S\,Token/cost convergence &
  Under fixed per-agent token budgets and a finite artifact state
  space, cumulative token expenditure is bounded above by a closed-form
  expression.  Promoted to a full proof once the token-budget policy
  constants are ratified in the governing documents. \\
\bottomrule
\end{tabularx}
\end{table}

% ----------------------------------------------------------------------------
\section{Summary of Key Formulas}
\label{app:key-formulas}

\begin{table}[H]
\centering
\caption{Closed-form expressions and bounds used in the proofs
         (Part~A: workflow and rework).}
\label{tab:formulas-a}
\begin{tabularx}{\linewidth}{@{}>{\(}p{3.6cm}<{\)} X@{}}
\toprule
\multicolumn{1}{@{}l}{\textbf{Formula}} & \textbf{Meaning / context} \\
\midrule
|V| = 11 &
  Fixed cardinality of the phase-gate node set; each node corresponds
  to one discipline gate in the MDAP workflow. \\[4pt]
B = R \cdot P &
  Pooled rework budget: total retry steps available across all $P$
  phases, each phase permitting at most $R$ retries (Theorem~2). \\[4pt]
0 \le b_k \le B,\; b_{k+1} < b_k &
  Monotone budget sequence used in the termination argument; strict
  decrease on every retry step guarantees finiteness. \\[4pt]
D_{\max} = 2 &
  Architectural constant bounding sub-orchestrator nesting depth;
  enforced by the governance layer, not derived. \\
\bottomrule
\end{tabularx}
\end{table}

\begin{table}[H]
\centering
\caption{Closed-form expressions and bounds used in the proofs
         (Part~B: agent multiplicity and fan-out).}
\label{tab:formulas-b}
\begin{tabularx}{\linewidth}{@{}>{\(}p{3.6cm}<{\)} X@{}}
\toprule
\multicolumn{1}{@{}l}{\textbf{Formula}} & \textbf{Meaning / context} \\
\midrule
\lceil |V| / F_{\max} \rceil &
  Minimum number of specialist-agent slots required so that all
  phase-gate nodes are covered without exceeding the per-turn fan-out
  ceiling $F_{\max}$ (Theorem~3, pigeonhole argument). \\[4pt]
|E_C| \le \binom{|\text{Agents}|}{2} &
  Trivial upper bound on conflict-graph edges; the actual edge set is
  determined by the overlap structure of write-lock regions. \\[4pt]
\tau(a) \in \{1,2,3\} &
  Co-domain of the trust-tier function; tier values are ordinal, not
  cardinal (no arithmetic on tiers is defined). \\
\bottomrule
\end{tabularx}
\end{table}

% ----------------------------------------------------------------------------
\section{Summary of Architectural Findings}
\label{app:findings}

\noindent The following two tables summarise the principal qualitative
findings of the paper, split by category to stay within column-width limits.

\begin{table}[H]
\centering
\caption{Structural findings: graph and automaton properties.}
\label{tab:findings-structural}
\begin{tabularx}{\linewidth}{@{}>{\bfseries\RaggedRight}p{3.6cm} X@{}}
\toprule
Finding & Implication for the Engineering Studio AI architecture \\
\midrule
$G$ is a DAG &
  A unique topological execution order exists for all eleven
  phase-gate nodes.  The orchestrator may schedule phases without
  cycle detection at runtime; the ordering is a compile-time invariant
  of the architecture. \\[4pt]
$\mathcal{M}$ is finite and deterministic &
  Every artifact has a well-defined current state at all times.  No
  artifact can be simultaneously \textsc{Approved} and
  \textsc{Rejected}; the automaton disallows ambiguous governance
  records. \\[4pt]
Conflict graph $C$ is acyclic in lock order &
  The write-lock discipline derived from $E_C$ is deadlock-free under
  a consistent lock-ordering policy; no circular wait can arise between
  specialist agents sharing artifact regions. \\
\bottomrule
\end{tabularx}
\end{table}

\begin{table}[H]
\centering
\caption{Governance findings: liveness, safety, and resource bounds.}
\label{tab:findings-governance}
\begin{tabularx}{\linewidth}{@{}>{\bfseries\RaggedRight}p{3.6cm} X@{}}
\toprule
Finding & Implication for the Engineering Studio AI architecture \\
\midrule
Rework terminates in $\le B$ steps &
  The architecture provides a hard, statically-computable bound on the
  number of LLM calls in the rework loop.  Operators can translate $B$
  into a worst-case token-cost ceiling for budgeting and SLA purposes. \\[4pt]
Trust-tier invariant is inductive &
  Tier assignments survive every state transition; the governance
  layer cannot be bypassed by a sequence of individually-valid
  transitions.  This makes the trust hierarchy tamper-evident by
  construction. \\[4pt]
Minimum agent-slot bound is tight &
  The pigeonhole lower bound $\lceil|V|/F_{\max}\rceil$ is achieved
  by a round-robin allocation; no design with fewer slots can
  satisfy the fan-out constraint, establishing this as an absolute
  resource floor. \\[4pt]
Token-cost convergence is a proof sketch &
  The sketch identifies the missing precondition (ratified per-agent
  token-budget constants) and provides an upgrade path to a full
  proof, honestly disclosing the boundary of current formal coverage. \\
\bottomrule
\end{tabularx}
\end{table}

% ----------------------------------------------------------------------------
\section{Scope Exclusions at a Glance}
\label{app:exclusions}

\begin{table}[H]
\centering
\caption{Explicit scope exclusions and their rationale
         (see Section~\ref{sec:intro:scope} for full discussion).}
\label{tab:exclusions}
\begin{tabularx}{\linewidth}{@{}>{\bfseries}p{1.0cm} >{\RaggedRight}p{3.2cm} X@{}}
\toprule
ID & Exclusion & Rationale \\
\midrule
E1 & Semantic correctness of specialist output &
  Domain-specific correctness predicates (e.g., electrical-validity
  for PCB layout) require formal semantics outside the scope of this
  paper; the model governs workflow structure, not artifact content. \\[4pt]
E2 & LLM internals and prompt engineering &
  The behaviour of individual LLM calls is treated as a black box;
  the formal model is agnostic to the underlying model provider or
  prompt strategy. \\[4pt]
E3 & Distributed systems failure modes &
  Network partitions, node crashes, and message loss are not modelled;
  the formal model assumes reliable, ordered message delivery within
  each phase. \\
\bottomrule
\end{tabularx}
\end{table}

% ----------------------------------------------------------------------------
\section{Full Proof Migration Note}
\label{app:migration}

\begin{sloppypar}
Per \ProofStd{} \S5.1.1, formal
proof documents proper reside under \code{MATHEMATICAL\_\allowbreak PROOFS/}, organized
by domain (e.g. \code{correctness/}, \code{convergence/}). Should
Theorems~\ref{thm:acyclic}--\ref{thm:trust} be promoted to standalone,
independently versioned proof artifacts, they belong under
\code{MATHEMATICAL\_\allowbreak PROOFS/misc/} (cross-domain: graph theory, loop safety,
and agent-governance combined) with their own SemVer-versioned filenames
and version-history tables, per \S5.2 of that standard. This paper does not
perform that migration; it embeds the proofs directly for readability of
the accompanying narrative.
\end{sloppypar}

% ============================================================================
\newpage
\section*{Changelog}
\label{sec:changelog}

\begin{table}[H]
\centering
\caption{Document Revision History (YPSV: YEAR.MAJOR.MINOR.PATCH).}
\begin{tabular}{@{}lllp{6.5cm}@{}}
\toprule
\textbf{Version} & \textbf{Date} & \textbf{Author} & \textbf{Description} \\
\midrule
2026.1.0.0 & 2026-07-05 & Hadrian Hu & Initial creation: formal model of the
Engineering Studio AI / MDAP workflow (graph, automaton, trust-tier
function, conflict graph) and four full proofs (acyclicity, bounded
rework termination, minimum agent multiplicity, trust-tier invariant) plus
one explicitly disclosed proof sketch (token/cost convergence). \\
\bottomrule
\end{tabular}
\end{table}



\end{document}
% ============================================================================
% END OF DOCUMENT
% ============================================================================
